
\chapter{Computer Architecture, Digital Logic Design, and Boolean Algebra}
\label{ch:computer_architecture_digital_logic}

This chapter establishes the theoretical, architectural, and circuit-level foundations of modern digital computing systems. It bridges the abstraction layers from semiconductor logic gates and Boolean minimization up through processor execution units, memory hierarchies, pipelined datapaths, and microarchitectural performance analysis.

\begin{enumerate}[label=\bfseries\arabic*.]
    \item \textbf{Computer Architecture and Digital Fundamentals:} Distinctions between architecture, organization, and digital logic; the classic Von Neumann and Harvard models; system buses; instruction cycles; and register transfer fundamentals.
    \item \textbf{Number Systems, Signed Arithmetic, and Codes:} Positional bases, radix conversions, two's complement arithmetic, hardware overflow detection, BCD, Gray codes, and character encodings.
    \item \textbf{Boolean Algebra and Logic Gates:} Axiomatic postulates, fundamental algebraic theorems, De Morgan's laws, standard and universal gates (NAND/NOR implementations).
    \item \textbf{Combinational Logic Design and Minimization:} Adders, subtractors, lookahead carry generators, decoders, encoders, multiplexers, demultiplexers, tri-state buffers, and Karnaugh map (K-map) optimization.
    \item \textbf{Sequential Circuits and Memory Elements:} Latches, flip-flops, setup/hold timing margins, metastability, synchronous vs. asynchronous counters, shift registers, and Finite State Machines (FSM: Mealy and Moore models).
    \item \textbf{Processor Design, Addressing Modes, and Pipelining:} Machine instruction formats, operand addressing modes, execution datapaths, 5-stage RISC pipelining, and hazard mitigation (structural, data, control).
    \item \textbf{Processor Families and Storage Hierarchies:} RISC vs. CISC, Intel 8085/8086 microprocessors, memory hierarchy, cache mapping architectures, replacement policies, cache coherence (MESI), and primary/secondary storage technologies.
    \item \textbf{Processor Performance and Comprehensive Examination Bank:} CPI, MIPS, Amdahl's Law, multi-core concurrency, Flynn's taxonomy, worked numerical problems, and viva examination interview questions.
\end{enumerate}


\section{Computer Architecture and Organization}
\label{sec:arch_and_org}

The study of computing hardware is partitioned into distinct abstraction levels. Engineers must distinguish between the programmer-visible interface and the underlying physical implementation.

\subsection{What Is Computer Architecture?}
\textbf{Computer Architecture} defines the functional behavior, attributes, and programmer-visible structure of a computer system. It represents the conceptual contract between software (operating systems, compilers, applications) and the hardware execution engine.

Key architectural dimensions include:
\begin{itemize}
    \item \textbf{Instruction Set Architecture (ISA):} The collection of machine-level instructions, operation codes (opcodes), and supported instruction formats.
    \item \textbf{Register Set:} The programmer-visible general-purpose and special-purpose registers, their data bit-widths, and specific roles.
    \item \textbf{Data Types and Formats:} The fundamental primitives natively recognized by hardware (e.g., 8-bit integers, 32-bit signed integers, IEEE 754 floating-point values).
    \item \textbf{Memory Addressing Modes:} The mathematical mechanisms available for computing effective operand addresses in primary memory.
    \item \textbf{Memory Management and Protection:} Address space layout, virtual memory paging/segmentation models, and privilege execution levels (e.g., User vs. Kernel mode).
    \item \textbf{Interrupt and Exception Handling:} The architectural protocol for diverting processor control upon synchronous traps or asynchronous hardware interrupts.
\end{itemize}

\textbf{Practical Example:} The Intel/AMD x86-64 architecture mandates sixteen 64-bit general-purpose registers (RAX, RBX, RCX, RDX, RSI, RDI, RBP, RSP, R8--R15), a flat 64-bit virtual address space, and standard condition code flags. Any program compiled to conform to x86-64 will execute deterministically on an Intel Core i9, an AMD Ryzen 9, or a cloud-hosted virtual machine, despite radical differences in their internal microarchitectures.

\subsection{What Is Computer Organization?}
\textbf{Computer Organization} (also termed \textbf{Microarchitecture}) refers to the operational hardware units and their physical interconnections that realize the architectural specifications. It focuses on the internal engineering implementation transparent to the assembly-level programmer.

Key organizational dimensions include:
\begin{itemize}
    \item \textbf{Datapath Implementation:} The physical design of the Arithmetic Logic Unit (ALU), internal buses, multiplexers, and hardware shifters.
    \item \textbf{Control Unit Design:} The mechanism used to generate control signals (hardwired state-machine logic vs. microprogrammed ROM control stores).
    \item \textbf{Pipeline Depth and Hazards:} The number of execution stages, speculative branch predictors, and register renaming tables.
    \item \textbf{Memory Hierarchy Realization:} The number, sizes, associativity, and latency of Level-1, Level-2, and Level-3 on-chip cache memories.
    \item \textbf{Bus Signaling and Interconnects:} Clock frequencies, bus arbitration schemes, PCIe lanes, and memory controller timing.
\end{itemize}

\begin{table}[htbp]
\centering
\small
\caption{Comparative Analysis: Computer Architecture versus Computer Organization}
\label{tab:arch_vs_org}
\begin{tabularx}{\textwidth}{l X X}
\toprule
\textbf{Dimension} & \textbf{Computer Architecture} & \textbf{Computer Organization (Microarchitecture)} \\
\midrule
\textbf{Definition} & Programmer-visible functional specifications and hardware capabilities. & Physical realization and hardware component interconnections. \\
\textbf{Visibility} & Directly visible to the system programmer, compiler writer, and assembly engineer. & Completely transparent to application and assembly-level software. \\
\textbf{Core Elements} & ISA, addressing modes, register set, data representations, interrupt models. & ALU design, pipeline stages, control signals, cache structures, bus clocking. \\
\textbf{Longevity} & Highly stable across decades (e.g., x86 ISA has persisted since 1978). & Changes frequently across CPU generations (e.g., Skylake, Zen 4, Golden Cove). \\
\textbf{Design Goal} & Maximize software compatibility, programmer ergonomics, and expressive power. & Maximize instructions per cycle (IPC), clock frequency, and energy efficiency. \\
\textbf{Classic Analogy} & Architectural blueprint of a house (number of rooms, door placements). & Structural building materials used (timber vs. steel beams, HVAC ducting). \\
\bottomrule
\end{tabularx}
\end{table}

\subsection{Computer Architecture versus Digital Logic Design}
Digital systems operate across a hierarchical pyramid of abstractions:
\begin{enumerate}[label=\bfseries\arabic*.]
    \item \textbf{High-Level Software:} C/C++, Java, Python, and Operating Systems.
    \item \textbf{Instruction Set Architecture (ISA):} Assembly language, machine code, register models.
    \item \textbf{Microarchitecture (Organization):} Datapaths, pipeline registers, branch predictors, cache controllers.
    \item \textbf{Digital Logic Design:} Combinational gates (AND, OR, XOR), sequential elements (latches, flip-flops), adders, multiplexers, and finite state machines.
    \item \textbf{Analog/Transistor Circuits:} CMOS inverter topologies, pull-up pMOS networks, pull-down nMOS networks, propagation delays.
    \item \textbf{Solid-State Physics:} Silicon doping, bandgaps, electron mobility, lithographic gate pitch.
\end{enumerate}

\textbf{Digital Logic Design} serves as the hardware substrate upon which computer microarchitecture is constructed. A computer architect designs an Instruction Set; a computer organization engineer maps that ISA into pipeline registers and execution units; and a digital logic engineer translates those units into gate-level netlists and layout schematics.


\section{Fundamental Computer System Components and Architecture}
\label{sec:system_components}

Every general-purpose digital computer requires five classical functional units operating under synchronous timing control to execute algorithmic instructions.

\subsection{Basic Functional Units}
\begin{enumerate}[label=\bfseries\alph*.]
    \item \textbf{Input Unit:} Transduces external electro-mechanical or optical information into binary electrical representations readable by memory and processors (e.g., keyboard encoders, network interface controllers, NVMe storage controllers).
    \item \textbf{Output Unit:} Converts binary computational results into external human-readable or machine-actuated physical formats (e.g., display framebuffers, DAC sound generators, network transmitters).
    \item \textbf{Memory Unit:} Stores program instructions and data operands. Partitioned hierarchically into primary working storage (semiconductor RAM) and auxiliary persistent storage (solid-state drives, magnetic media).
    \item \textbf{Arithmetic Logic Unit (ALU):} The computational core responsible for executing arithmetic operations (addition, subtraction, multiplication, division) and bitwise logical operations (AND, OR, XOR, NOT, bit-shifts).
    \item \textbf{Control Unit (CU):} The administrative coordinator that orchestrates instruction fetching, operand decoding, and timing signal distribution across all system datapaths.
\end{enumerate}

\subsection{Central Processing Unit (CPU)}
The \textbf{Central Processing Unit (CPU)} integrates the ALU, the Control Unit, and an internal register file onto a single microchip (microprocessor).

\begin{table}[htbp]
\centering
\small
\caption{Functional Differences: Arithmetic Logic Unit (ALU) versus Control Unit (CU)}
\label{tab:alu_vs_cu}
\begin{tabularx}{\textwidth}{l X X}
\toprule
\textbf{Feature} & \textbf{Arithmetic Logic Unit (ALU)} & \textbf{Control Unit (CU)} \\
\midrule
\textbf{Primary Function} & Executes numerical computation and Boolean logic on binary data operands. & Fetches, decodes, and sequences instructions; directs hardware operations. \\
\textbf{Inputs} & Binary data operands from registers or memory; operational control codes. & Machine instructions from the Instruction Register (IR); clock and status flags. \\
\textbf{Outputs} & Computational results, condition code status flags (Zero, Carry, Overflow). & Multiplexer select signals, ALU function codes, memory read/write strobes. \\
\textbf{Analogy} & The computational worker/factory floor of the computer. & The executive management/orchestrator of the enterprise. \\
\bottomrule
\end{tabularx}
\end{table}

\subsection{Core CPU Registers}
Registers are high-speed semiconductor storage locations situated directly inside the processor datapath, accessible within a single clock phase without bus latency.

\begin{table}[htbp]
\centering
\small
\caption{Standard CPU Register Roles and Architectural Functions}
\label{tab:cpu_registers}
\begin{tabularx}{\textwidth}{l l X}
\toprule
\textbf{Register} & \textbf{Full Name} & \textbf{Architectural Function and Operation} \\
\midrule
\textbf{PC} & Program Counter & Holds the memory address of the next instruction to be fetched. Automatically increments after fetch unless overridden by a branch, jump, or interrupt vector. \\
\textbf{IR} & Instruction Register & Latches the binary machine instruction fetched from memory while the control unit decodes its opcode and operand specifiers. \\
\textbf{MAR} & Memory Address Register & Connected to the address bus; holds the physical memory address for an impending read or write bus cycle. \\
\textbf{MDR / MBR} & Memory Data / Buffer Register & Connected to the bidirectional data bus; acts as a synchronous holding buffer for data written to or read from memory. \\
\textbf{ACC} & Accumulator & Primary working register in single-operand accumulator architectures; holds one operand and receives the result of ALU operations. \\
\textbf{SP} & Stack Pointer & Stores the current memory offset of the top of the call stack; dynamically decremented/incremented during \texttt{PUSH} and \texttt{POP} operations. \\
\textbf{Flags / PSW} & Program Status Word & Collection of individual flip-flops recording arithmetic status (Carry \texttt{C}, Zero \texttt{Z}, Sign/Negative \texttt{S}, Overflow \texttt{V}, Parity \texttt{P}) and system state (Interrupt Enable). \\
\textbf{GPRs} & General Purpose Registers & High-speed scratchpad registers (e.g., R0--R15, EAX, EBX) used by the compiler to minimize external memory traffic. \\
\bottomrule
\end{tabularx}
\end{table}

\subsection{System Bus Architecture}
A \textbf{bus} is a shared collection of physical conducting wires transmitting parallel electrical signals between system modules (CPU, Memory, I/O Controllers). The system bus is partitioned into three distinct functional buses:

\begin{enumerate}[label=\bfseries\arabic*.]
    \item \textbf{Address Bus:} Unidirectional bus driven exclusively by the bus master (CPU or DMA controller) to specify the memory location or I/O port being accessed.
    \item \textbf{Data Bus:} Bidirectional bus over which actual data bytes, operands, and machine instructions travel between the CPU, memory, and peripheral controllers.
    \item \textbf{Control Bus:} Bidirectional collection of individual control lines delivering timing strobes, command directives, and status acknowledgments (e.g., \texttt{MEMR}, \texttt{MEMW}, \texttt{IOR}, \texttt{IOW}, \texttt{READY}, \texttt{RESET}, \texttt{INTR}, \texttt{INTA}).
\end{enumerate}

\begin{table}[htbp]
\centering
\small
\caption{System Bus Partitioning: Directionality, Width, and Signaling}
\label{tab:system_buses}
\begin{tabularx}{\textwidth}{l c l X}
\toprule
\textbf{Bus Component} & \textbf{Directionality} & \textbf{Typical Width} & \textbf{Functional Responsibility} \\
\midrule
\textbf{Address Bus} & Unidirectional (CPU $\to$ Memory/IO) & 16, 32, 48, 64 bits & Defines maximum addressable memory capacity: $N$ lines address $2^N$ locations. \\
\textbf{Data Bus} & Bidirectional & 8, 16, 32, 64, 128 bits & Determines word transfer width and memory bus bandwidth per cycle. \\
\textbf{Control Bus} & Bidirectional (mixed lines) & 10--30 discrete lines & Synchronizes bus cycles, specifies read vs. write, and coordinates arbitration. \\
\bottomrule
\end{tabularx}
\end{table}

\subsection{Bus Width and Addressable Memory Calculation}
The number of address lines determines the maximum addressable memory space. If an address bus has $k$ lines, it can generate $2^k$ unique binary addresses:

\begin{equation}
\label{eq:addressable_space}
\text{Total Addressable Locations} = 2^k
\end{equation}

\begin{itemize}
    \item \textbf{Byte-Addressable Memory:} Each distinct address references exactly 1 byte (8 bits). A 32-bit address bus addresses $2^{32} \text{ bytes} = 4\text{ GiB}$. A 64-bit bus theoretical maximum is $2^{64} \text{ bytes} = 16\text{ Exabytes (EiB)}$.
    \item \textbf{Word-Addressable Memory:} Each address references a full word of width $W$ bytes. Total capacity is $2^k \times W$ bytes.
\end{itemize}

\subsection{Clock Speed, Clock Period, and Cycle Time}
Synchronous digital processors transition their state on the rising or falling edges of a continuous periodic master clock square wave.
\begin{equation}
\label{eq:clock_frequency}
f = \frac{1}{T} \quad \Longleftrightarrow \quad T = \frac{1}{f}
\end{equation}
where $f$ is clock frequency (in Hertz, Hz) and $T$ is clock cycle period (in seconds, s).
\begin{itemize}
    \item A $1\text{ GHz}$ CPU has a clock cycle time of $T = \frac{1}{10^9\text{ Hz}} = 1\text{ ns}$.
    \item A $4\text{ GHz}$ CPU has a clock cycle time of $T = \frac{1}{4 \times 10^9\text{ Hz}} = 0.25\text{ ns} = 250\text{ ps}$.
\end{itemize}

\subsection{The Instruction Execution Cycle}
Every machine instruction undergoes a multi-phase operational sequence:
\begin{enumerate}[label=\bfseries\arabic*.]
    \item \textbf{Instruction Fetch (IF):}
    \[ \text{MAR} \leftarrow \text{PC} \]
    \[ \text{IR} \leftarrow \text{Memory}[\text{MAR}], \quad \text{PC} \leftarrow \text{PC} + 4 \]
    \item \textbf{Instruction Decode (ID):} The control unit examines the opcode bits in the IR, determines the required operation, decodes register specifiers, and fetches operands from the register file.
    \item \textbf{Execute (EX):} The ALU operates on the fetched operands, computing an arithmetic result or an effective memory address.
    \item \textbf{Memory Access (MEM):} If the instruction is a \texttt{LOAD} or \texttt{STORE}, the computed address is sent to the MAR to read or write data memory.
    \item \textbf{Write-Back (WB):} The computational result or loaded memory value is latched into the destination general-purpose register.
\end{enumerate}

\subsection{Von Neumann versus Harvard Architecture}
\begin{table}[htbp]
\centering
\small
\caption{Architectural Comparison: Von Neumann versus Harvard Model}
\label{tab:von_neumann_vs_harvard}
\begin{tabularx}{\textwidth}{l X X}
\toprule
\textbf{Dimension} & \textbf{Von Neumann Architecture} & \textbf{Harvard Architecture} \\
\midrule
\textbf{Memory Structure} & Unified physical memory for both program instructions and data operands. & Separate physical memory spaces and separate buses for instructions and data. \\
\textbf{Bus Infrastructure} & Single shared address and data bus for instructions and data. & Dual independent bus sets (instruction bus and data bus). \\
\textbf{Concurrent Access} & Cannot fetch an instruction and read/write data in the same clock cycle. & Instruction fetch and data memory access occur concurrently in parallel. \\
\textbf{Hardware Cost} & Simpler microarchitectural design; lower pin count and circuit complexity. & Increased silicon area and bus routing complexity; higher chip pin count. \\
\textbf{Bottleneck} & Severe throughput ceiling (Von Neumann Bottleneck). & Bottleneck substantially eliminated at the level of primary memory access. \\
\textbf{Modern Usage} & Conceptual model of general-purpose computers (RAM holds code and data). & Internal design of modern CPUs (separate Level-1 Instruction and Data Caches). \\
\bottomrule
\end{tabularx}
\end{table}

\subsection{The Von Neumann Bottleneck}
The \textbf{Von Neumann Bottleneck} refers to the fundamental throughput limitation caused by the shared bus between the CPU and memory. Because the CPU processing speed historically scaled much faster than DRAM bus bandwidth and latency, the processor spent an increasing fraction of its execution cycles idling waiting for memory transfers (the \textit{Memory Wall}).

\textbf{Modern Architectural Mitigations:}
\begin{enumerate}[label=\bfseries\arabic*.]
    \item \textbf{Multi-Level On-Chip Caching (L1/L2/L3):} High-speed static RAM caches integrated on-die to satisfy 95\%+ of memory references locally.
    \item \textbf{Modified Harvard Cache Design:} Modern microprocessors split the L1 cache into independent L1 Instruction (L1-I) and L1 Data (L1-D) caches, providing concurrent dual-port bandwidth to the execution pipeline.
    \item \textbf{Hardware Prefetching:} Architectural predictors detect sequential memory access patterns and stream cache lines ahead of CPU execution.
    \item \textbf{Out-of-Order and Speculative Execution:} Processors execute independent instructions downstream while stalled on an outstanding cache miss.
\end{enumerate}


\section{Number Systems, Signed Representations, and Binary Arithmetic}
\label{sec:number_systems}

Digital computers utilize two-state semiconductor switches (transistors operating in saturation and cut-off regions), mandating binary base-2 representation for all internal quantities.

\subsection{Positional Radix Representation}
A number $N$ expressed in a positional base $r$ with digits $d_i$ is evaluated by the weighted polynomial:
\begin{equation}
\label{eq:radix_polynomial}
(N)_r = \sum_{i=-m}^{n-1} d_i \cdot r^i = d_{n-1} r^{n-1} + \dots + d_1 r^1 + d_0 r^0 + d_{-1} r^{-1} + \dots + d_{-m} r^{-m}
\end{equation}

\begin{table}[htbp]
\centering
\small
\caption{Fundamental Digital Number Systems}
\label{tab:number_systems}
\begin{tabularx}{\textwidth}{l c c X}
\toprule
\textbf{System} & \textbf{Base ($r$)} & \textbf{Allowed Digits} & \textbf{Primary Engineering Utility} \\
\midrule
\textbf{Binary} & 2 & 0, 1 & Native hardware representation; directly maps to transistor switch states. \\
\textbf{Octal} & 8 & 0--7 & Compact representation of binary (each octal digit equals exactly 3 bits). \\
\textbf{Decimal} & 10 & 0--9 & Standard human-readable numerical notation. \\
\textbf{Hexadecimal} & 16 & 0--9, A--F & Industry standard shorthand for binary bytes and memory addresses (4 bits/digit). \\
\bottomrule
\end{tabularx}
\end{table}

\subsection{Radix Conversion Procedures}
\begin{itemize}
    \item \textbf{Base $r$ to Decimal:} Expand using the weighted power series (Eq.~\ref{eq:radix_polynomial}).
    \[ (101101.101)_2 = 1\cdot 2^5 + 0\cdot 2^4 + 1\cdot 2^3 + 1\cdot 2^2 + 0\cdot 2^1 + 1\cdot 2^0 + 1\cdot 2^{-1} + 0\cdot 2^{-2} + 1\cdot 2^{-3} = 45.625_{10} \]
    \[ (2A3.8)_{16} = 2\cdot 16^2 + 10\cdot 16^1 + 3\cdot 16^0 + 8\cdot 16^{-1} = 512 + 160 + 3 + 0.5 = 675.5_{10} \]
    \item \textbf{Decimal to Binary (Integer Part):} Repeated division by 2, recording remainders from least significant bit (LSB) up to most significant bit (MSB).
    \item \textbf{Decimal to Binary (Fractional Part):} Repeated multiplication by 2, recording the generated integer digit (0 or 1) sequentially until termination or required precision.
    \item \textbf{Binary to Octal:} Group integer bits into triplets from right to left; replace each triplet with its octal equivalent ($000_2=0$ to $111_2=7$).
    \item \textbf{Binary to Hexadecimal:} Group integer bits into nibbles (quads) from right to left; replace each nibble with its hex character ($0000_2=0$ to $1111_2=\text{F}$).
\end{itemize}

\subsection{Signed Integer Representations}
Fixed-width $n$-bit computer registers represent signed integers via three primary conventions:

\begin{enumerate}[label=\bfseries\arabic*.]
    \item \textbf{Sign-Magnitude:} The MSB represents sign ($0 = +, 1 = -$), and the remaining $n-1$ bits denote magnitude.
    \begin{itemize}
        \item \textit{Flaws:} Contains two representations for zero ($+0 = 0000_2$, $-0 = 1000_2$); requires complex dual-mode addition/subtraction hardware.
        \item \textit{Range:} $[-(2^{n-1}-1), +(2^{n-1}-1)]$.
    \end{itemize}
    \item \textbf{One's Complement:} Positive numbers represented in standard binary. Negative numbers obtained by inverting all bits (bitwise NOT).
    \begin{itemize}
        \item \textit{Flaws:} Still has dual zeros ($+0 = 0000_2$, $-0 = 1111_2$); requires an end-around carry addition step.
        \item \textit{Range:} $[-(2^{n-1}-1), +(2^{n-1}-1)]$.
    \end{itemize}
    \item \textbf{Two's Complement:} Universal standard in modern microprocessors. Positive values have $\text{MSB}=0$. Negative values computed by taking the one's complement and adding 1:
    \begin{equation}
    \label{eq:twos_complement_def}
    [-X]_2 = [\sim X] + 1 = 2^n - X
    \end{equation}
    \begin{itemize}
        \item \textit{Mathematical Weight of MSB:} The MSB possesses negative weight: $-2^{n-1}$.
        \[ \text{Value} = -b_{n-1} 2^{n-1} + \sum_{i=0}^{n-2} b_i 2^i \]
        \item \textit{Advantages:} Exactly one representation for zero ($00000000_2$); subtraction $A - B$ executed directly by standard addition circuitry: $A + (-B)_2$.
        \item \textit{Range:} $[-2^{n-1}, +(2^{n-1}-1)]$. Note the asymmetric extra negative value (e.g., for 8 bits: $-128$ to $+127$).
    \end{itemize}
\end{enumerate}

\begin{table}[htbp]
\centering
\small
\caption{8-Bit Signed Representation Comparison Across Critical Numerical Values}
\label{tab:signed_representations}
\begin{tabularx}{\textwidth}{l c c c}
\toprule
\textbf{Decimal Value} & \textbf{Sign-Magnitude} & \textbf{One's Complement} & \textbf{Two's Complement} \\
\midrule
$+127$ & \texttt{01111111} & \texttt{01111111} & \texttt{01111111} \\
$+1$ & \texttt{00000001} & \texttt{00000001} & \texttt{00000001} \\
$+0$ & \texttt{00000000} & \texttt{00000000} & \texttt{00000000} \\
$-0$ & \texttt{10000000} & \texttt{11111111} & (Does not exist; unique zero) \\
$-1$ & \texttt{10000001} & \texttt{11111110} & \texttt{11111111} \\
$-127$ & \texttt{11111111} & \texttt{10000000} & \texttt{10000001} \\
$-128$ & (Out of range) & (Out of range) & \texttt{10000000} \\
\bottomrule
\end{tabularx}
\end{table}

\subsection{Two's Complement Subtraction}
To compute $A - B$ using an $n$-bit adder:
\begin{enumerate}[label=\bfseries\arabic*.]
    \item Invert all bits of $B$ (One's complement: $\sim B$).
    \item Assert the carry-in bit ($C_{in} = 1$) at the LSB position of the adder.
    \item Perform binary addition: $\text{Result} = A + \sim B + 1$.
    \item Discard any end-carry out of the MSB position ($C_{out}$).
\end{enumerate}

\textbf{Worked Example:} Compute $43_{10} - 25_{10}$ in 8-bit two's complement arithmetic.
\begin{itemize}
    \item $A = 43_{10} = 00101011_2$
    \item $B = 25_{10} = 00011001_2 \implies -B = \sim(00011001_2) + 1 = 11100110_2 + 1 = 11100111_2$
    \item Sum:
    \begin{center}
    \begin{tabular}{r c}
      & \texttt{00101011} \quad ($+43$) \\
    $+$ & \texttt{11100111} \quad ($-25$) \\
    \hline
    (1) & \texttt{00010010} \quad ($+18_{10}$)
    \end{tabular}
    \end{center}
    \item End-carry (1) is discarded. The 8-bit result is $00010010_2 = 18_{10}$. Correct.
\end{itemize}

\subsection{Arithmetic Overflow versus Hardware Carry}
In technical examinations, candidates frequently confuse the \textbf{Carry Flag (C)} and the \textbf{Overflow Flag (V)}:

\begin{itemize}
    \item \textbf{Carry Flag ($C$):} Indicates an arithmetic carry-out of the MSB during \textit{unsigned} addition, or a borrow during unsigned subtraction. Used for multi-precision arithmetic.
    \item \textbf{Overflow Flag ($V$):} Indicates that the signed arithmetic result exceeded the representable numerical range $[-2^{n-1}, +(2^{n-1}-1)]$, yielding a mathematically erroneous sign bit.
\end{itemize}

\textbf{Universal Overflow Detection Rule:} Overflow in two's complement addition occurs if and only if:
\begin{enumerate}[label=\bfseries\alph*.]
    \item Two positive numbers are added and produce a negative result:
    \[ (+A) + (+B) = (-C) \]
    \item Two negative numbers are added and produce a positive result:
    \[ (-A) + (-B) = (+C) \]
\end{enumerate}
\textit{Note:} Adding a positive number and a negative number can \textbf{never} produce an overflow.

\textbf{Hardware Logic Equation for Overflow:}
In an $n$-bit parallel adder, overflow is calculated by the XOR of the carry-in and carry-out of the MSB position:
\begin{equation}
\label{eq:overflow_xor}
V = C_{n-1} \oplus C_{n}
\end{equation}
where $C_{n-1}$ is the carry into the sign bit and $C_n$ is the carry out of the sign bit.

\subsection{Specialized Binary Codes}
\begin{enumerate}[label=\bfseries\arabic*.]
    \item \textbf{Binary-Coded Decimal (BCD / 8421 Code):} Encodes each decimal digit ($0$--$9$) with its 4-bit binary equivalent ($0000_2$--$1001_2$).
    \begin{itemize}
        \item \textit{Invalid States:} Six binary combinations ($1010_2$ to $1111_2$, values 10 to 15) are strictly invalid.
        \item \textit{BCD Correction Rule:} If a 4-bit BCD addition generates an invalid combination ($> 9$) or produces an auxiliary carry out ($AC=1$), add $6_{10} = 0110_2$ to skip the 6 invalid states.
    \end{itemize}
    \item \textbf{Gray Code (Reflected Binary Code):} A non-weighted, cyclic code possessing the \textbf{unit-distance property}: adjacent code words differ in exactly one bit position.
    \begin{itemize}
        \item \textit{Engineering Utility:} Eliminates false transient switching glitches in optical shaft encoders, angle transducers, and asynchronous FIFO pointer synchronization across clock domains.
        \item \textit{Binary to Gray Conversion:}
        \[ G_{n-1} = B_{n-1}, \quad G_i = B_{i+1} \oplus B_i \quad (0 \le i < n-1) \]
        \item \textit{Gray to Binary Conversion:}
        \[ B_{n-1} = G_{n-1}, \quad B_i = B_{i+1} \oplus G_i \quad (0 \le i < n-1) \]
    \end{itemize}
    \item \textbf{Character Encodings (ASCII and Unicode):}
    \begin{itemize}
        \item \textbf{ASCII:} 7-bit standard encoding $2^7 = 128$ characters (letters, digits, punctuation, and 33 non-printable control characters).
        \item \textbf{UTF-8:} Variable-length Unicode encoding using 1 to 4 bytes per character; backward-compatible with 7-bit ASCII in byte 1.
    \end{itemize}
\end{enumerate}



\section{Boolean Algebra Fundamentals, Axiomatic Laws, and Minimization}
\label{sec:boolean_algebra}

Digital circuits operate entirely on two discrete voltage states (representing binary logic $0$ and $1$). \textbf{Boolean Algebra}, formulated by George Boole and formalized axiomatically by E.~V.~Huntington, provides the mathematical framework for analyzing, optimizing, and minimizing digital switching networks.

\subsection{Boolean Variables and Fundamental Operations}
A Boolean variable $A \in \{0, 1\}$ represents a two-state signal. The three primitive operations are:
\begin{enumerate}[label=\bfseries\arabic*.]
    \item \textbf{Conjunction (AND):} Denoted $A \cdot B$ or $AB$. Output is $1$ if and only if all operands are $1$.
    \item \textbf{Disjunction (OR):} Denoted $A + B$. Output is $1$ if at least one operand is $1$.
    \item \textbf{Negation / Inversion (NOT):} Denoted $\overline{A}$ or $A'$. Flips $0 \to 1$ and $1 \to 0$.
\end{enumerate}

\subsection{Fundamental Laws and Theorems of Boolean Algebra}
Every Boolean algebraic manipulation relies on ten fundamental pairs of dual identities:

\begin{enumerate}[label=\bfseries\arabic*.]
    \item \textbf{Identity Law:}
    \[ A + 0 = A \qquad \text{and} \qquad A \cdot 1 = A \]
    \item \textbf{Null / Annihilation / Dominance Law:}
    \[ A + 1 = 1 \qquad \text{and} \qquad A \cdot 0 = 0 \]
    \item \textbf{Idempotent Law:}
    \[ A + A = A \qquad \text{and} \qquad A \cdot A = A \]
    \item \textbf{Complement / Inverse Law:}
    \[ A + \overline{A} = 1 \qquad \text{and} \qquad A \cdot \overline{A} = 0 \]
    \item \textbf{Involution Law (Double Negation):}
    \[ \overline{\overline{A}} = A \]
    \item \textbf{Commutative Law:}
    \[ A + B = B + A \qquad \text{and} \qquad A \cdot B = B \cdot A \]
    \item \textbf{Associative Law:}
    \[ (A + B) + C = A + (B + C) \qquad \text{and} \qquad (A \cdot B) \cdot C = A \cdot (B \cdot C) \]
    \item \textbf{Distributive Law:}
    \[ A \cdot (B + C) = (A \cdot B) + (A \cdot C) \]
    \[ A + (B \cdot C) = (A + B) \cdot (A + C) \quad \text{(Crucial dual form unique to Boolean algebra)} \]
    \item \textbf{Absorption Law:}
    \[ A + (A \cdot B) = A \qquad \text{and} \qquad A \cdot (A + B) = A \]
    \[ A + \overline{A} B = A + B \qquad \text{and} \qquad A \cdot (\overline{A} + B) = A B \]
    \item \textbf{Consensus Theorem (Redundancy Law):}
    \begin{equation}
    \label{eq:consensus_theorem}
    AB + \overline{A}C + BC = AB + \overline{A}C
    \end{equation}
    \[ \text{Dual Form:} \quad (A + B)(\overline{A} + C)(B + C) = (A + B)(\overline{A} + C) \]
    \textit{Explanation:} The term $BC$ is redundant because if $B=1$ and $C=1$, the expression evaluates to $1$ regardless of whether $A=1$ (making $AB=1$) or $A=0$ (making $\overline{A}C=1$).
\end{enumerate}

\subsection{De Morgan's Theorems}
De Morgan's laws establish the algebraic bridge between conjunction, disjunction, and inversion:

\begin{enumerate}[label=\bfseries\arabic*.]
    \item \textbf{First Theorem (NOR Equivalence):} The complement of a sum is equal to the product of the individual complements:
    \begin{equation}
    \label{eq:demorgan_1}
    \overline{A + B} = \overline{A} \cdot \overline{B}
    \end{equation}
    \item \textbf{Second Theorem (NAND Equivalence):} The complement of a product is equal to the sum of the individual complements:
    \begin{equation}
    \label{eq:demorgan_2}
    \overline{A \cdot B} = \overline{A} + \overline{B}
    \end{equation}
\end{enumerate}

\begin{table}[htbp]
\centering
\small
\caption{Truth Table Proof of De Morgan's Dual Theorems}
\label{tab:demorgan_proof}
\begin{tabularx}{\textwidth}{c c c c c c c c}
\toprule
$A$ & $B$ & $\overline{A}$ & $\overline{B}$ & $\overline{A + B}$ & $\overline{A} \cdot \overline{B}$ & $\overline{A \cdot B}$ & $\overline{A} + \overline{B}$ \\
\midrule
0 & 0 & 1 & 1 & \textbf{1} & \textbf{1} & \textbf{1} & \textbf{1} \\
0 & 1 & 1 & 0 & \textbf{0} & \textbf{0} & \textbf{1} & \textbf{1} \\
1 & 0 & 0 & 1 & \textbf{0} & \textbf{0} & \textbf{1} & \textbf{1} \\
1 & 1 & 0 & 0 & \textbf{0} & \textbf{0} & \textbf{0} & \textbf{0} \\
\bottomrule
\end{tabularx}
\end{table}

\subsection{The Principle of Duality}
The \textbf{Principle of Duality} states that any true Boolean algebraic identity remains valid if:
\begin{enumerate}[label=\bfseries\alph*.]
    \item Every AND ($\cdot$) operator is replaced by an OR ($+$) operator.
    \item Every OR ($+$) operator is replaced by an AND ($\cdot$) operator.
    \item Every binary constant $0$ is replaced by $1$, and every $1$ is replaced by $0$.
\end{enumerate}
\textit{Note:} Variable literals ($A, \overline{A}$) are left uncomplemented. For example, the dual of $A + \overline{A}B = A + B$ is $A \cdot (\overline{A} + B) = A \cdot B$.

\subsection{Algebraic Simplification Examples}
\textbf{Example 1:} Simplify $F = AB + A(B + C) + B(B + C)$.
\begin{align*}
F &= AB + AB + AC + BB + BC && \text{(Distributive Law)} \\
  &= AB + AC + B + BC && \text{(Idempotent: $AB+AB=AB, BB=B$)} \\
  &= AB + AC + B && \text{(Absorption: $B + BC = B$)} \\
  &= B(A + 1) + AC && \text{(Distributive Law)} \\
  &= B(1) + AC = \mathbf{B + AC} && \text{(Null Law: $A+1=1$)}
\end{align*}

\textbf{Example 2:} Simplify $F = (A + B)(A + \overline{B})(A + C)$.
\begin{align*}
(A + B)(A + \overline{B}) &= A + B\overline{B} = A + 0 = A && \text{(Distributive dual: $X+YZ = (X+Y)(X+Z)$)} \\
F &= A \cdot (A + C) = \mathbf{A} && \text{(Absorption Law: $A(A+C) = A$)}
\end{align*}


\section{Standard and Derived Digital Logic Gates}
\label{sec:logic_gates}

Logic gates are physical electronic switching circuits (typically CMOS transistor topologies) that implement primitive Boolean operations.

\subsection{Fundamental and Derived Logic Gates}
\begin{table}[htbp]
\centering
\small
\caption{Comprehensive Summary of Standard Logic Gates}
\label{tab:standard_logic_gates}
\begin{tabularx}{\textwidth}{l c l X}
\toprule
\textbf{Gate Name} & \textbf{Symbolic Equation} & \textbf{Behavioral Logic Condition} & \textbf{CMOS Transistor Count} \\
\midrule
\textbf{NOT (Inverter)} & $Y = \overline{A}$ & Output is active when input is low. & 2 transistors (1 pMOS, 1 nMOS) \\
\textbf{AND} & $Y = A \cdot B$ & Output is high only when all inputs are high. & 6 transistors (NAND + Inverter) \\
\textbf{OR} & $Y = A + B$ & Output is high if any input is high. & 6 transistors (NOR + Inverter) \\
\textbf{NAND} & $Y = \overline{AB}$ & Output is low only when all inputs are high. & 4 transistors (2 pMOS parallel, 2 nMOS series) \\
\textbf{NOR} & $Y = \overline{A + B}$ & Output is high only when all inputs are low. & 4 transistors (2 pMOS series, 2 nMOS parallel) \\
\textbf{XOR} & $Y = A \oplus B = \overline{A}B + A\overline{B}$ & Output is high when inputs differ (odd parity). & 8--10 transistors \\
\textbf{XNOR} & $Y = \overline{A \oplus B} = AB + \overline{A}\,\overline{B}$ & Output is high when inputs are identical (coincidence). & 8--10 transistors \\
\bottomrule
\end{tabularx}
\end{table}

\begin{table}[htbp]
\centering
\small
\caption{Unified 2-Input Logic Gate Truth Table}
\label{tab:unified_truth_table}
\begin{tabularx}{\textwidth}{c c c c c c c c}
\toprule
$A$ & $B$ & \textbf{AND} & \textbf{OR} & \textbf{NAND} & \textbf{NOR} & \textbf{XOR} & \textbf{XNOR} \\
\midrule
0 & 0 & 0 & 0 & 1 & 1 & 0 & 1 \\
0 & 1 & 0 & 1 & 1 & 0 & 1 & 0 \\
1 & 0 & 0 & 1 & 1 & 0 & 1 & 0 \\
1 & 1 & 1 & 1 & 0 & 0 & 0 & 1 \\
\bottomrule
\end{tabularx}
\end{table}

\subsection{Properties and Applications of XOR and XNOR Gates}
The Exclusive-OR (XOR) gate exhibits algebraic properties essential for arithmetic adders, parity checkers, and pseudo-random number generators:

\begin{enumerate}[label=\bfseries\alph*.]
    \item \textbf{Invert-Buffer Control Property:}
    \[ A \oplus 0 = A \quad (\text{Buffer Mode}) \qquad \text{and} \qquad A \oplus 1 = \overline{A} \quad (\text{Controlled Inverter Mode}) \]
    This property enables the standard adder/subtractor circuit to negate operands dynamically based on a single subtraction mode control line ($M$).
    \item \textbf{Self-Inverse Identity:}
    \[ A \oplus A = 0 \qquad \text{and} \qquad A \oplus \overline{A} = 1 \]
    \item \textbf{Commutative and Associative:}
    \[ A \oplus B = B \oplus A, \qquad (A \oplus B) \oplus C = A \oplus (B \oplus C) \]
    \item \textbf{Multi-Input Parity Generation:} A chain of XOR gates evaluates to $1$ if and only if an \textbf{odd number} of inputs are $1$. An XNOR gate chain detects an \textbf{even number} of $1$s.
\end{enumerate}


\section{Universal Gates: NAND and NOR Logic Syntheses}
\label{sec:universal_gates}

A logic gate is termed \textbf{Universal} (functionally complete) if any arbitrary Boolean switching function can be synthesized entirely using only that single gate type, with no external inverters or secondary gates required.

\subsection{Functional Completeness Proof: The NAND Gate}
To prove universality, we must demonstrate the synthesis of the complete set of primitive operations $\{\text{NOT}, \text{AND}, \text{OR}\}$ using exclusively 2-input NAND gates ($Y = \overline{AB}$):

\begin{enumerate}[label=\bfseries\arabic*.]
    \item \textbf{Synthesis of NOT using NAND (1 Gate):}
    Tie both inputs together:
    \[ Y = \overline{A \cdot A} = \overline{A} \]
    \item \textbf{Synthesis of AND using NAND (2 Gates):}
    Feed a NAND gate into a NAND-based inverter:
    \[ Y = \overline{\overline{A \cdot B}} = A \cdot B \]
    \item \textbf{Synthesis of OR using NAND (3 Gates):}
    Apply De Morgan's theorem $\overline{\overline{A} \cdot \overline{B}} = A + B$:
    Invert $A$ (Gate 1), invert $B$ (Gate 2), and feed both into Gate 3:
    \[ Y = \overline{\overline{A} \cdot \overline{B}} = A + B \]
    \item \textbf{Synthesis of NOR using NAND (4 Gates):}
    Generate OR (3 gates) and invert the output (Gate 4):
    \[ Y = \overline{A + B} \]
    \item \textbf{Synthesis of XOR using NAND (4 Gates):}
    \begin{align*}
    G_1 &= \overline{AB} \\
    G_2 &= \overline{A \cdot G_1} = \overline{A \cdot \overline{AB}} = \overline{A} + AB = \overline{A} + B \\
    G_3 &= \overline{B \cdot G_1} = \overline{B \cdot \overline{AB}} = \overline{B} + AB = \overline{B} + A \\
    Y &= \overline{G_2 \cdot G_3} = \overline{(\overline{A} + B)(\overline{B} + A)} = AB + \overline{A}\,\overline{B} = A \oplus B
    \end{align*}
    \item \textbf{Synthesis of XNOR using NAND (5 Gates):}
    Invert the 4-gate XOR output with a 5th NAND gate, or use an optimized 5-gate configuration.
\end{enumerate}

\subsection{Functional Completeness Proof: The NOR Gate}
Similarly, we synthesize $\{\text{NOT}, \text{OR}, \text{AND}\}$ using exclusively 2-input NOR gates ($Y = \overline{A + B}$):

\begin{enumerate}[label=\bfseries\arabic*.]
    \item \textbf{Synthesis of NOT using NOR (1 Gate):}
    Tie both inputs together:
    \[ Y = \overline{A + A} = \overline{A} \]
    \item \textbf{Synthesis of OR using NOR (2 Gates):}
    Feed a NOR gate into a NOR-based inverter:
    \[ Y = \overline{\overline{A + B}} = A + B \]
    \item \textbf{Synthesis of AND using NOR (3 Gates):}
    Apply De Morgan's theorem $\overline{\overline{A} + \overline{B}} = A \cdot B$:
    Invert $A$ (Gate 1), invert $B$ (Gate 2), and feed both into Gate 3:
    \[ Y = \overline{\overline{A} + \overline{B}} = A \cdot B \]
    \item \textbf{Synthesis of NAND using NOR (4 Gates):}
    Generate AND (3 gates) and invert the output (Gate 4):
    \[ Y = \overline{A \cdot B} \]
    \item \textbf{Synthesis of XNOR using NOR (4 Gates):}
    Synthesizes XNOR natively in 4 NOR gates using dual topology to XOR.
    \item \textbf{Synthesis of XOR using NOR (5 Gates):}
    Invert the 4-gate XNOR output with a 5th NOR gate.
\end{enumerate}

\begin{table}[htbp]
\centering
\small
\caption{Minimum Gate Count for Universal Gate Logic Syntheses}
\label{tab:universal_gate_counts}
\begin{tabularx}{\textwidth}{l c c}
\toprule
\textbf{Target Logic Function} & \textbf{Minimum 2-Input NAND Gates} & \textbf{Minimum 2-Input NOR Gates} \\
\midrule
\textbf{NOT} & 1 & 1 \\
\textbf{AND} & 2 & 3 \\
\textbf{OR} & 3 & 2 \\
\textbf{NAND} & 1 & 4 \\
\textbf{NOR} & 4 & 1 \\
\textbf{XOR} & 4 & 5 \\
\textbf{XNOR} & 5 & 4 \\
\bottomrule
\end{tabularx}
\end{table}



\section{Combinational Logic Circuits}
\label{sec:combinational_logic}

Digital circuits are broadly classified into two grand architectural categories: \textbf{Combinational Circuits} and \textbf{Sequential Circuits}.

\subsection{Definition of Combinational Logic}
A \textbf{combinational circuit} is a connected arrangement of logic gates whose outputs at any given instant depend \textit{solely and strictly} upon the present combination of input values applied at that exact moment. It contains no internal storage elements (latches or flip-flops) and possesses no closed feedback paths from outputs back to inputs.

\begin{table}[htbp]
\centering
\small
\caption{Architectural Comparison: Combinational versus Sequential Logic}
\label{tab:comb_vs_seq}
\begin{tabularx}{\textwidth}{l X X}
\toprule
\textbf{Dimension} & \textbf{Combinational Circuit} & \textbf{Sequential Circuit} \\
\midrule
\textbf{Output Dependency} & Strictly depends on current inputs: $Y = f(X)$. & Depends on current inputs and past history (stored internal state). \\
\textbf{Memory Elements} & None. No latches, registers, or capacitors. & Requires memory elements (flip-flops, latches, RAM cells). \\
\textbf{Feedback Loops} & None (acyclic directed graph). & Closed feedback loops from outputs back to combinational inputs. \\
\textbf{Clock Synchronization} & Purely asynchronous; outputs change after propagation delay. & Clocked (synchronous) or unclocked (asynchronous state changes). \\
\textbf{Design Complexity} & Analyzed via truth tables, K-maps, Boolean algebra. & Analyzed via state tables, state diagrams, and excitation equations. \\
\textbf{Standard Examples} & Adders, subtractors, decoders, multiplexers, ALUs. & Counters, shift registers, CPU pipeline registers, RAM, FSMs. \\
\bottomrule
\end{tabularx}
\end{table}

\subsection{Systematic Combinational Design Workflow}
The formal synthesis of combinational circuits follows five structured phases:
\begin{enumerate}[label=\bfseries\arabic*.]
    \item \textbf{Specification:} Define input/output variable names, polarities, and required operational logic.
    \item \textbf{Truth Table Formulation:} Enumerate all $2^n$ possible input combinations and specify desired outputs.
    \item \textbf{Boolean Minimization:} Map minterms/maxterms into Karnaugh Maps or apply Quine-McCluskey tabulation to find the Minimal Sum of Products (SOP) or Product of Sums (POS).
    \item \textbf{Gate-Level Schematic Synthesis:} Map minimized equations into target gate libraries (e.g., standard CMOS NAND/NOR cells).
    \item \textbf{Timing and Glitch Verification:} Analyze propagation delay paths and verify the absence of static/dynamic timing hazards.
\end{enumerate}


\section{Half Adder Circuit Design}
\label{sec:half_adder}

\subsection{Functional Definition and Truth Table}
A \textbf{Half Adder} is an arithmetic combinational building block that computes the sum of two single-bit binary inputs ($A$ and $B$). It produces two binary outputs: a \textbf{Sum ($S$)} bit and a \textbf{Carry ($C$)} bit. It has no provision for receiving a carry bit from a preceding stage.

\begin{table}[htbp]
\centering
\small
\caption{Half Adder Truth Table}
\label{tab:half_adder_truth}
\begin{tabularx}{0.6\textwidth}{c c c c}
\toprule
\textbf{Input $A$} & \textbf{Input $B$} & \textbf{Sum ($S$)} & \textbf{Carry ($C$)} \\
\midrule
0 & 0 & 0 & 0 \\
0 & 1 & 1 & 0 \\
1 & 0 & 1 & 0 \\
1 & 1 & 0 & 1 \\
\bottomrule
\end{tabularx}
\end{table}

\subsection{Boolean Equations and Gate Implementation}
From Table~\ref{tab:half_adder_truth}, the canonical Boolean equations are:
\begin{equation}
\label{eq:half_adder_sum}
S = \overline{A}B + A\overline{B} = A \oplus B
\end{equation}
\begin{equation}
\label{eq:half_adder_carry}
C = A \cdot B
\end{equation}
\textbf{Hardware Realization:} A Half Adder requires exactly one 2-input XOR gate to generate the Sum bit, and one 2-input AND gate to generate the Carry bit. When synthesized using universal NAND gates, a Half Adder requires exactly \textbf{5 NAND gates}.


\section{Full Adder Circuit Design}
\label{sec:full_adder}

\subsection{Functional Definition and Truth Table}
A \textbf{Full Adder} computes the arithmetic sum of three single-bit inputs: two significant operand bits ($A$ and $B$) and an incoming carry bit ($C_{in}$) generated by a lower-order stage. It produces a \textbf{Sum ($S$)} and an outgoing \textbf{Carry ($C_{out}$)}.

\begin{table}[htbp]
\centering
\small
\caption{Full Adder Truth Table}
\label{tab:full_adder_truth}
\begin{tabularx}{0.7\textwidth}{c c c c c}
\toprule
\textbf{Input $A$} & \textbf{Input $B$} & \textbf{Carry In ($C_{in}$)} & \textbf{Sum ($S$)} & \textbf{Carry Out ($C_{out}$)} \\
\midrule
0 & 0 & 0 & 0 & 0 \\
0 & 0 & 1 & 1 & 0 \\
0 & 1 & 0 & 1 & 0 \\
0 & 1 & 1 & 0 & 1 \\
1 & 0 & 0 & 1 & 0 \\
1 & 0 & 1 & 0 & 1 \\
1 & 1 & 0 & 0 & 1 \\
1 & 1 & 1 & 1 & 1 \\
\bottomrule
\end{tabularx}
\end{table}

\subsection{Boolean Expressions and K-Map Simplification}
From Table~\ref{tab:full_adder_truth}, the Sum output is an odd parity function:
\begin{equation}
\label{eq:full_adder_sum}
S = \sum m(1, 2, 4, 7) = A \oplus B \oplus C_{in}
\end{equation}
The Carry Out is a majority function:
\begin{equation}
\label{eq:full_adder_cout}
C_{out} = \sum m(3, 5, 6, 7) = AB + BC_{in} + AC_{in}
\end{equation}
Factoring Eq.~\ref{eq:full_adder_cout} in terms of the Half Adder intermediate product:
\begin{equation}
\label{eq:full_adder_cout_factored}
C_{out} = AB + C_{in}(A \oplus B)
\end{equation}

\subsection{Synthesizing a Full Adder Using Two Half Adders}
A Full Adder can be constructed hierarchically using \textbf{two Half Adders (HA) and one OR gate}:
\begin{enumerate}[label=\bfseries\arabic*.]
    \item \textbf{First Half Adder ($HA_1$):} Receives $A$ and $B$.
    \[ S_1 = A \oplus B, \qquad C_1 = AB \]
    \item \textbf{Second Half Adder ($HA_2$):} Receives $S_1$ and $C_{in}$.
    \[ S = S_1 \oplus C_{in} = (A \oplus B) \oplus C_{in} \]
    \[ C_2 = S_1 \cdot C_{in} = (A \oplus B) \cdot C_{in} \]
    \item \textbf{Final OR Gate:} Generates $C_{out} = C_1 + C_2 = AB + C_{in}(A \oplus B)$.
\end{enumerate}
\textit{Universal NAND Count:} A Full Adder requires exactly \textbf{9 NAND gates}.

\begin{table}[htbp]
\centering
\small
\caption{Half Adder versus Full Adder Architectural Summary}
\label{tab:ha_vs_fa}
\begin{tabularx}{\textwidth}{l X X}
\toprule
\textbf{Characteristic} & \textbf{Half Adder} & \textbf{Full Adder} \\
\midrule
\textbf{Inputs} & 2 operand bits ($A, B$). & 3 bits ($A, B$, and carry-in $C_{in}$). \\
\textbf{Outputs} & 2 bits (Sum $S$, Carry $C$). & 2 bits (Sum $S$, Carry-out $C_{out}$). \\
\textbf{Cascadability} & Cannot be cascaded for multi-bit addition. & Cascadable indefinitely to form $n$-bit parallel adders. \\
\textbf{Hardware Cost} & 1 XOR, 1 AND (5 NAND gates). & 2 XOR, 2 AND, 1 OR (9 NAND gates). \\
\textbf{Propagation Delay} & 1 XOR gate delay ($t_{XOR}$). & 2 XOR gate delays for Sum; $t_{XOR} + t_{AND} + t_{OR}$ for Carry. \\
\bottomrule
\end{tabularx}
\end{table}


\section{Subtractors, Multi-Bit Adders, and Comparators}
\label{sec:arithmetic_circuits}

\subsection{Half and Full Subtractors}
Binary subtraction operates symmetrically to addition, substituting a \textbf{Borrow} signal for the Carry.
\begin{itemize}
    \item \textbf{Half Subtractor:} Computes $A - B$.
    \begin{equation}
    \label{eq:half_subtractor}
    \text{Difference } D = A \oplus B, \qquad \text{Borrow Out } B_{out} = \overline{A} \cdot B
    \end{equation}
    \item \textbf{Full Subtractor:} Computes $A - B - B_{in}$.
    \begin{equation}
    \label{eq:full_subtractor}
    D = A \oplus B \oplus B_{in}, \qquad B_{out} = \overline{A}B + \overline{(A \oplus B)}B_{in}
    \end{equation}
\end{itemize}

\subsection{The 4-Bit Parallel Adder/Subtractor Circuit}
Rather than building separate subtractor hardware, processors utilize an $n$-bit parallel adder augmented with 2-input XOR gates as \textbf{controlled inverters}:
\begin{itemize}
    \item Each bit $B_i$ is fed into an XOR gate with control line $M$ ($B_i \oplus M$).
    \item The control line $M$ is wired directly to the carry-in of the LSB ($C_0 = M$).
    \item \textbf{When $M = 0$ (Addition Mode):} $B_i \oplus 0 = B_i$ and $C_0 = 0$. The circuit computes $A + B + 0 = A + B$.
    \item \textbf{When $M = 1$ (Subtraction Mode):} $B_i \oplus 1 = \overline{B_i}$ and $C_0 = 1$. The circuit computes $A + \overline{B} + 1 = A - B$ in two's complement.
\end{itemize}

\subsection{Ripple Carry Adder (RCA) versus Carry Lookahead Adder (CLA)}
\begin{enumerate}[label=\bfseries\arabic*.]
    \item \textbf{Ripple Carry Adder (RCA):} An $n$-bit RCA cascades $n$ Full Adders such that $C_{out}$ of stage $i$ feeds $C_{in}$ of stage $i+1$.
    \begin{itemize}
        \item \textit{Delay Bottleneck:} The MSB cannot compute its final Sum or Carry until the carry has ripple-propagated linearly through all $n$ stages.
        \begin{equation}
        \label{eq:rca_delay}
        T_{\text{RCA}} = n \cdot t_{\text{carry}} = O(n)
        \end{equation}
        For a 64-bit adder with $t_{\text{carry}} = 2\text{ gate delays}$, total carry latency is $128\text{ gate delays}$, which limits CPU clock speed.
    \end{itemize}
    \item \textbf{Carry Lookahead Adder (CLA):} Solves the ripple latency bottleneck by computing carry signals in parallel using two auxiliary Boolean definitions:
    \begin{itemize}
        \item \textbf{Carry Generate ($G_i$):} Stage $i$ generates a carry unconditionally if both inputs are $1$:
        \[ G_i = A_i \cdot B_i \]
        \item \textbf{Carry Propagate ($P_i$):} Stage $i$ propagates an incoming carry if at least one input is $1$:
        \[ P_i = A_i \oplus B_i \quad (\text{or } A_i + B_i) \]
    \end{itemize}
    The carry out of any stage $i$ is formulated recursively in terms of $G_i$, $P_i$, and initial carry $C_0$:
    \begin{align*}
    C_1 &= G_0 + P_0 C_0 \\
    C_2 &= G_1 + P_1 C_1 = G_1 + P_1 G_0 + P_1 P_0 C_0 \\
    C_3 &= G_2 + P_2 G_1 + P_2 P_1 G_0 + P_2 P_1 P_0 C_0 \\
    C_4 &= G_3 + P_3 G_2 + P_3 P_2 G_1 + P_3 P_2 P_1 G_0 + P_3 P_2 P_1 P_0 C_0
    \end{align*}
    \textit{Latency Advantage:} Every carry bit is evaluated within \textbf{2 gate delays} (one level of ANDs, one level of OR) regardless of adder width. The asymptotic latency of CLA is $O(\log n)$ when constructed hierarchically (Block CLA).
\end{enumerate}

\subsection{Magnitude Comparator}
A \textbf{Magnitude Comparator} compares two binary words $A$ and $B$ and asserts one of three mutually exclusive output flags: $A > B$, $A = B$, or $A < B$.
\begin{itemize}
    \item For 1-bit inputs:
    \[ (A > B) = A\overline{B}, \qquad (A = B) = \overline{A \oplus B} = A \odot B, \qquad (A < B) = \overline{A}B \]
    \item For $n$-bit words, equality requires all bit positions to match simultaneously:
    \[ (A = B) = x_{n-1} \cdot x_{n-2} \cdots x_1 \cdot x_0 \quad \text{where } x_i = \overline{A_i \oplus B_i} \]
\end{itemize}


\section{Encoders, Priority Encoders, and Decoders}
\label{sec:encoders_decoders}

\subsection{Binary Encoders}
An \textbf{Encoder} has $2^n$ (or fewer) input lines and $n$ output lines. It generates the unique $n$-bit binary code corresponding to whichever single input line is currently asserted high.
\begin{itemize}
    \item \textit{Limitation of Simple Encoders:} If two or more inputs are asserted simultaneously, or if no inputs are active, the circuit generates an erroneous, undefined binary output.
\end{itemize}

\subsection{Priority Encoders}
A \textbf{Priority Encoder} resolves simultaneous inputs by establishing a strict priority hierarchy (typically with the highest index having highest priority). If multiple inputs are asserted, only the one with the highest priority is encoded. A \textbf{Valid ($V$)} output bit is provided to distinguish an active $D_0$ input from an all-idle state.

\begin{table}[htbp]
\centering
\small
\caption{4-to-2 Priority Encoder Truth Table (Higher Index = Higher Priority)}
\label{tab:priority_encoder}
\begin{tabularx}{\textwidth}{c c c c c c c}
\toprule
$D_3$ & $D_2$ & $D_1$ & $D_0$ & $Y_1$ & $Y_0$ & \textbf{Valid ($V$)} \\
\midrule
0 & 0 & 0 & 0 & X & X & \textbf{0} \\
0 & 0 & 0 & 1 & 0 & 0 & \textbf{1} \\
0 & 0 & 1 & X & 0 & 1 & \textbf{1} \\
0 & 1 & X & X & 1 & 0 & \textbf{1} \\
1 & X & X & X & 1 & 1 & \textbf{1} \\
\bottomrule
\end{tabularx}
\end{table}

\textbf{Minimized Logic Equations:}
\begin{equation}
\label{eq:priority_encoder_eqs}
Y_1 = D_3 + D_2, \qquad Y_0 = D_3 + \overline{D_2} D_1, \qquad V = D_3 + D_2 + D_1 + D_0
\end{equation}

\subsection{Binary Decoders}
A \textbf{Decoder} performs the inverse operation of an encoder. It has $n$ input lines and up to $2^n$ output lines. Exactly one output line is activated for each of the $2^n$ unique input minterm combinations.
\begin{itemize}
    \item \textbf{Enable Control ($E$):} Decoders include an Enable input. When $E=0$, all output lines are disabled (deasserted).
    \item \textbf{Minterm Generator Property:} Because a decoder's outputs correspond directly to minterms $m_0, m_1, \dots, m_{2^n-1}$, \textbf{any arbitrary Boolean function can be implemented by simply ORing the appropriate decoder outputs}.
    \item \textbf{Memory Address Decoding:} In microprocessor memory subsystems, the high-order address lines feed a decoder to generate chip-select ($\overline{CS}$) signals for individual RAM/ROM IC banks.
\end{itemize}


\section{Multiplexers (MUX) and Demultiplexers (DEMUX)}
\label{sec:mux_demux}

\subsection{Multiplexer (Data Selector)}
A \textbf{Multiplexer (MUX)} has $2^n$ data inputs ($I_0, \dots, I_{2^n-1}$), $n$ select lines ($S_{n-1}, \dots, S_0$), and a single output line ($Y$). It routes the data from one selected input to the output based on the binary value of the select lines.

\begin{equation}
\label{eq:mux_general}
Y = \sum_{k=0}^{2^n-1} m_k \cdot I_k
\end{equation}
where $m_k$ is the minterm formed by the $n$ select lines.

\textbf{4-to-1 Multiplexer Equation:}
\begin{equation}
\label{eq:mux_4to1}
Y = \overline{S_1}\,\overline{S_0} I_0 + \overline{S_1} S_0 I_1 + S_1 \overline{S_0} I_2 + S_1 S_0 I_3
\end{equation}

\subsection{Synthesizing Arbitrary Logic Functions Using Multiplexers}
Multiplexers are universal combinational logic building blocks.
\begin{enumerate}[label=\bfseries\arabic*.]
    \item \textbf{Using a $2^n$-to-1 MUX for an $n$-Variable Function:}
    Connect the $n$ Boolean variables to the select lines. Connect the data inputs $I_k$ directly to $1$ or $0$ corresponding to the function's truth table values.
    \item \textbf{Using a $2^{n-1}$-to-1 MUX for an $n$-Variable Function:}
    Connect $n-1$ variables to the select lines. The data inputs $I_k$ are assigned to the remaining variable ($C$), its complement ($\overline{C}$), $0$, or $1$ via algebraic residues.
\end{enumerate}

\subsection{Demultiplexer (Data Distributor)}
A \textbf{Demultiplexer (DEMUX)} takes a single input data line $D$ and routes it to one of $2^n$ possible output lines ($Y_0, \dots, Y_{2^n-1}$) based on $n$ select lines.
\begin{itemize}
    \item \textbf{Duality with Decoders:} A decoder with an active-high Enable line operates identically to a demultiplexer whose input data $D$ is fed into the Enable pin.
\end{itemize}


\section{Tri-State Buffers and Bus Contention}
\label{sec:tristate_buffers}

\subsection{The Three Operational Logic States}
Standard digital gates drive outputs to either Logic Low ($0\text{ V}$) or Logic High ($+V_{DD}$). A \textbf{Tri-State Buffer} introduces a third electrical state:
\begin{enumerate}[label=\bfseries\arabic*.]
    \item \textbf{Logic 0:} Low-impedance path to Ground.
    \item \textbf{Logic 1:} Low-impedance path to Supply Voltage ($V_{DD}$).
    \item \textbf{High-Impedance State (Hi-Z):} Both the internal pull-up pMOS and pull-down nMOS output transistors are driven into cut-off. The output pin is electrically disconnected from the bus, floating with infinite effective resistance ($>10\text{ M}\Omega$).
\end{enumerate}

\subsection{Preventing Bus Contention on Shared System Buses}
In microprocessor architectures, multiple hardware modules (CPU, DMA controller, multiple RAM chips) connect physically to the same shared parallel data bus.
\begin{itemize}
    \item \textbf{Bus Contention (Deadly Short Circuit):} If two different drivers attempt to drive the same physical line simultaneously—one asserting Logic 1 ($5\text{ V}$) and the other asserting Logic 0 ($0\text{ V}$)—a direct low-resistance short circuit is formed between $V_{DD}$ and Ground. This draws destructive surge currents ($>100\text{ mA}$), causing permanent thermal silicon burnout and invalid signal levels.
    \item \textbf{Tri-State Bus Governance:} Address decoders guarantee that at any given instant, \textbf{only one driver's tri-state enable line is asserted}. All other dormant drivers are held in the High-Impedance (Hi-Z) state, allowing the single active device to transmit clean, uncorrupted binary data across the bus.
\end{itemize}



\section{Karnaugh Maps (K-Maps), Canonical Forms, and Boolean Simplification}
\label{sec:karnaugh_maps}

The \textbf{Karnaugh Map (K-Map)}, invented by Maurice Karnaugh in 1953, is a graphical technique for minimizing Boolean algebraic expressions. It organizes a truth table into a multi-dimensional grid where adjacent cells differ in exactly one binary variable, enabling visual identification and elimination of redundant literals via the algebraic identity $A + \overline{A} = 1$.

\subsection{Canonical Boolean Forms: Minterms and Maxterms}
Any arbitrary switching function $F(A, B, C, \dots)$ can be represented in two canonical algebraic formats:

\begin{enumerate}[label=\bfseries\arabic*.]
    \item \textbf{Minterm (Standard Product Term):} A product of all $n$ literals in either true or complemented form that evaluates to $1$ for exactly one row of the truth table.
    \begin{itemize}
        \item Variable convention: $A = 1$ is uncomplemented ($A$), $A = 0$ is complemented ($\overline{A}$).
        \item Example for 3 variables: Row $101_2$ corresponds to minterm $m_5 = A \overline{B} C$.
        \item \textbf{Canonical Sum of Products (SOP):} Formed by ORing all minterms where $F = 1$:
        \begin{equation}
        \label{eq:canonical_sop}
        F(A, B, C) = \sum m(1, 3, 5, 7) = \overline{A}\,\overline{B}C + \overline{A}BC + A\overline{B}C + ABC
        \end{equation}
    \end{itemize}
    \item \textbf{Maxterm (Standard Sum Term):} A sum of all $n$ literals in either true or complemented form that evaluates to $0$ for exactly one row of the truth table.
    \begin{itemize}
        \item Variable convention: $A = 0$ is uncomplemented ($A$), $A = 1$ is complemented ($\overline{A}$).
        \item Example for 3 variables: Row $101_2$ corresponds to maxterm $M_5 = \overline{A} + B + \overline{C}$.
        \item \textbf{Canonical Product of Sums (POS):} Formed by ANDing all maxterms where $F = 0$:
        \begin{equation}
        \label{eq:canonical_pos}
        F(A, B, C) = \prod M(0, 2, 4, 6) = (A + B + C)(A + \overline{B} + C)(\overline{A} + B + C)(\overline{A} + \overline{B} + C)
        \end{equation}
    \end{itemize}
\end{enumerate}

\begin{table}[htbp]
\centering
\small
\caption{Canonical Forms Comparison: Sum of Products (SOP) versus Product of Sums (POS)}
\label{tab:sop_vs_pos}
\begin{tabularx}{\textwidth}{l X X}
\toprule
\textbf{Dimension} & \textbf{Sum of Products (SOP)} & \textbf{Product of Sums (POS)} \\
\midrule
\textbf{Constituent Terms} & Minterms ($m_i$): AND product terms. & Maxterms ($M_i$): OR sum terms. \\
\textbf{Truth Table Focus} & Rows where output $F = 1$. & Rows where output $F = 0$. \\
\textbf{Algebraic Form} & AND gates feeding a final OR gate. & OR gates feeding a final AND gate. \\
\textbf{K-Map Grouping} & Groups of adjacent 1s. & Groups of adjacent 0s. \\
\textbf{Active Variable Polarities} & $0 = \overline{A}$ (complemented), $1 = A$ (true). & $0 = A$ (true), $1 = \overline{A}$ (complemented). \\
\textbf{Universal Gate Mapping} & Natively synthesized via 2-level NAND-NAND. & Natively synthesized via 2-level NOR-NOR. \\
\bottomrule
\end{tabularx}
\end{table}

\subsection{Why Gray Code Is Mandatory in K-Maps}
A K-map organizes $2^n$ cells into a grid where the row and column coordinates follow a \textbf{reflected Gray Code sequence} ($00, 01, 11, 10$ instead of natural binary $00, 01, 10, 11$).
\begin{itemize}
    \item \textbf{Unit-Distance Property:} Because adjacent cells differ by exactly one binary bit, any pair of adjacent cells representing $m_1$ and $m_2$ shares $n-1$ identical literals and one complemented pair ($X$ and $\overline{X}$).
    \item Grouping them algebraically applies the complement identity:
    \[ W \cdot X + W \cdot \overline{X} = W \cdot (X + \overline{X}) = W \cdot (1) = W \]
    This eliminates the varying variable $X$ from the product term entirely.
\end{itemize}

\subsection{K-Map Structural Topologies and Cell Adjacency}
An $n$-variable K-map contains $2^n$ cells:
\begin{enumerate}[label=\bfseries\arabic*.]
    \item \textbf{Two-Variable K-Map (4 Cells):}
    \begin{center}
    \small
    \begin{tabular}{|c|c|c|}
    \hline
    $A \backslash B$ & \textbf{0} & \textbf{1} \\
    \hline
    \textbf{0} & $m_0$ ($\overline{A}\,\overline{B}$) & $m_1$ ($\overline{A}B$) \\
    \hline
    \textbf{1} & $m_2$ ($A\overline{B}$) & $m_3$ ($AB$) \\
    \hline
    \end{tabular}
    \end{center}
    \item \textbf{Three-Variable K-Map (8 Cells):}
    \begin{center}
    \small
    \begin{tabular}{|c|c|c|c|c|}
    \hline
    $A \backslash BC$ & \textbf{00} & \textbf{01} & \textbf{11} & \textbf{10} \\
    \hline
    \textbf{0} & $m_0$ & $m_1$ & $m_3$ & $m_2$ \\
    \hline
    \textbf{1} & $m_4$ & $m_5$ & $m_7$ & $m_6$ \\
    \hline
    \end{tabular}
    \end{center}
    \textit{Note cell order:} The third column is $11$ ($m_3, m_7$) and the fourth column is $10$ ($m_2, m_6$).
    \item \textbf{Four-Variable K-Map (16 Cells):}
    \begin{center}
    \small
    \begin{tabular}{|c|c|c|c|c|}
    \hline
    $AB \backslash CD$ & \textbf{00} & \textbf{01} & \textbf{11} & \textbf{10} \\
    \hline
    \textbf{00} & $m_0$ & $m_1$ & $m_3$ & $m_2$ \\
    \hline
    \textbf{01} & $m_4$ & $m_5$ & $m_7$ & $m_6$ \\
    \hline
    \textbf{11} & $m_{12}$ & $m_{13}$ & $m_{15}$ & $m_{14}$ \\
    \hline
    \textbf{10} & $m_8$ & $m_9$ & $m_{11}$ & $m_{10}$ \\
    \hline
    \end{tabular}
    \end{center}
    \textit{Note row order:} The third row is $11$ ($m_{12}, m_{13}, m_{15}, m_{14}$) and the fourth row is $10$ ($m_8, m_9, m_{11}, m_{10}$).
\end{enumerate}

\subsection{Rigorous K-Map Grouping Rules}
To obtain the mathematically minimal expression, groupings must strictly adhere to the following laws:
\begin{enumerate}[label=\bfseries\arabic*.]
    \item \textbf{Power of Two Rule:} Every group must contain exactly $2^k$ cells ($1, 2, 4, 8, \text{ or } 16$ cells). Groups of 3, 5, 6, or 7 cells are strictly invalid.
    \begin{itemize}
        \item A group of 1 cell represents an $n$-literal term (no variable eliminated).
        \item A group of 2 cells (a \textbf{Pair}) eliminates $1$ variable.
        \item A group of 4 cells (a \textbf{Quad}) eliminates $2$ variables.
        \item A group of 8 cells (an \textbf{Octet}) eliminates $3$ variables.
        \item A group of 16 cells eliminates all $4$ variables (evaluates to constant $1$).
    \end{itemize}
    \item \textbf{Geometric Rectangularity:} Groups must form contiguous rectangles or squares. Diagonal, L-shaped, or triangular groupings are physically invalid.
    \item \textbf{Toroidal Wrap-Around Adjacency:} The K-map grid wraps around horizontally and vertically:
    \begin{itemize}
        \item The extreme left column ($00$) is adjacent to the extreme right column ($10$).
        \item The top row ($00$) is adjacent to the bottom row ($10$).
        \item \textbf{The Four-Corner Quad:} Cells $m_0, m_2, m_8, m_{10}$ form a valid 4-cell quad.
    \end{itemize}
    \item \textbf{Maximal Sizing:} Always form the largest possible group first. Larger groups eliminate more variables, yielding simpler, lower-cost logic gates.
    \item \textbf{Redundancy Elimination:} Overlapping groups are permitted and encouraged to enlarge groups, but \textbf{a group is completely redundant and must be eliminated} if all of its constituent $1$s are already covered by other groups.
\end{enumerate}

\subsection{Don't-Care Conditions ($d$ or $X$)}
In many digital systems, certain input combinations can never physically occur during normal operation (e.g., in BCD arithmetic, inputs $1010_2$ through $1111_2$ are impossible). The circuit designer does not care whether the output evaluates to $0$ or $1$ for these states.
\begin{itemize}
    \item \textbf{Strategic Exploitation:} A Don't-Care ($X$) cell may be treated as a $1$ if doing so enables the formation of a larger group (reducing gate count), or treated as a $0$ if it does not help simplify the logic.
    \item \textbf{Crucial Examination Rule:} \textbf{Never} form a group consisting exclusively of Don't-Cares. A group must contain at least one essential $1$.
\end{itemize}

\subsection{Worked K-Map Simplification Examples}
\textbf{Example 1 (3-Variable with Wrap-Around):} Minimize $F(A, B, C) = \sum m(0, 2, 4, 6)$.
\begin{center}
\small
\begin{tabular}{|c|c|c|c|c|}
\hline
$A \backslash BC$ & \textbf{00} & \textbf{01} & \textbf{11} & \textbf{10} \\
\hline
\textbf{0} & \textbf{1} & 0 & 0 & \textbf{1} \\
\hline
\textbf{1} & \textbf{1} & 0 & 0 & \textbf{1} \\
\hline
\end{tabular}
\end{center}
\textit{Minimization:} Columns $00$ and $10$ across rows $0$ and $1$ wrap around to form a single \textbf{Quad of 4 cells} ($m_0, m_2, m_4, m_6$).
\begin{itemize}
    \item For rows $0$ and $1$: $A$ changes from $0$ to $1 \implies A$ is eliminated.
    \item For columns $00$ ($B=0, C=0$) and $10$ ($B=1, C=0$): $B$ changes ($0 \to 1$) $\implies B$ is eliminated; $C=0$ remains constant.
    \item Minimized Result: $\mathbf{F = \overline{C}}$.
\end{itemize}

\textbf{Example 2 (4-Variable with Don't-Care):}
Minimize $F(A, B, C, D) = \sum m(1, 3, 7, 11, 15) + d(0, 2, 5)$.
\begin{center}
\small
\begin{tabular}{|c|c|c|c|c|}
\hline
$AB \backslash CD$ & \textbf{00} & \textbf{01} & \textbf{11} & \textbf{10} \\
\hline
\textbf{00} & $X$ & \textbf{1} & \textbf{1} & $X$ \\
\hline
\textbf{01} & 0 & $X$ & \textbf{1} & 0 \\
\hline
\textbf{11} & 0 & 0 & \textbf{1} & 0 \\
\hline
\textbf{10} & 0 & 0 & \textbf{1} & 0 \\
\hline
\end{tabular}
\end{center}
\textit{Minimization:}
\begin{enumerate}[label=\bfseries\alph*.]
    \item Column $CD = 11$ contains $1$s in all four rows ($m_3, m_7, m_{15}, m_{11}$). This forms an \textbf{Octet-like Quad of 4 cells} eliminating $A$ and $B$, giving the term $\mathbf{CD}$.
    \item Minterm $m_1$ ($0001$) remains. By including don't-cares $d_0$ ($0000$), $d_2$ ($0010$), and $d_5$ ($0101$) along with $m_3$ ($0011$), we can form a \textbf{Quad} in row $00$ and row $01$ under columns $00$ and $01$:
    Cells $\{m_0, m_1, m_4(\text{zero, cannot use})\}$.
    Instead, look at row $00$: cells $m_0(X), m_1(1), m_3(1), m_2(X)$ span the entire top row!
    This forms a complete \textbf{Row Quad of 4 cells} ($m_0, m_1, m_3, m_2$). In this row, $A=0, B=0$ are constant, while $C$ and $D$ vary across all four combinations. This yields the term $\mathbf{\overline{A}\,\overline{B}}$.
    \item Minimized SOP Expression: $\mathbf{F = CD + \overline{A}\,\overline{B}}$.
\end{enumerate}

\subsection{Essential Examination Traps in K-Maps}
\begin{enumerate}[label=\bfseries\arabic*.]
    \item \textbf{Inverting Gray Code Coordinates:} Labeling columns as $00, 01, 10, 11$ instead of $00, 01, 11, 10$. This destroys spatial cell adjacency and produces completely wrong groupings.
    \item \textbf{Non-Power of Two Groups:} Grouping 3 or 6 adjacent cells. Group sizes must always strictly be $1, 2, 4, 8, \text{ or } 16$.
    \item \textbf{Overlooking Corner Adjacency:} Missing the 4-corner quad ($m_0, m_2, m_8, m_{10}$) in 4-variable maps, which simplifies to $\overline{B}\,\overline{D}$.
    \item \textbf{Retaining Redundant Implicants:} Keeping an isolated pair or quad whose elements are already covered by larger groups. Every retained group must contain at least one \textbf{Essential Prime Implicant} ($1$ covered by no other group).
\end{enumerate}



\section{Sequential Logic Fundamentals}
\label{sec:sequential_logic}

Unlike purely combinational circuits, \textbf{Sequential Logic Circuits} incorporate internal feedback paths and bistable storage elements. Their present outputs are a function of both the current external inputs and the historical sequence of past inputs (the internal stored state).

\subsection{Synchronous versus Asynchronous Sequential Systems}
\begin{enumerate}[label=\bfseries\arabic*.]
    \item \textbf{Synchronous Sequential Circuits:} State transitions occur simultaneously across all storage elements, regulated by a periodic global timing signal called the \textbf{System Clock}.
    \begin{itemize}
        \item \textit{Advantages:} Highly predictable, immune to internal propagation race conditions, and supported by automated static timing analysis (STA) in EDA synthesis tools.
        \item \textit{Predominance:} Represents $>99\%$ of modern microprocessor and ASIC architectures.
    \end{itemize}
    \item \textbf{Asynchronous Sequential Circuits:} State transitions are event-driven, triggered directly by changes in external input levels without a coordinating clock signal.
    \begin{itemize}
        \item \textit{Challenges:} Highly vulnerable to critical race conditions, dynamic glitches, and timing hazards; requires complex differential delay modeling.
    \end{itemize}
\end{enumerate}


\section{Latches: Unclocked Bistable Storage Elements}
\label{sec:latches}

A \textbf{Latch} is a level-sensitive bistable multivibrator capable of storing one bit of binary data. Its outputs follow input changes continuously as long as the enable signal remains asserted high (transparent mode).

\subsection{The Cross-Coupled NOR SR Latch}
Constructed using two cross-coupled NOR gates where each gate's output feeds the second gate's input:
\begin{itemize}
    \item \textbf{Set ($S = 1, R = 0$):} Forces output $Q = 1$ and $\overline{Q} = 0$.
    \item \textbf{Reset ($S = 0, R = 1$):} Forces output $Q = 0$ and $\overline{Q} = 1$.
    \item \textbf{Hold / Quiescent ($S = 0, R = 0$):} Feedback maintains the previous stored state ($Q^+ = Q$).
    \item \textbf{Forbidden / Invalid State ($S = 1, R = 1$):} Both NOR outputs are driven to $0$ simultaneously ($Q = 0, \overline{Q} = 0$), violating the complementary definition of $Q$ and $\overline{Q}$. If both inputs transition to $0$ simultaneously, the latch enters an unpredictable metastable race.
\end{itemize}

\subsection{The Cross-Coupled NAND $\overline{\text{S}}\,\overline{\text{R}}$ Latch}
Constructed using cross-coupled NAND gates with active-low inputs:
\begin{itemize}
    \item Inputs are asserted with logic $0$: $\overline{S}=0$ sets $Q=1$; $\overline{R}=0$ resets $Q=0$.
    \item The quiescent hold state is $\overline{S}=1, \overline{R}=1$.
    \item The forbidden state is $\overline{S}=0, \overline{R}=0$.
\end{itemize}

\subsection{Gated D Latch (Transparent Latch)}
The Gated D Latch eliminates the invalid input state by inserting an inverter between the Set and Reset inputs ($S = D, R = \overline{D}$), controlled by an Enable pin ($E$):
\begin{itemize}
    \item When $E = 1$ (Transparent): Output $Q$ follows input $D$ directly ($Q^+ = D$).
    \item When $E = 0$ (Opaque / Latched): Output holds its last valid state, immune to fluctuations on $D$.
\end{itemize}


\section{Flip-Flops, Characteristic Equations, and Excitation Tables}
\label{sec:flip_flops}

A \textbf{Flip-Flop} is an edge-triggered binary storage cell that samples its inputs and updates its output state \textit{only during a discrete transition} (rising edge $0 \to 1$ or falling edge $1 \to 0$) of the clock signal.

\subsection{Comparative Analysis: Latch versus Flip-Flop}
\begin{table}[htbp]
\centering
\small
\caption{Architectural Distinctions: Latches versus Flip-Flops}
\label{tab:latch_vs_flipflop}
\begin{tabularx}{\textwidth}{l X X}
\toprule
\textbf{Dimension} & \textbf{Latch} & \textbf{Flip-Flop} \\
\midrule
\textbf{Triggering Mechanism} & Level-sensitive (active during high or low clock phase). & Edge-triggered (active during rising or falling voltage transition). \\
\textbf{Data Transparency} & Fully transparent while enable is high. & Completely opaque; samples inputs only at instantaneous clock edges. \\
\textbf{Race Conditions} & Highly vulnerable to race-around condition if feedback exists. & Immune to race-around; output updates after edge has passed. \\
\textbf{Hardware Complexity} & Simple (typically 4--8 transistors). & Requires two cascaded latches (Master-Slave) or 6 NAND gates (12--20 transistors). \\
\textbf{Primary Role} & High-speed asynchronous FIFO buffers, clock gating cells. & CPU pipeline registers, synchronous counters, finite state machines. \\
\bottomrule
\end{tabularx}
\end{table}

\subsection{The Four Standard Flip-Flop Types}
\begin{enumerate}[label=\bfseries\arabic*.]
    \item \textbf{SR Flip-Flop:} Clocked version of the SR latch.
    \begin{itemize}
        \item Characteristic Equation:
        \begin{equation}
        \label{eq:sr_char}
        Q^+ = S + \overline{R} Q \quad (\text{Constraint: } S \cdot R = 0)
        \end{equation}
    \end{itemize}
    \item \textbf{D Flip-Flop (Data / Delay Flip-Flop):} Captures the input $D$ at the clock edge and delays it to the output by exactly one clock cycle.
    \begin{itemize}
        \item Characteristic Equation:
        \begin{equation}
        \label{eq:d_char}
        Q^+ = D
        \end{equation}
    \end{itemize}
    \item \textbf{JK Flip-Flop:} Overcomes the invalid state of the SR flip-flop by feeding internal outputs back to the input gating circuitry. When $J=K=1$, the output \textbf{toggles} ($Q^+ = \overline{Q}$).
    \begin{itemize}
        \item Characteristic Equation:
        \begin{equation}
        \label{eq:jk_char}
        Q^+ = J \overline{Q} + \overline{K} Q
        \end{equation}
    \end{itemize}
    \item \textbf{T Flip-Flop (Toggle Flip-Flop):} Formed by tying the $J$ and $K$ inputs of a JK flip-flop together ($J = K = T$). When $T=0$, it holds state ($Q^+ = Q$). When $T=1$, it inverts state ($Q^+ = \overline{Q}$).
    \begin{itemize}
        \item Characteristic Equation:
        \begin{equation}
        \label{eq:t_char}
        Q^+ = T \oplus Q = T \overline{Q} + \overline{T} Q
        \end{equation}
    \end{itemize}
\end{enumerate}

\subsection{Characteristic Tables and Excitation Tables}
Digital synthesis engineers use two distinct tabular representations:
\begin{itemize}
    \item \textbf{Characteristic Table:} Defines next state $Q^+$ as a function of current inputs and current state $Q$. Used for circuit \textit{analysis}.
    \item \textbf{Excitation Table:} Defines the required input values needed to force a desired state transition from present state $Q$ to next state $Q^+$. Used for circuit \textit{synthesis}.
\end{itemize}

\begin{table}[htbp]
\centering
\small
\caption{Comprehensive Master Excitation Table for All Flip-Flop Architectures}
\label{tab:master_excitation}
\begin{tabularx}{\textwidth}{c c c c c c c c}
\toprule
\textbf{Present State $Q$} & \textbf{Next State $Q^+$} & \textbf{SR: $S$} & \textbf{SR: $R$} & \textbf{JK: $J$} & \textbf{JK: $K$} & \textbf{D: $D$} & \textbf{T: $T$} \\
\midrule
0 & 0 & 0 & X & 0 & X & 0 & 0 \\
0 & 1 & 1 & 0 & 1 & X & 1 & 1 \\
1 & 0 & 0 & 1 & X & 1 & 0 & 1 \\
1 & 1 & X & 0 & X & 0 & 1 & 0 \\
\bottomrule
\end{tabularx}
\end{table}


\section{Timing Parameters, Metastability, and the Race-Around Condition}
\label{sec:flipflop_timing}

\subsection{Critical Dynamic Timing Specifications}
Reliable state transitions in clocked flip-flops mandate adherence to three strict temporal boundaries:
\begin{enumerate}[label=\bfseries\arabic*.]
    \item \textbf{Setup Time ($t_{\text{setup}}$):} The minimum interval of time immediately \textit{prior to the active clock transition} during which data inputs must remain completely stable and static.
    \item \textbf{Hold Time ($t_{\text{hold}}$):} The minimum interval of time immediately \textit{following the active clock transition} during which data inputs must continue to be held stable.
    \item \textbf{Clock-to-Q Propagation Delay ($t_{cq}$):} The elapsed physical latency between the triggering clock edge arriving at the clock pin and the valid output data appearing at pin $Q$.
\end{enumerate}

\begin{equation}
\label{eq:max_clock_freq}
T_{\text{clock\_min}} \ge t_{cq} + t_{\text{comb\_max}} + t_{\text{setup}} \quad \implies \quad f_{\text{max}} = \frac{1}{T_{\text{clock\_min}}}
\end{equation}

\subsection{Metastability}
If an asynchronous input violates either setup or hold time requirements (e.g., changes logic states directly inside the forbidden aperture $[t_{\text{clk}} - t_{\text{setup}}, t_{\text{clk}} + t_{\text{hold}}]$), the internal regenerative feedback transistors become balanced at an unstable intermediate voltage between $V_{OH}$ and $V_{OL}$.
\begin{itemize}
    \item The output will oscillate or hover at an indeterminate voltage for an unpredictable duration before settling to a random binary state ($0$ or $1$).
    \item \textbf{Mitigation:} Multi-stage synchronizers (cascading two or three D flip-flops clocked by the destination domain) exponentially reduce the Mean Time Between Failures (MTBF):
    \begin{equation}
    \label{eq:mtbf}
    \text{MTBF} = \frac{e^{t_{\text{resolve}} / \tau}}{T_0 \cdot f_{\text{clk}} \cdot f_{\text{data}}}
    \end{equation}
\end{itemize}

\subsection{The Race-Around Condition and the Master-Slave Solution}
\textbf{The Race-Around Problem:} In a level-triggered JK flip-flop with inputs $J = 1$ and $K = 1$, the output toggles every propagation delay time $t_{pd}$.
\begin{itemize}
    \item If the clock pulse width ($t_p$) is greater than the internal gate propagation delay ($t_{pd}$), the output will toggle repeatedly and continuously back and forth between $0$ and $1$ throughout the duration of the high clock phase:
    \begin{equation}
    \label{eq:race_around_condition}
    t_p > t_{pd}
    \end{equation}
    \item When the clock drops low, the final resting state of $Q$ is completely unpredictable.
\end{itemize}

\textbf{Solutions to Race-Around:}
\begin{enumerate}[label=\bfseries\arabic*.]
    \item \textbf{True Edge-Triggering:} Use edge-triggered flip-flop topologies (such as the 6-NAND positive-edge triggered D flip-flop) where the active sampling aperture is restricted to a few picoseconds.
    \item \textbf{Master-Slave JK Architecture:} Cascades two distinct latches driven by complementary clocks:
    \begin{itemize}
        \item \textit{Master Latch:} Enabled when $\text{CLK} = 1$. Samples inputs $J$ and $K$, while the Slave is disabled.
        \item \textit{Slave Latch:} Enabled when $\text{CLK} = 0$ via an inverter. Transfers the Master's state to final output $Q$, while the Master is disabled.
        \item Because the Master and Slave are never active simultaneously, no continuous closed feedback loop can exist, completely eliminating race-around.
    \end{itemize}
\end{enumerate}


\section{Asynchronous and Synchronous Counters}
\label{sec:counters}

A \textbf{Counter} is a sequential state register that cycles through a prescribed sequence of binary states upon receiving clock pulses. The number of unique states in the cycle is the counter's \textbf{Modulus (MOD)}. An $n$-flip-flop counter can achieve a maximum modulus of $\text{MOD} \le 2^n$.

\subsection{Asynchronous (Ripple) Counters}
In an \textbf{Asynchronous Counter}, the external clock is applied \textit{only to the LSB flip-flop}. The output $Q$ (or $\overline{Q}$) of each preceding flip-flop serves as the clock input for the subsequent higher-order stage.
\begin{itemize}
    \item \textit{Ripple Propagation Delay:} Because each stage must wait for the preceding stage to settle, cumulative propagation delay accumulates linearly:
    \begin{equation}
    \label{eq:ripple_delay}
    T_{\text{total\_delay}} = n \cdot t_{pd}
    \end{equation}
    \item \textit{Frequency Ceiling:} To avoid counting errors, the clock period must satisfy $T_{\text{clock}} \ge n \cdot t_{pd}$, severely restricting high-speed operation.
    \item \textit{Decoding Glitches:} Unequal stage switching times create false intermediate transient states, causing spurious spikes in combinational decoders.
\end{itemize}

\subsection{Synchronous Counters}
In a \textbf{Synchronous Counter}, the global clock is connected in parallel to the clock inputs of \textit{all flip-flops simultaneously}.
\begin{itemize}
    \item State transitions occur in unison within a single flip-flop propagation delay ($t_{cq}$).
    \item \textit{High Frequency:} Maximum clock rate is independent of counter bit-width $n$:
    \[ T_{\text{clock}} \ge t_{cq} + t_{\text{logic\_gate}} \]
\end{itemize}

\begin{table}[htbp]
\centering
\small
\caption{Comparative Analysis: Asynchronous (Ripple) versus Synchronous Counters}
\label{tab:async_vs_sync_counters}
\begin{tabularx}{\textwidth}{l X X}
\toprule
\textbf{Feature} & \\textbf{Asynchronous (Ripple) Counter} & \\textbf{Synchronous Counter} \\
\midrule
\textbf{Clock Distribution} & Applied only to LSB; ripples through stages. & Applied simultaneously in parallel to all flip-flops. \\
\textbf{Operating Speed} & Slow; total latency scales as $O(n \cdot t_{pd})$. & Extremely fast; latency is $O(1)$ stage delay. \\
\textbf{Transient Glitches} & High susceptibility to spurious decoding spikes. & Glitch-free synchronous state transitions. \\
\textbf{Design Complexity} & Minimal; no combinational next-state gates. & Requires AND gating networks for higher-order toggle inputs. \\
\bottomrule
\end{tabularx}
\end{table}

\subsection{Specialized Shift Register Counters}
\begin{enumerate}[label=\bfseries\arabic*.]
    \item \textbf{Ring Counter:} An $n$-stage shift register where output $Q_{n-1}$ is fed back directly to input $D_0$, circulating a single active bit (e.g., $1000 \to 0100 \to 0010 \to 0001$).
    \begin{itemize}
        \item $\text{Modulus} = n$ unique states.
        \item Requires no decoding gates (each flip-flop directly indicates one timing phase).
    \end{itemize}
    \item \textbf{Johnson Counter (Twisted Ring / Moebius Counter):} An $n$-stage shift register where the \textit{inverted} output $\overline{Q}_{n-1}$ is fed back to input $D_0$.
    \begin{itemize}
        \item $\text{Modulus} = 2n$ unique states.
        \item For example, a 4-bit Johnson counter provides a MOD-8 cycle: $0000 \to 1000 \to 1100 \to 1110 \to 1111 \to 0111 \to 0011 \to 0001$.
        \item Yields twice the state capacity of a standard ring counter using identical flip-flop count.
    \end{itemize}
\end{enumerate}


\section{Shift Registers}
\label{sec:shift_registers}

A \textbf{Shift Register} consists of an array of flip-flops connected in cascade, designed to shift binary bits laterally from one stage to adjacent stages upon each clock cycle.

\begin{table}[htbp]
\centering
\small
\caption{Shift Register Architectures and Operational Characteristics}
\label{tab:shift_registers}
\begin{tabularx}{\textwidth}{l l X l}
\toprule
\textbf{Acronym} & \textbf{Full Topology Name} & \textbf{Operational Behavior} & \textbf{Clock Cycles for $n$ Bits} \\
\midrule
\textbf{SISO} & Serial-In Serial-Out & Data enters 1 bit/cycle; leaves 1 bit/cycle. & $n$ to load, $n$ to read. \\
\textbf{SIPO} & Serial-In Parallel-Out & Data enters serially; all $n$ bits read concurrently. & $n$ to load, 1 to read. \\
\textbf{PISO} & Parallel-In Serial-Out & All $n$ bits loaded at once; shifted out serially. & 1 to load, $n$ to read. \\
\textbf{PIPO} & Parallel-In Parallel-Out & Parallel data loaded and read simultaneously. & 1 to load, 1 to read. \\
\bottomrule
\end{tabularx}
\end{table}

\textbf{Universal Shift Register:} Integrates parallel load, serial right shift, serial left shift, and hold modes into a single device using multiplexers on each flip-flop input controlled by 2 mode bits ($S_1, S_0$).


\section{Finite State Machines: Mealy and Moore Models}
\label{sec:fsm}

A \textbf{Finite State Machine (FSM)} is the formal mathematical paradigm used to model, specify, and synthesize complex sequential control units in digital processors.

\subsection{Moore Machine versus Mealy Machine}
\begin{enumerate}[label=\bfseries\arabic*.]
    \item \textbf{Moore Machine Model:} The circuit outputs are a mathematical function \textbf{exclusively of the current internal state}:
    \begin{equation}
    \label{eq:moore_output}
    Y(t) = g(S(t))
    \end{equation}
    \begin{itemize}
        \item Outputs are strictly synchronized with clock edges, isolating downstream circuits from input glitches.
        \item Often requires more internal states to model complex reactive control flows.
    \end{itemize}
    \item \textbf{Mealy Machine Model:} The circuit outputs are a mathematical function of \textbf{both the current internal state and the current external inputs}:
    \begin{equation}
    \label{eq:mealy_output}
    Y(t) = g(S(t), X(t))
    \end{equation}
    \begin{itemize}
        \item Outputs react immediately to asynchronous changes in external inputs within the current clock cycle (faster response time).
        \item Typically synthesizes with fewer internal states.
        \item Vulnerable to input glitches propagating directly through combinational output logic to subsequent stages.
    \end{itemize}
\end{enumerate}

\begin{table}[htbp]
\centering
\small
\caption{Comprehensive Architectural Comparison: Moore versus Mealy Finite State Machines}
\label{tab:moore_vs_mealy}
\begin{tabularx}{\textwidth}{l X X}
\toprule
\textbf{Dimension} & \textbf{Moore Machine} & \textbf{Mealy Machine} \\
\midrule
\textbf{Output Dependence} & Strictly present state: $Y = g(S)$. & Present state and present inputs: $Y = g(S, X)$. \\
\textbf{State Diagram Labeling} & Outputs written inside the state circle ($S / Y$). & Outputs written on the transition arcs ($X / Y$). \\
\textbf{Number of States} & Generally requires equal or greater state count. & Frequently minimizes to fewer states for identical tasks. \\
\textbf{Output Timing} & Changes synchronously on active clock edge. & Changes asynchronously as soon as inputs transition. \\
\textbf{Glitch Immunity} & High; outputs are buffered by state flip-flops. & Low; input glitches pass directly to outputs. \\
\textbf{Hardware Synthesis} & Clean timing closure; ideal for CPU control units. & Fast reactive controllers, protocol handshake engines. \\
\bottomrule
\end{tabularx}
\end{table}



\section{Instruction Set Architecture and Machine Instructions}
\label{sec:isa_machine_instructions}

The \textbf{Instruction Set Architecture (ISA)} defines the programmer-visible interface of the processor. Every high-level programming construct is translated by a compiler into a sequence of binary \textbf{Machine Instructions}.

\subsection{Machine Instruction Formats}
A machine instruction consists of a sequence of bit fields:
\begin{enumerate}[label=\bfseries\arabic*.]
    \item \textbf{Opcode (Operation Code):} Specifies the operation to be performed (e.g., \texttt{ADD}, \texttt{SUB}, \texttt{LOAD}, \texttt{BRANCH}). An $n$-bit opcode can specify up to $2^n$ distinct instructions.
    \item \textbf{Addressing Mode Specifier:} Designates how each operand's physical location (register or memory) is computed.
    \item \textbf{Operand Fields:} Encodes source register indices, destination registers, memory address offsets, or immediate constants.
\end{enumerate}

\subsection{Instruction Architectural Styles: Address Count}
Architectures are classified by the number of explicit operand addresses permitted per instruction:
\begin{enumerate}[label=\bfseries\arabic*.]
    \item \textbf{3-Address Machines (e.g., MIPS, ARM, RISC-V):}
    \begin{center}
    \texttt{ADD R1, R2, R3} \quad $\implies \quad R_1 \leftarrow R_2 + R_3$
    \end{center}
    Allows non-destructive operations; preserves both source operands.
    \item \textbf{2-Address Machines (e.g., Intel x86):}
    \begin{center}
    \texttt{ADD EAX, EBX} \quad $\implies \quad \text{EAX} \leftarrow \text{EAX} + \text{EBX}$
    \end{center}
    Destructive operation; the first operand acts as both a source and the final destination.
    \item \textbf{1-Address Machines (Accumulator Architectures, e.g., Intel 8085):}
    \begin{center}
    \texttt{ADD B} \quad $\implies \quad \text{ACC} \leftarrow \text{ACC} + B$
    \end{center}
    One operand and destination are implicitly the hardware Accumulator (\texttt{ACC}).
    \item \textbf{0-Address Machines (Stack Architectures, e.g., JVM, Forth):}
    \begin{center}
    \texttt{ADD} \quad $\implies \quad \text{TOS} \leftarrow \text{POP}() + \text{POP}()$
    \end{center}
    Operands are implicitly popped from the Top of Stack (\texttt{TOS}), and the result is pushed back.
\end{enumerate}

\subsection{Instruction Categories}
\begin{table}[htbp]
\centering
\small
\caption{Classification of Machine Instruction Categories}
\label{tab:instruction_types}
\begin{tabularx}{\textwidth}{l X X}
\toprule
\textbf{Category} & \textbf{Operational Function} & \textbf{Representative Assembly Mnemonics} \\
\midrule
\textbf{Data Movement} & Transfers bytes between registers and memory. & \texttt{MOV}, \texttt{LDR}, \texttt{STR}, \texttt{PUSH}, \texttt{POP} \\
\textbf{Arithmetic} & Integer and floating-point computations. & \texttt{ADD}, \texttt{SUB}, \texttt{MUL}, \texttt{DIV}, \texttt{INC}, \texttt{DEC} \\
\textbf{Logical} & Bitwise Boolean and manipulation operations. & \texttt{AND}, \texttt{OR}, \texttt{XOR}, \texttt{NOT}, \texttt{TEST} \\
\textbf{Shift / Rotate} & Bit-level lateral displacement. & \texttt{SHL}, \texttt{SHR}, \texttt{SAR}, \texttt{ROL}, \texttt{ROR} \\
\textbf{Control Transfer} & Updates PC to alter program execution flow. & \texttt{JMP}, \texttt{BEQ}, \texttt{BNE}, \texttt{CALL}, \texttt{RET} \\
\textbf{System / Privilege} & Traps into OS kernel; disables interrupts. & \texttt{SYSCALL}, \texttt{SYSRET}, \texttt{CLI}, \texttt{STI}, \texttt{HLT} \\
\bottomrule
\end{tabularx}
\end{table}


\section{Addressing Modes and Effective Address Calculations}
\label{sec:addressing_modes}

An \textbf{Addressing Mode} specifies the mathematical rule for determining the \textbf{Effective Address ($EA$)} of an operand in memory or within the CPU register file.

\subsection{Comprehensive Catalog of Addressing Modes}
\begin{enumerate}[label=\bfseries\arabic*.]
    \item \textbf{Immediate Addressing:} The operand is embedded directly inside the instruction bytes.
    \[ \text{Operand} = \text{Value embedded in instruction} \]
    \textit{Example:} \texttt{MOV R1, \#42}. Requires zero memory data accesses.
    \item \textbf{Register Direct Addressing:} The operand resides inside a CPU general-purpose register.
    \[ \text{Operand} = R[r] \]
    \textit{Example:} \texttt{MOV R1, R2}. High-speed; zero memory accesses.
    \item \textbf{Direct / Absolute Addressing:} The instruction contains the complete physical address of the operand in memory.
    \[ EA = A, \qquad \text{Operand} = \text{Memory}[A] \]
    \textit{Example:} \texttt{LOAD R1, [2000H]}. Requires 1 memory data access.
    \item \textbf{Register Indirect Addressing:} The register holds the memory address of the operand.
    \[ EA = R[r], \qquad \text{Operand} = \text{Memory}[R[r]] \]
    \textit{Example:} \texttt{LOAD R1, [R2]}. Enables pointer manipulation in C.
    \item \textbf{Indexed Addressing:} The effective address is formed by adding an offset/displacement to an Index Register ($IX$):
    \[ EA = A + R[IX], \qquad \text{Operand} = \text{Memory}[EA] \]
    \textit{Utility:} Ideal for traversing linear arrays where $A$ points to the array base and $IX$ holds the loop index.
    \item \textbf{Base-Register / Displacement Addressing:} The effective address is computed by adding an offset to a Base Register ($BR$):
    \[ EA = R[BR] + \text{Displacement} \]
    \textit{Utility:} Essential for position-independent code and relocatable code segments.
    \item \textbf{PC-Relative Addressing:} The displacement is added to the Program Counter:
    \[ EA = \text{PC} + \text{Displacement} \]
    \textit{Utility:} Used in conditional branch instructions (\texttt{BEQ}, \texttt{BNE}) to jump within local code blocks without knowing absolute addresses.
    \item \textbf{Implied / Inherent Addressing:} The operands are implicitly defined by the opcode itself.
    \textit{Example:} \texttt{CLC} (Clear Carry Flag), \texttt{RET} (Return from Subroutine).
    \item \textbf{Stack Addressing:} Operands are implicitly referenced from the memory location pointed to by the Stack Pointer (\texttt{SP}).
    \textit{Example:} \texttt{PUSH R1}, \texttt{POP R2}.
    \item \textbf{Auto-Increment and Auto-Decrement:} The register is used as a pointer ($EA = R[r]$) and is automatically incremented or decremented by the operand size (e.g., by 4 bytes for 32-bit words).
\end{enumerate}

\begin{table}[htbp]
\centering
\small
\caption{Comprehensive Summary of Processor Addressing Modes}
\label{tab:addressing_modes}
\begin{tabularx}{\textwidth}{l l l c}
\toprule
\textbf{Addressing Mode} & \textbf{Effective Address ($EA$)} & \textbf{Assembly Syntax} & \textbf{Data Memory Accesses} \\
\midrule
\textbf{Immediate} & None (Embedded operand) & \texttt{MOV R1, \#100} & 0 \\
\textbf{Register Direct} & None (Register holds operand) & \texttt{ADD R1, R2} & 0 \\
\textbf{Direct (Absolute)} & $EA = A$ & \texttt{LOAD R1, [1050H]} & 1 \\
\textbf{Register Indirect} & $EA = R[r]$ & \texttt{LOAD R1, [R2]} & 1 \\
\textbf{Memory Indirect} & $EA = \text{Memory}[A]$ & \texttt{LOAD R1, @[1050H]} & 2 \\
\textbf{Indexed} & $EA = R[IX] + A$ & \texttt{LOAD R1, [R2 + 40]} & 1 \\
\textbf{Base-Displacement} & $EA = R[BR] + \text{Offset}$ & \texttt{MOV R1, [RBP - 8]} & 1 \\
\textbf{PC-Relative} & $EA = \text{PC} + \text{Offset}$ & \texttt{BNE +16} & 0 (Branch target) \\
\textbf{Implied} & Defined implicitly & \texttt{NOP}, \texttt{STC} & 0 \\
\bottomrule
\end{tabularx}
\end{table}


\section{Stacks, Subroutines, and Control Flow Management}
\label{sec:stack_subroutines}

\subsection{Call Stack Architecture}
The \textbf{Call Stack} is a Last-In, First-Out (LIFO) contiguous region of memory utilized by the operating system and compilers to manage nested procedure calls, pass function parameters, allocate local automatic variables, and preserve register state.

\begin{itemize}
    \item \textbf{Stack Growth Direction:} In modern x86 and ARM platforms, the stack grows \textbf{downward} from high memory addresses toward lower memory addresses.
    \item \textbf{Stack Pointer (\texttt{SP}):} Always points to the current top of the stack (the lowest address currently allocated).
    \begin{align*}
    \text{\texttt{PUSH R1}} &\implies \text{SP} \leftarrow \text{SP} - 4, \quad \text{Memory}[\text{SP}] \leftarrow R_1 \\
    \text{\texttt{POP R1}}  &\implies R_1 \leftarrow \text{Memory}[\text{SP}], \quad \text{SP} \leftarrow \text{SP} + 4
    \end{align*}
    \item \textbf{Frame Pointer (\texttt{FP} / \texttt{EBP}):} Fixed base register pointing to the base of the current function's activation record (stack frame), allowing stable offsets to access parameters ($[\text{FP} + 8]$) and local variables ($[\text{FP} - 4]$).
\end{itemize}

\subsection{Subroutine Call (\texttt{CALL}) and Return (\texttt{RET}) Protocol}
\begin{enumerate}[label=\bfseries\arabic*.]
    \item \textbf{\texttt{CALL Target}} Execution:
    \begin{align*}
    \text{SP} &\leftarrow \text{SP} - 4 \\
    \text{Memory}[\text{SP}] &\leftarrow \text{PC}_{\text{next}} \quad \text{(Saves Return Address)} \\
    \text{PC} &\leftarrow \text{Target Address}
    \end{align*}
    \item \textbf{\texttt{RET}} Execution:
    \begin{align*}
    \text{PC} &\leftarrow \text{Memory}[\text{SP}] \quad \text{(Restores Return Address)} \\
    \text{SP} &\leftarrow \text{SP} + 4
    \end{align*}
\end{enumerate}

\begin{table}[htbp]
\centering
\small
\caption{Comparative Analysis: Hardware Interrupt versus Software Subroutine Call}
\label{tab:interrupt_vs_call}
\begin{tabularx}{\textwidth}{l X X}
\toprule
\textbf{Dimension} & \textbf{Hardware Interrupt} & \textbf{Subroutine Call (\texttt{CALL})} \\
\midrule
\textbf{Trigger Origin} & Asynchronous external hardware pin (e.g., timer, NIC, disk). & Synchronous software instruction explicitly executed in program code. \\
\textbf{Predictability} & Can occur at any random clock cycle between any two instructions. & Occurs at a strictly deterministic, programmed location in the code. \\
\textbf{Saved State} & CPU automatically pushes PC \textit{and} Program Status Word (PSW/Flags). & Automatically pushes only the Return Address (PC); flags not preserved. \\
\textbf{Target Address} & Fetched from Interrupt Vector Table (IVT) indexed by interrupt type. & Target address encoded directly in the instruction or in a register. \\
\textbf{Return Directive} & Uses \texttt{IRET} (Interrupt Return), restoring PC and Flags simultaneously. & Uses standard \texttt{RET}, restoring only PC. \\
\bottomrule
\end{tabularx}
\end{table}


\section{Input/Output Organization and Data Transfer Paradigms}
\label{sec:io_organization}

Peripheral I/O devices (keyboards, solid-state storage, network controllers) exhibit speeds and electrical characteristics vastly different from the high-speed CPU core. Three primary architectural data transfer mechanisms govern I/O interaction:

\begin{enumerate}[label=\bfseries\arabic*.]
    \item \textbf{Programmed I/O (Polling / Busy Waiting):}
    \begin{itemize}
        \item The CPU continuously interrogates a peripheral status register in a tight software loop until the device signals ready ($100\%$ CPU utilization wasted on polling).
        \item Simple hardware, but completely halts computational progress during slow I/O operations.
    \end{itemize}
    \item \textbf{Interrupt-Driven I/O:}
    \begin{itemize}
        \item The CPU issues an I/O command and immediately proceeds to execute other user processes.
        \item When the device completes the transfer, it asserts an asynchronous hardware interrupt request line (\texttt{INTR}).
        \item The CPU suspends the active thread, saves state, executes the device's \textbf{Interrupt Service Routine (ISR)}, and resumes the interrupted program.
    \end{itemize}
    \item \textbf{Direct Memory Access (DMA):}
    \begin{itemize}
        \item Designed for high-speed bulk data transfers (e.g., disk sectors, gigabit network packets).
        \item A specialized hardware coprocessor (the \textbf{DMA Controller}, DMAC) takes control of the system buses from the CPU.
        \item Data transfers occur directly between the peripheral controller and main memory at bus wire speed without routing data bytes through CPU registers.
    \end{itemize}
\end{enumerate}

\subsection{DMA Operating Modes}
\begin{itemize}
    \item \textbf{Burst Mode (Block Transfer):} The DMAC seizes the system bus and transfers an entire block of data (e.g., 4 KB) continuously before releasing the bus back to the CPU. High throughput, but locks out the CPU during the burst.
    \item \textbf{Cycle Stealing Mode:} The DMAC takes control of the system bus for exactly one bus cycle per byte/word, interleaving memory cycles between CPU instruction phases. Lowers peak transfer rate but prevents CPU starvation.
    \item \textbf{Transparent Mode (Hidden DMA):} The DMAC transfers data only during clock phases when the CPU is decoding or executing internally and not using the external memory bus. Zero impact on CPU execution time.
\end{itemize}

\subsection{Memory-Mapped I/O versus Isolated (Port-Mapped) I/O}
\begin{table}[htbp]
\centering
\small
\caption{Architectural Comparison: Memory-Mapped I/O versus Isolated (Port) I/O}
\label{tab:memory_vs_isolated_io}
\begin{tabularx}{\textwidth}{l X X}
\toprule
\textbf{Dimension} & \textbf{Memory-Mapped I/O} & \textbf{Isolated (Port-Mapped) I/O} \\
\midrule
\textbf{Address Space} & I/O device registers share the same unified memory address space. & Completely separate, isolated address space for I/O ports. \\
\textbf{Instruction Set} & Standard memory instructions (\texttt{MOV}, \texttt{LOAD}, \texttt{STORE}) access I/O. & Requires dedicated I/O machine instructions (\texttt{IN}, \texttt{OUT}). \\
\textbf{Control Signals} & Uses standard memory control lines (\texttt{MEMR}, \texttt{MEMW}). & Requires distinct I/O control lines (\texttt{IOR}, \texttt{IOW}). \\
\textbf{Memory Capacity} & Reduces addressable RAM capacity by dedicating a block of addresses to I/O. & Full address space available for RAM; no capacity sacrifice. \\
\textbf{Addressing Modes} & All processor addressing modes are applicable to I/O operations. & Severely restricted addressing modes (usually register or port direct). \\
\textbf{Architectures} & ARM, RISC-V, MIPS, Motorola 68000. & Intel x86 (supports both, but provides 64K port space). \\
\bottomrule
\end{tabularx}
\end{table}


\section{Pipelining, Datapath Hazards, and Performance Optimization}
\label{sec:pipelining}

\textbf{Pipelining} is an architectural implementation technique where multiple instructions are overlapped in execution, analogous to an industrial automotive assembly line. Pipelining increases instruction \textbf{throughput} (instructions completed per second) without decreasing individual instruction \textbf{latency}.

\subsection{The Classical 5-Stage RISC Pipeline}
\begin{enumerate}[label=\bfseries\arabic*.]
    \item \textbf{IF (Instruction Fetch):} Fetch instruction from instruction cache into the IR; increment $\text{PC} \leftarrow \text{PC} + 4$.
    \item \textbf{ID (Instruction Decode \& Register Fetch):} Decode opcode; read source register operands $R_s, R_t$ from the register file; sign-extend immediate fields.
    \item \textbf{EX (Execution \& Address Calculation):} Perform ALU computation or compute effective memory address for load/store instructions.
    \item \textbf{MEM (Memory Access):} Read data from or write data to the data cache for load/store instructions; bypass stage for register-register operations.
    \item \textbf{WB (Write-Back):} Write computational result or loaded data into the destination register $R_d$ in the register file.
\end{enumerate}

\subsection{Theoretical Pipeline Speedup Formula}
Let $n$ be the number of instructions executed, and $k$ be the number of pipeline stages.
\begin{itemize}
    \item \textbf{Non-Pipelined Execution Time:}
    \begin{equation}
    \label{eq:non_pipelined_time}
    T_{\text{non-pipe}} = n \cdot k \cdot \tau
    \end{equation}
    where $\tau$ is the clock cycle duration.
    \item \textbf{Pipelined Execution Time:}
    The first instruction requires $k$ cycles to fill the pipe. Each subsequent instruction completes at the rate of 1 instruction per cycle:
    \begin{equation}
    \label{eq:pipelined_time}
    T_{\text{pipe}} = [k + (n - 1)] \cdot \tau
    \end{equation}
    \item \textbf{Theoretical Pipeline Speedup ($S$):}
    \begin{equation}
    \label{eq:pipeline_speedup}
    S = \frac{T_{\text{non-pipe}}}{T_{\text{pipe}}} = \frac{n \cdot k \cdot \tau}{[k + (n - 1)] \cdot \tau} = \frac{n \cdot k}{k + n - 1}
    \end{equation}
    As $n \to \infty$ (large program size):
    \begin{equation}
    \label{eq:pipeline_speedup_limit}
    \lim_{n \to \infty} S = k
    \end{equation}
    \textbf{Amdahl's Pipeline Law:} The maximum theoretical speedup achievable by a pipeline is strictly bounded by the number of pipeline stages $k$.
\end{itemize}

\subsection{Pipeline Hazards}
A \textbf{Hazard} is any condition that prevents the pipeline from executing the next instruction in its designated clock cycle, mandating the insertion of empty idle cycles (\textbf{Pipeline Stalls} or \textbf{Bubbles}). Hazards fall into three distinct classes:

\subsubsection{A. Structural Hazards (Resource Conflicts)}
Occur when two or more instructions in different pipeline stages attempt to access the identical physical hardware resource simultaneously.
\begin{itemize}
    \item \textit{Classic Example:} A single-ported unified memory accessed by the IF stage (fetching an instruction) and the MEM stage (loading an operand) during the same clock cycle.
    \item \textit{Mitigation:} Modified Harvard architecture with independent L1 Instruction and L1 Data caches.
\end{itemize}

\subsubsection{B. Data Hazards (Data Dependencies)}
Occur when an instruction depends on the result of an earlier instruction that is still in flight in the pipeline and has not yet written its result to the register file.

\begin{table}[htbp]
\centering
\small
\caption{Classification of Data Dependency Hazards}
\label{tab:data_hazards}
\begin{tabularx}{\textwidth}{l l X}
\toprule
\textbf{Hazard Type} & \textbf{Formal Dependency Name} & \textbf{Architectural Description and Phenomenon} \\
\midrule
\textbf{RAW} & Read-After-Write (True Dependency) & Instruction $J$ attempts to read operand before Instruction $I$ writes it. (Most common hazard in in-order pipelines). \\
\textbf{WAR} & Write-After-Read (Anti-Dependency) & Instruction $J$ attempts to write destination before Instruction $I$ reads it. (Occurs in out-of-order execution engines). \\
\textbf{WAW} & Write-After-Write (Output Dependency) & Instruction $J$ attempts to write destination before Instruction $I$ writes it, destroying program order. (Out-of-order engines). \\
\bottomrule
\end{tabularx}
\end{table}

\textbf{Mitigations for Data Hazards:}
\begin{enumerate}[label=\bfseries\arabic*.]
    \item \textbf{Hardware Data Forwarding (Bypassing):} Routes the computational result directly from the ALU output of the EX or MEM stage back to the ALU input of the EX stage via dedicated bypass multiplexers, eliminating stalls for ALU-to-ALU dependencies.
    \item \textbf{Load-Use Hazard Stall:} Even with forwarding, a \texttt{LOAD} instruction produces its data at the end of the MEM stage. If the immediately following instruction requires that data in its EX stage, \textbf{one bubble cycle is mandatory}. Compilers perform instruction scheduling to fill load-delay slots.
\end{enumerate}

\subsubsection{C. Control Hazards (Branch Hazards)}
Occur when the pipeline fetches instructions following a conditional branch (\texttt{BEQ}, \texttt{BNE}) before the branch outcome (taken vs. not taken) and target address have been computed in the EX or MEM stage.
\begin{itemize}
    \item \textbf{Branch Penalty:} Flushes wrong-path instructions from the pipe, discarding 2--3 cycles of work.
    \item \textbf{Mitigations:}
    \begin{itemize}
        \item \textbf{Static Branch Prediction:} Predict always not-taken, or predict backward branches taken (loops) and forward branches not-taken (error checks).
        \item \textbf{Dynamic Branch Prediction:} Uses a \textbf{Branch History Table (BHT)} with 2-bit saturating counters (Strongly Taken, Weakly Taken, Weakly Not-Taken, Strongly Not-Taken), achieving $>95\%$ prediction accuracy.
        \item \textbf{Branch Target Buffer (BTB):} Caches target addresses to enable zero-cycle branch redirections.
        \item \textbf{Delayed Branching:} Compiler places a useful instruction inside the branch delay slot that executes unconditionally.
    \end{itemize}
\end{itemize}



\section{RISC versus CISC Architectural Philosophies}
\label{sec:risc_vs_cisc}

Instruction Set Architectures historically evolved along two opposing philosophical trajectories: \textbf{Reduced Instruction Set Computer (RISC)} and \textbf{Complex Instruction Set Computer (CISC)}.

\subsection{Architectural Core Principles}
\begin{enumerate}[label=\bfseries\arabic*.]
    \item \textbf{CISC (Complex Instruction Set Computer):}
    \begin{itemize}
        \item \textit{Historical Genesis:} Emerged in the 1970s when magnetic core memory was slow and expensive, and compilers were primitive. The goal was to provide rich, expressive, multi-cycle machine instructions that closely mirrored high-level language constructs to minimize program memory footprint.
        \item \textit{Key Attributes:} Variable-length instructions (e.g., x86 instructions range from 1 to 15 bytes); memory-to-memory operations (\texttt{ADD [EBX], EAX}); specialized addressing modes; microprogrammed ROM control units.
    \end{itemize}
    \item \textbf{RISC (Reduced Instruction Set Computer):}
    \begin{itemize}
        \item \textit{Empirical Discovery (Patterson \& Hennessy):} In the 1980s, compiler traces revealed that $>80\%$ of typical execution time was spent executing simple instructions (load, store, branch, integer add), while complex instructions were rarely used by optimizing compilers.
        \item \textit{Key Attributes:} Fixed-length instructions (typically 32 bits); strict \textbf{Load/Store Architecture} (computational instructions operate exclusively on registers; only \texttt{LOAD} and \texttt{STORE} touch memory); single-cycle execution target ($\text{CPI} \approx 1$); large general-purpose register file (32+ registers); hardwired control units.
    \end{itemize}
\end{enumerate}

\begin{table}[htbp]
\centering
\small
\caption{Comprehensive Comparison: RISC versus CISC Architectures}
\label{tab:risc_vs_cisc}
\begin{tabularx}{\textwidth}{l X X}
\toprule
\textbf{Architectural Dimension} & \textbf{RISC Architecture} & \textbf{CISC Architecture} \\
\midrule
\textbf{Instruction Length} & Fixed width (typically 32 bits uniform). & Highly variable (1 to 15 bytes in x86-64). \\
\textbf{Instruction Count per Task} & Higher (more simple instructions needed). & Lower (single instruction encodes complex loop/copy). \\
\textbf{Cycles Per Instruction (CPI)} & Targets $\approx 1.0$ cycle per instruction. & Multi-cycle; ranges from 1 to over 50 cycles. \\
\textbf{Memory Operands} & Load/Store only; ALU operates on registers. & Direct memory-to-register and memory-to-memory. \\
\textbf{Addressing Modes} & Few, simple modes (Immediate, Register, Base+Disp). & Many complex modes (Base+Index+Scale+Disp). \\
\textbf{Register File} & Large general-purpose register set (32--128 registers). & Historically smaller register set (8 in x86; 16 in x64). \\
\textbf{Control Unit Implementation} & Hardwired logic (combinational gates, ultra-fast). & Microprogrammed control store (microcode ROM). \\
\textbf{Pipelining Feasibility} & Extremely straightforward; uniform stage timings. & Complex due to variable instruction decode lengths. \\
\textbf{Compiler Burden} & Heavy; optimizer schedules instructions and pipelines. & Moderate; hardware handles complex operations. \\
\textbf{Representative Systems} & ARM (Cortex-A/M), RISC-V, MIPS, Apple Silicon. & Intel x86, AMD64, IBM z/Architecture, VAX. \\
\bottomrule
\end{tabularx}
\end{table}

\subsection{Modern Architectural Convergence}
The historic binary distinction between RISC and CISC has largely dissolved in modern microarchitectures:
\begin{itemize}
    \item Modern CISC processors (Intel Core, AMD Ryzen) maintain a CISC x86-64 outer ISA for binary backward compatibility, but their internal frontend decoders translate CISC instructions on the fly into fixed-length, RISC-like primitive micro-operations (\textbf{$\mu$-ops}).
    \item These internal $\mu$-ops are subsequently scheduled and executed on a high-throughput, out-of-order, pipelined superscalar execution engine that is purely RISC in character.
\end{itemize}


\section{Microprocessors versus Microcontrollers}
\label{sec:mpu_vs_mcu}

\subsection{Architectural Distinctions}
\begin{enumerate}[label=\bfseries\arabic*.]
    \item \textbf{Microprocessor (MPU):} An integrated circuit that contains only the Central Processing Unit (ALU, Control Unit, and registers) on-chip. All primary storage (DRAM), persistent storage (Flash, SSD), and peripheral controllers must be connected externally across printed circuit board (PCB) buses.
    \begin{itemize}
        \item \textit{Primary Focus:} High-performance, high-clock-rate general-purpose computing executing full operating systems (Linux, Windows, macOS).
    \end{itemize}
    \item \textbf{Microcontroller (MCU):} A complete \textbf{System-on-Chip (SoC)} integrating a CPU core along with RAM, non-volatile program memory (Flash/EEPROM), clock oscillators, timers, and diverse peripheral controllers (GPIO, ADC, DAC, UART, SPI, I2C, CAN, PWM) onto a single silicon die.
    \begin{itemize}
        \item \textit{Primary Focus:} Low-cost, low-power, dedicated real-time control applications (automotive ECUs, medical implants, industrial robotics, IoT nodes).
    \end{itemize}
\end{enumerate}

\begin{table}[htbp]
\centering
\small
\caption{Comparative Analysis: Microprocessor (MPU) versus Microcontroller (MCU)}
\label{tab:mpu_vs_mcu}
\begin{tabularx}{\textwidth}{l X X}
\toprule
\textbf{Feature} & \textbf{Microprocessor (MPU)} & \textbf{Microcontroller (MCU)} \\
\midrule
\textbf{Integration Level} & CPU core only; external RAM, ROM, I/O. & Complete monolithic SoC (CPU + RAM + Flash + Peripherals). \\
\textbf{Clock Speed} & Hundreds of MHz to $>5.0\text{ GHz}$. & Typically $1\text{ MHz}$ to $500\text{ MHz}$. \\
\textbf{Power Consumption} & High ($10\text{ W}$ to $>250\text{ W}$); requires active cooling. & Low ($<1\text{ mW}$ to $5\text{ W}$); optimized for battery operation. \\
\textbf{RAM Capacity} & Gigabytes to Terabytes of external DRAM. & Kilobytes to Megabytes of internal SRAM. \\
\textbf{Operating Systems} & Full multitasking OS (Linux, Windows, BSD). & Bare-metal firmware or Real-Time OS (FreeRTOS, Zephyr). \\
\textbf{Hardware Interrupts} & Routed via external APIC / PIC controllers. & Built-in Nested Vectored Interrupt Controller (NVIC). \\
\textbf{System Cost} & High system BOM (requires multi-layer PCB, PMIC). & Extremely low (\$0.50 to \$10.00 for single-chip design). \\
\textbf{Examples} & Intel Core i9, AMD EPYC, Apple M-series. & STM32 (ARM Cortex-M), Microchip PIC, ATmega328P. \\
\bottomrule
\end{tabularx}
\end{table}


\section{Intel 8085 versus 8086 Microprocessors and Memory Segmentation}
\label{sec:8085_vs_8086}

The Intel 8085 and 8086 represent the watershed historical architectures foundational to all modern x86 computing systems.

\subsection{The Intel 8085 Microprocessor}
Introduced in 1976, the 8085 is an 8-bit accumulator-based NMOS processor operating at $3.072\text{ MHz}$ on a single $+5\text{ V}$ power supply.
\begin{itemize}
    \item \textbf{Bus Architecture:}
    \begin{itemize}
        \item 16-bit address bus addressing $2^{16} = 64\text{ KiB}$ of memory.
        \item Multiplexed lower address/data bus ($\text{AD}_0$--$\text{AD}_7$). The \textbf{Address Latch Enable (ALE)} signal demultiplexes this bus via an external 74LS373 octal latch.
    \end{itemize}
    \item \textbf{Register Set:} Accumulator A (8-bit), temporary registers B, C, D, E, H, L (usable as 16-bit register pairs: BC, DE, HL), Stack Pointer (SP, 16-bit), Program Counter (PC, 16-bit).
    \item \textbf{Status Flags:} Sign (S), Zero (Z), Auxiliary Carry (AC, used for BCD), Parity (P), Carry (CY).
    \item \textbf{Hardware Interrupt Priority:}
    \[ \textbf{TRAP} \text{ (NMI, Non-Maskable, Highest)} > \textbf{RST 7.5} > \textbf{RST 6.5} > \textbf{RST 5.5} > \textbf{INTR} \text{ (Maskable, Lowest)} \]
\end{itemize}

\subsection{The Intel 8086 Microprocessor}
Introduced in 1978, the 8086 is a true 16-bit processor with a 20-bit address bus, establishing the foundation of the x86 family.
\begin{itemize}
    \item \textbf{Bus Architecture:}
    \begin{itemize}
        \item 20-bit address bus addressing $2^{20} = 1\text{ MiB}$ of physical memory.
        \item 16-bit true data bus ($\text{AD}_0$--$\text{AD}_{15}$).
    \end{itemize}
    \item \textbf{Pipelined Internal Architecture:} Partitioned into two independent processing units:
    \begin{enumerate}[label=\bfseries\arabic*.]
        \item \textbf{Bus Interface Unit (BIU):} Fetches instruction bytes from memory ahead of time and stores them in a \textbf{6-byte Instruction Prefetch Queue}. Manages segment registers and physical address generation.
        \item \textbf{Execution Unit (EU):} Decodes instructions and executes them via the 16-bit ALU and register file.
    \end{enumerate}
    This parallel overlap represents primitive 2-stage instruction pipelining.
\end{itemize}

\subsection{8086 Segmented Memory Architecture and Physical Address Generation}
Because internal registers are 16 bits wide, an offset register can generate addresses only up to $2^{16} = 64\text{ KiB}$. To access a 20-bit physical address space ($1\text{ MiB}$), the 8086 divides memory into logical \textbf{Segments} of up to $64\text{ KiB}$ each.

\textbf{Segment Registers (16-bit):}
\begin{itemize}
    \item \textbf{CS (Code Segment):} Paired with \textbf{IP (Instruction Pointer)} for instruction fetching.
    \item \textbf{DS (Data Segment):} Paired with general registers, \textbf{SI (Source Index)}, or \textbf{DI (Destination Index)} for data variables.
    \item \textbf{SS (Stack Segment):} Paired with \textbf{SP (Stack Pointer)} or \textbf{BP (Base Pointer)} for stack operations.
    \item \textbf{ES (Extra Segment):} Paired with \textbf{DI} for string destination operations.
\end{itemize}

\textbf{Mathematical Formula for Physical Address Calculation:}
To calculate the 20-bit physical address, the 16-bit Segment Base is shifted left by 4 bits (multiplied by $16_{10}$ or $10_{16}$) and added to the 16-bit Offset:
\begin{equation}
\label{eq:physical_address_8086}
\text{Physical Address} = (\text{Segment} \times 16_{10}) + \text{Offset} = (\text{Segment} \ll 4) + \text{Offset}
\end{equation}

\textbf{Worked Example:} Compute the physical memory address if $\text{CS} = 348A_{16}$ and $\text{IP} = 4214_{16}$.
\begin{align*}
\text{Segment Base} &= 348A_{16} \times 10_{16} = 348A0_{16} \\
\text{Offset (IP)}  &= \phantom{0000}4214_{16} \\
\text{Physical Address} &= 348A0_{16} + 4214_{16} = \mathbf{38AB4_{16}}
\end{align*}

\begin{table}[htbp]
\centering
\small
\caption{Comparative Analysis: Intel 8085 versus Intel 8086}
\label{tab:8085_vs_8086}
\begin{tabularx}{\textwidth}{l X X}
\toprule
\textbf{Specification} & \textbf{Intel 8085} & \textbf{Intel 8086} \\
\midrule
\textbf{Word Size} & 8-bit processor. & 16-bit processor. \\
\textbf{Data Bus Width} & 8 bits (multiplexed $\text{AD}_0$--$\text{AD}_7$). & 16 bits (multiplexed $\text{AD}_0$--$\text{AD}_{15}$). \\
\textbf{Address Bus Width} & 16 bits. & 20 bits. \\
\textbf{Addressable Memory} & $2^{16} = 64\text{ KiB}$. & $2^{20} = 1\text{ MiB}$. \\
\textbf{Memory Organization} & Single flat linear address space. & Segmented memory space (4 active segments of 64 KB). \\
\textbf{Pipelining / Queue} & None. No instruction prefetch queue. & 6-byte instruction prefetch queue in BIU. \\
\textbf{Internal Units} & Unified processing core. & Partitioned into BIU (Bus Interface) and EU (Execution). \\
\textbf{Operating Modes} & Single standard mode. & Minimum Mode (single CPU) and Maximum Mode (multiprocessor). \\
\textbf{Arithmetic Support} & Only 8-bit integer add/sub. & Hardware 16-bit multiply and divide instructions (\texttt{MUL}, \texttt{DIV}). \\
\bottomrule
\end{tabularx}
\end{table}


\section{32-Bit versus 64-Bit Processor Architectures}
\label{sec:32bit_vs_64bit}

The designation ``$N$-bit processor'' designates the native register bit-width and data bus width of the integer execution datapath.

\subsection{Addressable Memory Space Scaling}
\begin{itemize}
    \item \textbf{32-Bit Systems (e.g., x86 IA-32, ARMv7):}
    A 32-bit register can generate $2^{32}$ unique binary memory addresses:
    \[ 2^{32} \text{ bytes} = 4,294,967,296\text{ bytes} = \mathbf{4\text{ GiB}} \]
    This 4 GiB ceiling severely constrains enterprise database servers, high-performance computing, and virtualization.
    \item \textbf{64-Bit Systems (e.g., x86-64 / AMD64, ARM64 / AArch64):}
    A 64-bit register provides a theoretical addressable space of $2^{64}$ bytes:
    \[ 2^{64} \text{ bytes} = 18,446,744,073,709,551,616\text{ bytes} = \mathbf{16\text{ Exabytes (EiB)}} \]
    Current commercial 64-bit microprocessors implement 48-bit or 57-bit virtual addressing ($256\text{ TiB}$ to $128\text{ PiB}$) and up to 52-bit physical addressing to economize on silicon address translation page tables.
\end{itemize}

\subsection{Architectural Trade-offs: 32-Bit versus 64-Bit}
\begin{table}[htbp]
\centering
\small
\caption{Architectural Comparison: 32-Bit versus 64-Bit Computing Platforms}
\label{tab:32bit_vs_64bit}
\begin{tabularx}{\textwidth}{l X X}
\toprule
\textbf{Dimension} & \textbf{32-Bit Platform} & \textbf{64-Bit Platform} \\
\midrule
\textbf{GPR Bit Width} & 32 bits (e.g., EAX, EBX, ECX). & 64 bits (e.g., RAX, RBX, RCX). \\
\textbf{Native Integer Range} & $[-2^{31}, 2^{31}-1]$ ($\approx \pm 2.14\times 10^9$). & $[-2^{63}, 2^{63}-1]$ ($\approx \pm 9.22\times 10^{18}$). \\
\textbf{Pointer Storage} & 4 bytes per pointer in memory. & 8 bytes per pointer in memory (\textit{Pointer Bloat}). \\
\textbf{Register File Count} & x86 provides 8 general registers. & x86-64 doubles the register count to 16 registers (RAX--R15). \\
\textbf{Calling Convention} & Parameters passed on the slow stack. & Fast ABI: first 6 parameters passed in CPU registers. \\
\textbf{Memory Footprint} & Smaller executable binary sizes. & 10--30\% larger binary footprint due to 8-byte pointers. \\
\textbf{Mathematical Throughput} & 64-bit math requires multi-instruction emulation. & 64-bit integers and SIMD vectors computed in single cycles. \\
\bottomrule
\end{tabularx}
\end{table}



\section{Memory Hierarchy and Cache Memory Fundamentals}
\label{sec:memory_hierarchy}

Microprocessor computation speed historically outpaced DRAM memory access latency by several orders of magnitude (the \textbf{Memory Wall}). To bridge this bandwidth-latency deficit, modern architectures organize storage into a stratified \textbf{Memory Hierarchy}.

\subsection{The Principle of Locality}
The memory hierarchy exploits the empirical reality that computer programs do not access their memory space uniformly. Memory accesses exhibit two fundamental dimensions of locality:
\begin{enumerate}[label=\bfseries\arabic*.]
    \item \textbf{Temporal Locality (Locality in Time):} If a particular memory location is referenced at time $t$, it is highly probable that the identical location will be re-referenced again in the near future ($t + \Delta t$).
    \begin{itemize}
        \item \textit{Manifestations:} Loop index variables, accumulator registers, instruction loop bodies, stack frame pointers.
    \end{itemize}
    \item \textbf{Spatial Locality (Locality in Space):} If a memory location at address $A$ is referenced, it is highly probable that memory locations physically proximate to $A$ ($A \pm 1, A \pm 2, \dots$) will be accessed soon.
    \begin{itemize}
        \item \textit{Manifestations:} Sequential instruction execution paths, linear array traversals, contiguous struct record member accesses.
    \end{itemize}
    \item \textbf{Sequential Locality:} A specialized subset of spatial locality wherein access patterns proceed strictly sequentially through monotonically ascending addresses (e.g., streaming audio/video, instruction prefetch pipelines).
\end{enumerate}

\begin{figure}[htbp]
\centering
\small
\begin{tcolorbox}[colback=blue!5!white,colframe=blue!75!black,title=\textbf{The Classical Computer Memory Hierarchy Pyramid}]
\centering
\textbf{CPU Registers} (Bit width: 64 bits $\cdot$ Latency: $<0.5\text{ ns}$ $\cdot$ Capacity: $\sim 1\text{ KB}$ $\cdot$ Cost: Highest) \\
$\Downarrow$ \\
\textbf{Level-1 Cache (L1-I / L1-D)} (Latency: $\sim 1\text{ ns}$ $\cdot$ Capacity: $32$--$64\text{ KB}$/core $\cdot$ SRAM) \\
$\Downarrow$ \\
\textbf{Level-2 Cache (L2)} (Latency: $\sim 3$--$5\text{ ns}$ $\cdot$ Capacity: $512\text{ KB}$--$2\text{ MB}$/core $\cdot$ SRAM) \\
$\Downarrow$ \\
\textbf{Level-3 Cache (L3)} (Latency: $\sim 10$--$20\text{ ns}$ $\cdot$ Capacity: $16$--$128\text{ MB}$ shared $\cdot$ SRAM) \\
$\Downarrow$ \\
\textbf{Main Memory (DRAM / DDR5)} (Latency: $\sim 50$--$80\text{ ns}$ $\cdot$ Capacity: $16$--$128\text{ GB}$ $\cdot$ Cost: Moderate) \\
$\Downarrow$ \\
\textbf{Secondary Storage (NVMe SSD / PCIe)} (Latency: $\sim 10$--$50\ \mu\text{s}$ $\cdot$ Capacity: $512\text{ GB}$--$4\text{ TB}$ $\cdot$ Flash NAND) \\
$\Downarrow$ \\
\textbf{Auxiliary Storage (Magnetic HDD / Cloud / Tape)} (Latency: $\sim 5$--$15\text{ ms}$ $\cdot$ Capacity: $10$--$100\text{ TB}$ $\cdot$ Cost: Lowest)
\end{tcolorbox}
\caption{Hierarchical Stratification: Balancing Access Latency, Capacity, and Cost}
\label{fig:memory_pyramid}
\end{figure}

\subsection{Cache Hit, Cache Miss, and Hit Ratio}
When the CPU core requests a word from memory:
\begin{itemize}
    \item \textbf{Cache Hit:} The requested word is currently present in the cache. The CPU reads or writes the data immediately at the cache's native clock speed (\textbf{Hit Time}).
    \item \textbf{Cache Miss:} The requested word is absent from the cache. The request is forwarded down the hierarchy to the next lower tier (L2, L3, or DRAM). The pipeline stalls while an entire fixed-size block of data containing the word (\textbf{Cache Line}) is fetched (\textbf{Miss Penalty}).
\end{itemize}

\begin{equation}
\label{eq:hit_ratio}
\text{Hit Ratio } (h) = \frac{\text{Number of Cache Hits}}{\text{Total Memory References}}
\end{equation}
\begin{equation}
\label{eq:miss_ratio}
\text{Miss Ratio } (m) = 1 - h = \frac{\text{Number of Cache Misses}}{\text{Total Memory References}}
\end{equation}

\subsection{The 3 Cs Model of Cache Misses}
Mark Hill categorized cache misses into three fundamental classes (the \textbf{3 Cs}):
\begin{enumerate}[label=\bfseries\arabic*.]
    \item \textbf{Compulsory Misses (Cold Misses):} The very first access to a cache block. Because the block has never been brought into the cache, the miss is unavoidable. (Reduced by larger cache line block sizes and hardware prefetching).
    \item \textbf{Capacity Misses:} Occur when the total working set of active program data exceeds the total physical capacity of the cache. Blocks are evicted and must be re-fetched later. (Reduced by increasing total cache size).
    \item \textbf{Conflict Misses (Collision Misses):} Occur in direct-mapped or set-associative caches when multiple active memory blocks map to the exact same cache set, causing thrashing and premature evictions even when the overall cache has abundant empty space. (Reduced by increasing associativity).
    \item \textbf{Coherence Misses (The 4th C in Multiprocessors):} Occur when a cache line is invalidated because another CPU core wrote to that address across the shared bus.
\end{enumerate}

\subsection{Average Memory Access Time (AMAT)}
The quantitative metric governing memory subsystem performance is \textbf{Average Memory Access Time (AMAT)}:

\begin{equation}
\label{eq:amat_single}
\boxed{\text{AMAT} = \text{Hit Time} + (\text{Miss Rate} \times \text{Miss Penalty})}
\end{equation}

\textbf{Multi-Level Cache AMAT (Two-Level System):}
\begin{equation}
\label{eq:amat_multilevel}
\boxed{\text{AMAT} = T_{\text{hit\_L1}} + m_{\text{L1}} \cdot \Big[ T_{\text{hit\_L2}} + m_{\text{L2\_local}} \cdot T_{\text{memory\_penalty}} \Big]}
\end{equation}
where $m_{\text{L1}}$ is the L1 miss rate, $T_{\text{hit\_L2}}$ is the L2 access latency, and $m_{\text{L2\_local}}$ is the local miss rate of the L2 cache.

\textbf{Worked Numerical Example:}
Given: L1 Hit Time $= 1\text{ ns}$, L1 Miss Rate $= 6\%$, L2 Hit Time $= 5\text{ ns}$, L2 Local Miss Rate $= 15\%$, Main Memory Access Penalty $= 80\text{ ns}$.
Calculate AMAT.
\begin{align*}
\text{AMAT} &= 1\text{ ns} + 0.06 \times [5\text{ ns} + (0.15 \times 80\text{ ns})] \\
            &= 1 + 0.06 \times [5 + 12] \\
            &= 1 + 0.06 \times 17 \\
            &= 1 + 1.02 = \mathbf{2.02\text{ ns}}
\end{align*}


\section{Cache Organization, Mapping Techniques, and Replacement Policies}
\label{sec:cache_mapping}

A cache is partitioned into equal-size storage blocks called \textbf{Cache Lines} (typically 64 bytes in modern x86/ARM processors). Each line stores the actual data bytes along with auxiliary architectural metadata:
\begin{itemize}
    \item \textbf{Valid Bit ($V$):} Indicates whether the line contains valid program data ($1$) or uninitialized garbage ($0$).
    \item \textbf{Dirty Bit ($D$):} In write-back caches, indicates that the line was modified by a store instruction and must be written back to DRAM upon eviction.
    \item \textbf{Tag Bits:} The high-order address bits identifying which physical memory block is currently housed in that line.
\end{itemize}

\subsection{Cache Mapping Techniques}
\begin{enumerate}[label=\bfseries\arabic*.]
    \item \textbf{Direct-Mapped Cache:} Each physical memory block maps to \textit{exactly one specific cache line}.
    \begin{equation}
    \label{eq:direct_map_set}
    \text{Line Index} = (\text{Block Address}) \pmod{\text{Total Cache Lines}}
    \end{equation}
    \begin{itemize}
        \item \textit{Hardware Cost:} Lowest; requires only 1 comparator. Fast hit latency.
        \item \textit{Flaw:} Highly vulnerable to \textbf{Cache Thrashing} (conflict misses) if two actively accessed arrays map to the same set index.
    \end{itemize}
    \item \textbf{Fully Associative Cache:} A memory block can reside in \textit{any arbitrary cache line}.
    \begin{itemize}
        \item \textit{Advantage:} Zero conflict misses; optimal utilization of cache space.
        \item \textit{Hardware Cost:} Prohibitive; requires parallel comparators for every single line in the cache. Restricted to small structures (e.g., TLBs).
    \end{itemize}
    \item \textbf{$N$-Way Set-Associative Cache:} The cache is divided into $S$ sets, where each set contains exactly $N$ cache lines (ways). A memory block maps to a specific set index, but can occupy \textit{any of the $N$ ways} within that set.
    \begin{equation}
    \label{eq:set_assoc_index}
    \text{Set Index} = (\text{Block Address}) \pmod{S}
    \end{equation}
    \begin{itemize}
        \item Balances fast access latency with dramatically reduced conflict misses.
        \item Standard industry configurations: 4-way, 8-way, or 16-way set associativity.
    \end{itemize}
\end{enumerate}

\subsection{Physical Address Breakdown Formulas}
A physical memory address of bit-width $W$ is mathematically partitioned into three fields:
\begin{equation}
\label{eq:address_partition}
W = \text{Tag Bits} + \text{Set Index Bits} + \text{Byte Offset Bits}
\end{equation}
\begin{enumerate}[label=\bfseries\arabic*.]
    \item \textbf{Byte Offset ($b$ bits):} Selects the individual byte within the cache line of size $B$ bytes:
    \begin{equation}
    \label{eq:offset_formula}
    b = \log_2(B)
    \end{equation}
    \item \textbf{Set Index ($s$ bits):} Selects the target set among $S$ total sets:
    \begin{equation}
    \label{eq:set_formula}
    S = \frac{\text{Total Cache Capacity}}{\text{Associativity } N \times \text{Line Size } B} \implies s = \log_2(S)
    \end{equation}
    \item \textbf{Tag ($t$ bits):} Uniquely identifies the physical block residing in the set:
    \begin{equation}
    \label{eq:tag_formula}
    t = W - s - b
    \end{equation}
\end{enumerate}

\begin{table}[htbp]
\centering
\small
\caption{Comparative Analysis: Cache Mapping Architectures}
\label{tab:cache_mapping}
\begin{tabularx}{\textwidth}{l X X X}
\toprule
\textbf{Dimension} & \textbf{Direct Mapped} & \textbf{Fully Associative} & \textbf{$N$-Way Set Associative} \\
\midrule
\textbf{Line Placement} & Exactly 1 predetermined line. & Any arbitrary line in cache. & Any of $N$ lines in target set. \\
\textbf{Index Bits} & $s = \log_2(\text{Total Lines})$. & $0$ index bits (No index). & $s = \log_2(\text{Total Sets})$. \\
\textbf{Hardware Comparators} & Exactly 1 comparator. & $K$ comparators (Total lines). & Exactly $N$ comparators. \\
\textbf{Conflict Miss Rate} & Highest (severe thrashing). & Zero conflict misses. & Low to moderate. \\
\textbf{Replacement Policy} & None (unconditional overwrite). & Mandatory (LRU, FIFO, etc.). & Mandatory within the set. \\
\bottomrule
\end{tabularx}
\end{table}

\subsection{Cache Replacement Policies}
When a miss occurs in an $N$-way set-associative cache and all $N$ ways in the target set are occupied, a line must be evicted:
\begin{itemize}
    \item \textbf{Least Recently Used (LRU):} Evicts the block that has gone unreferenced for the longest duration. Closest to optimal temporal locality exploitation, but requires tracking bits ($O(\log N)$ or matrix flip-flops).
    \item \textbf{First-In, First-Out (FIFO):} Evicts the block that has resided in the cache the longest, regardless of recent access frequency.
    \item \textbf{Least Frequently Used (LFU):} Evicts the block with the smallest reference count.
    \item \textbf{Random Replacement:} Selects an eviction victim pseudo-randomly. Minimal hardware overhead; performance approaches LRU under large caches.
\end{itemize}

\subsection{Cache Write Policies}
\begin{enumerate}[label=\bfseries\arabic*.]
    \item \textbf{Write-Through:} Every store operation writes immediately to both the cache and lower-tier DRAM.
    \begin{itemize}
        \item Ensures strict memory consistency, but consumes significant bus write bandwidth.
        \item Mitigated by a \textbf{Write Buffer} to avoid stalling the CPU.
    \end{itemize}
    \item \textbf{Write-Back (Copy-Back):} Stores update only the cache line, setting its internal \textbf{Dirty Bit} ($D=1$). DRAM is updated only when the dirty block is evicted.
    \begin{itemize}
        \item Drastically minimizes bus write bandwidth; preferred standard in high-performance processors.
    \end{itemize}
    \item \textbf{Write Allocation Policies (On Write Miss):}
    \begin{itemize}
        \item \textbf{Write-Allocate (Fetch on Write):} The missed block is fetched from DRAM into cache, then the write is applied. (Standard pairing with Write-Back).
        \item \textbf{No-Write-Allocate (Write-Around):} The write modifies DRAM directly without loading the block into cache. (Standard pairing with Write-Through).
    \end{itemize}
\end{enumerate}


\section{Cache Coherence Protocols and the MESI State Machine}
\label{sec:cache_coherence}

\subsection{The Multicore Cache Coherence Problem}
In symmetric multiprocessing (SMP) and multi-core processors, each CPU core possesses its own private Level-1 and Level-2 caches, but all cores share a unified physical DRAM address space.
\begin{itemize}
    \item If Core~0 writes a value ($X \leftarrow 42$) into its local write-back cache, Core~1 reading address $X$ from its own cache or DRAM will read stale data ($X = 10$).
    \item \textbf{Cache Coherence Protocols} guarantee that all processors observe a consistent view of memory.
\end{itemize}

\subsection{Coherence Hardware Paradigms}
\begin{itemize}
    \item \textbf{Snooping (Bus-Based):} Every cache controller continuously ``snoops'' (monitors) the shared broadcast bus for address transactions issued by peer cores. Used in small-scale multicore systems ($<16$ cores).
    \item \textbf{Directory-Based:} A centralized or distributed directory tracks the caching state and location of every memory block across the system. Scales efficiently to hundreds of nodes (NUMA supercomputers).
\end{itemize}

\subsection{The MESI Protocol (Illinois Protocol)}
The \textbf{MESI protocol} is an invalidation-based cache coherence state machine wherein every cache line resides in one of four mutually exclusive states:

\begin{enumerate}[label=\bfseries\arabic*.]
    \item \textbf{Modified (M):} The cache line is present \textit{only} in the current core's cache, and its data is \textbf{dirty} (has been modified relative to DRAM). The core holds exclusive ownership and can read/write without bus notification.
    \item \textbf{Exclusive (E):} The cache line is present \textit{only} in the current core's cache, but its data is \textbf{clean} (identical to DRAM). The core can transition silently to Modified on a write.
    \item \textbf{Shared (S):} The cache line may be replicated across multiple cores' caches. Its data is clean relative to DRAM (or relative to the owner). Writes are forbidden without issuing a bus invalidation request.
    \item \textbf{Invalid (I):} The cache line contains invalid, stale data; reading from it constitutes a cache miss.
\end{enumerate}

\begin{table}[htbp]
\centering
\small
\caption{The Four Operational States of the MESI Coherence Protocol}
\label{tab:mesi_states}
\begin{tabularx}{\textwidth}{l c c X}
\toprule
\textbf{MESI State} & \textbf{Line Valid?} & \textbf{Data Clean or Dirty?} & \textbf{Present in Peer Caches?} \\
\midrule
\textbf{M (Modified)} & Yes & \textbf{Dirty} (Must write back to DRAM on eviction) & No (Exclusively owned by this core) \\
\textbf{E (Exclusive)} & Yes & \textbf{Clean} (Matches DRAM) & No (Exclusively owned by this core) \\
\textbf{S (Shared)} & Yes & \textbf{Clean} (Matches DRAM or Owner) & Yes (Shared across $\ge 1$ peer cores) \\
\textbf{I (Invalid)} & No & Invalid / Out-of-date & May be present elsewhere \\
\bottomrule
\end{tabularx}
\end{table}


\section{Random Access Memory: SRAM versus DRAM and DDR Evolution}
\label{sec:ram_technologies}

\textbf{Random Access Memory (RAM)} provides volatile primary storage where any arbitrary memory location can be accessed directly in constant time $O(1)$, independent of physical location.

\subsection{Static RAM (SRAM) versus Dynamic RAM (DRAM)}
\begin{enumerate}[label=\bfseries\arabic*.]
    \item \textbf{Static RAM (SRAM):}
    \begin{itemize}
        \item \textit{Bit Cell Design:} Uses a 6-transistor (6T) CMOS cross-coupled bistable latch.
        \item \textit{Operation:} Retains its stored binary state continuously as long as DC power is supplied. No charge leakage; zero refresh cycles required.
        \item \textit{Performance:} Ultra-fast access times ($0.5$--$2.5\text{ ns}$).
        \item \textit{Trade-offs:} Large silicon area per bit cell ($>6\times$ larger than DRAM); low density; high power and thermal dissipation; expensive.
        \item \textit{Role:} Level-1, Level-2, and Level-3 on-chip CPU caches.
    \end{itemize}
    \item \textbf{Dynamic RAM (DRAM):}
    \begin{itemize}
        \item \textit{Bit Cell Design:} Uses a minimalistic 1-transistor, 1-capacitor (1T-1C) cell.
        \item \textit{Operation:} Stores binary data as an electrostatic charge on a micro-capacitor ($C \sim 20\text{ fF}$). Because charge leaks continuously through sub-threshold silicon paths, \textbf{periodic refresh cycles} (reading and rewriting every row every $64\text{ ms}$) are mandatory to prevent data loss.
        \item \textit{Destructive Read:} Reading a DRAM cell discharges the capacitor; the sensing amplifier must immediately rewrite the row buffer.
        \item \textit{Performance:} Slower access latency ($12$--$18\text{ ns}$).
        \item \textit{Trade-offs:} Extreme packing density (billions of bits per chip); low cost per gigabyte.
        \item \textit{Role:} Primary system memory (DIMMs).
    \end{itemize}
\end{enumerate}

\begin{table}[htbp]
\centering
\small
\caption{Comprehensive Comparison: Static RAM (SRAM) versus Dynamic RAM (DRAM)}
\label{tab:sram_vs_dram}
\begin{tabularx}{\textwidth}{l X X}
\toprule
\textbf{Specification} & \textbf{Static RAM (SRAM)} & \textbf{Dynamic RAM (DRAM)} \\
\midrule
\textbf{Cell Circuitry} & 6 CMOS transistors (cross-coupled inverters). & 1 transistor and 1 storage capacitor (1T-1C). \\
\textbf{Storage Mechanism} & Bistable flip-flop voltage latching. & Electrostatic charge on tiny capacitor. \\
\textbf{Refresh Requirement} & None. Data is static while powered. & Mandatory periodic refresh (every 32--64 ms). \\
\textbf{Access Speed} & Ultra-fast ($0.5$ to $3.0\text{ ns}$). & Moderate ($10$ to $20\text{ ns}$). \\
\textbf{Packing Density} & Low (large physical cell area). & Extremely high (compact 1T-1C cell). \\
\textbf{Power Dissipation} & High active power; low standby leakage. & Low active power; requires refresh power. \\
\textbf{Cost per Bit} & Very high. & Low. \\
\textbf{Primary Role} & CPU L1/L2/L3 cache, register files. & Main computer system memory (RAM). \\
\bottomrule
\end{tabularx}
\end{table}

\subsection{Evolution of Synchronous DRAM and DDR Standards}
\begin{itemize}
    \item \textbf{SDR SDRAM (Single Data Rate):} Synchronized with the external memory bus clock; transfers data strictly on the rising clock edge (1 transfer/cycle).
    \item \textbf{DDR SDRAM (Double Data Rate):} Doubles bus throughput by transferring data on \textbf{both the rising and falling edges} of the differential clock (2 transfers/cycle).
    \item \textbf{Generational Architectural Scaling:}
    \begin{itemize}
        \item \textbf{DDR1:} 2n-bit prefetch buffer $\cdot$ Operating voltage: $2.5\text{ V}$.
        \item \textbf{DDR2:} 4n-bit prefetch buffer $\cdot$ Operating voltage: $1.8\text{ V}$.
        \item \textbf{DDR3:} 8n-bit prefetch buffer $\cdot$ Operating voltage: $1.5\text{ V}$.
        \item \textbf{DDR4:} 8n-bit prefetch with bank groups $\cdot$ Operating voltage: $1.2\text{ V}$.
        \item \textbf{DDR5:} 16n-bit prefetch $\cdot$ Dual 32-bit sub-channels $\cdot$ Operating voltage: $1.1\text{ V}$ with on-die PMIC and on-die ECC.
    \end{itemize}
\end{itemize}


\section{Read-Only Memory and Non-Volatile Semiconductor Storage}
\label{sec:rom_technologies}

\textbf{Read-Only Memory (ROM)} provides persistent, non-volatile storage that retains its contents permanently without electrical power.

\begin{enumerate}[label=\bfseries\arabic*.]
    \item \textbf{Mask ROM (MROM):} Hardwired during semiconductor lithographic fabrication. Permanently unalterable. Used for high-volume, low-cost microcode.
    \item \textbf{PROM (Programmable ROM):} Shipped blank with microscopic semiconductor fuse links. Programmed once by burning open selected fuses using high-voltage current pulses (\textbf{One-Time Programmable, OTP}).
    \item \textbf{EPROM (Erasable PROM):} Uses floating-gate MOS transistors. Erasable by exposing the physical silicon chip through a transparent quartz window to intense \textbf{Ultraviolet (UV) radiation} for 20--30 minutes. All contents are erased uniformly.
    \item \textbf{EEPROM (Electrically Erasable PROM):} Utilizes Fowler-Nordheim quantum tunneling to erase and reprogram data \textbf{in-circuit electrically at the byte level}. Re-programmable $\sim 10^5$ to $10^6$ times.
    \item \textbf{Flash Memory:} A high-density evolution of EEPROM that erases data in large \textbf{blocks} rather than single bytes:
    \begin{itemize}
        \item \textbf{NOR Flash:} Allows byte-level random read access; enables \textbf{Execute-in-Place (XIP)} for CPU firmware. Low density, higher cost.
        \item \textbf{NAND Flash:} Page-based read/write and block-based erase. High packing density, low cost. Serves as the storage foundation for USB drives, SD cards, and Solid-State Drives (SSDs).
    \end{itemize}
\end{enumerate}


\section{RAM versus ROM: Comprehensive Comparative Analysis}
\label{sec:ram_vs_rom}

\begin{table}[htbp]
\centering
\small
\caption{Comprehensive Architectural Analysis: RAM versus ROM}
\label{tab:ram_vs_rom}
\begin{tabularx}{\textwidth}{l X X}
\toprule
\textbf{Dimension} & \textbf{Random Access Memory (RAM)} & \textbf{Read-Only Memory (ROM)} \\
\midrule
\textbf{Data Volatility} & Volatile (Loses all contents immediately upon power termination). & Non-volatile (Retains data indefinitely without power). \\
\textbf{Read/Write Operations} & High-speed read and write operations supported continuously. & Read-dominant; writing requires specialized block-erase cycles. \\
\textbf{Execution Speed} & Extremely fast ($0.5$ to $20\text{ ns}$). & Slower than RAM ($15$ to $100\text{ ns}$). \\
\textbf{Primary Role} & Working execution storage for active OS processes, stacks, and heaps. & Boot firmware (BIOS, UEFI), bootloader, embedded firmware. \\
\textbf{Storage Lifespan} & Infinite read/write cycles (no silicon wear). & Limited write/erase endurance cycles ($10^3$--$10^6$ writes). \\
\bottomrule
\end{tabularx}
\end{table}


\section{Solid-State Drives versus Mechanical Hard Disk Drives}
\label{sec:ssd_vs_hdd}

\subsection{Mechanical Hard Disk Drives (HDD)}
An HDD records data magnetically on concentric circular \textbf{tracks} divided into \textbf{sectors} across rapidly rotating rigid magnetic platters ($5400$--$15000\text{ RPM}$).
\begin{itemize}
    \item \textbf{Access Time Components:}
    \begin{equation}
    \label{eq:hdd_access_time}
    \text{Total HDD Access Time} = T_{\text{seek}} + T_{\text{rotational\_latency}} + T_{\text{transfer}}
    \end{equation}
    \begin{enumerate}[label=\bfseries\alph*.]
        \item \textbf{Seek Time ($T_{\text{seek}}$):} The mechanical time required for the actuator arm to physically move the read/write head over the desired track ($3$--$10\text{ ms}$).
        \item \textbf{Rotational Latency ($T_{\text{rotational}}$):} The time taken for the target sector to rotate underneath the head. On average, equals half a revolution:
        \begin{equation}
        \label{eq:rotational_latency}
        T_{\text{rotational\_avg}} = \frac{1}{2} \cdot \left( \frac{60}{\text{RPM}} \right) \text{ seconds}
        \end{equation}
        For a $7200\text{ RPM}$ disk: $T_{\text{rotational\_avg}} = \frac{60}{2 \times 7200} = 4.17\text{ ms}$.
        \item \textbf{Data Transfer Time ($T_{\text{transfer}}$):} Time to read the data bits at platter surface speed.
    \end{enumerate}
\end{itemize}

\subsection{Solid-State Drives (SSD)}
SSDs eliminate all moving mechanical parts, storing data across multi-channel arrays of \textbf{NAND Flash memory} chips managed by an intelligent embedded \textbf{Flash Controller}.
\begin{itemize}
    \item \textbf{Zero Mechanical Latency:} Access time is purely electronic ($10$--$50\ \mu\text{s}$), delivering hundreds of thousands of Input/Output Operations Per Second (\textbf{IOPS}).
    \item \textbf{Erase-Before-Write Constraint:} NAND flash cannot overwrite existing data. Data is written in \textbf{Pages} ($4$--$16\text{ KB}$), but can only be erased in whole \textbf{Blocks} ($2$--$8\text{ MB}$).
    \item \textbf{Wear Leveling:} The controller dynamically remaps writes across physical blocks to prevent localized cell degradation.
    \item \textbf{TRIM Command:} The OS informs the SSD which deleted file sectors are no longer needed, allowing background garbage collection.
\end{itemize}


\section{Unified Storage Hierarchy Matrix}
\label{sec:storage_matrix}

\begin{table}[htbp]
\centering
\small
\caption{Master Comparison of the Digital Storage Hierarchy}
\label{tab:master_storage_matrix}
\begin{tabularx}{\textwidth}{l c c c X}
\toprule
\textbf{Tier} & \textbf{Technology} & \textbf{Typical Latency} & \textbf{Typical Capacity} & \textbf{Primary System Purpose} \\
\midrule
\textbf{CPU Registers} & CMOS Flip-Flops & $<0.5\text{ ns}$ & $<2\text{ KiB}$ & Active CPU execution operands. \\
\textbf{L1 Cache} & SRAM (6T) & $1\text{ ns}$ & $32$--$64\text{ KiB}$/core & Immediate instruction/data caching. \\
\textbf{L2 Cache} & SRAM (6T) & $3$--$5\text{ ns}$ & $512\text{ KiB}$--$2\text{ MiB}$ & Secondary core caching. \\
\textbf{L3 Cache} & SRAM (6T) & $10$--$20\text{ ns}$ & $16$--$128\text{ MiB}$ & Last-level shared multicore cache. \\
\textbf{Main Memory} & DRAM (1T-1C DDR5) & $50$--$80\text{ ns}$ & $16$--$128\text{ GiB}$ & OS kernel and application workspace. \\
\textbf{NVMe SSD} & NAND Flash & $10$--$50\ \mu\text{s}$ & $512\text{ GiB}$--$4\text{ TiB}$ & High-speed persistent operating system. \\
\textbf{Mechanical HDD} & Magnetic Platters & $5$--$15\text{ ms}$ & $2$--$24\text{ TiB}$ & Bulk storage and archival backup. \\
\bottomrule
\end{tabularx}
\end{table}



\section{Processor Performance Metrics and Amdahl's Law}
\label{sec:processor_performance}

Quantitative computer architecture evaluates processor efficiency using rigorous execution time formulations rather than superficial clock frequency figures.

\subsection{The Fundamental CPU Performance Equation}
The total CPU execution time ($T_{\text{CPU}}$) required to execute a software program is governed by the \textbf{Iron Law of Processor Performance}:

\begin{equation}
\label{eq:cpu_time_formula}
\boxed{T_{\text{CPU}} = \text{Instruction Count (IC)} \times \text{Cycles Per Instruction (CPI)} \times \text{Clock Cycle Time } \tau = \frac{\text{IC} \times \text{CPI}}{f_{\text{clock}}}}
\end{equation}

Each factor is governed by distinct engineering disciplines:
\begin{itemize}
    \item \textbf{Instruction Count (IC):} Determined by the algorithm design, high-level programming language, and compiler optimization flags.
    \item \textbf{Cycles Per Instruction (CPI):} Determined by the Instruction Set Architecture (ISA) and processor microarchitecture (pipelining depth, branch prediction, cache hit rates, superscalar issue width).
    \item \textbf{Clock Frequency ($f_{\text{clock}} = 1/\tau$):} Determined by semiconductor fabrication process node, transistor gate delay, circuit design, and thermal limits.
\end{itemize}

\subsection{Weighted Average CPI}
In modern processors, different instruction classes require varying numbers of execution cycles:
\begin{equation}
\label{eq:weighted_cpi}
\text{CPI}_{\text{avg}} = \sum_{i=1}^{m} \left( \text{CPI}_i \times \frac{\text{IC}_i}{\text{IC}_{\text{total}}} \right) = \sum_{i=1}^{m} (\text{CPI}_i \times f_i)
\end{equation}
where $f_i$ is the fractional frequency of instruction class $i$.

\textbf{Worked Example:} An application workload consists of:
\begin{itemize}
    \item $50\%$ Arithmetic instructions with $\text{CPI} = 1$.
    \item $30\%$ Memory instructions with $\text{CPI} = 3$.
    \item $20\%$ Branch instructions with $\text{CPI} = 2$.
\end{itemize}
Calculate the weighted average CPI:
\[ \text{CPI}_{\text{avg}} = (0.50 \times 1) + (0.30 \times 3) + (0.20 \times 2) = 0.50 + 0.90 + 0.40 = \mathbf{1.80} \]

\subsection{MIPS and FLOPS Metrics}
\begin{itemize}
    \item \textbf{MIPS (Million Instructions Per Second):}
    \begin{equation}
    \label{eq:mips_formula}
    \text{MIPS} = \frac{\text{Instruction Count}}{\text{Execution Time} \times 10^6} = \frac{f_{\text{clock}}}{\text{CPI} \times 10^6}
    \end{equation}
    \textit{Limitation:} MIPS does not account for differing instruction capabilities between ISAs (a RISC machine may require more instructions than a CISC machine for the same task).
    \item \textbf{FLOPS (Floating-Point Operations Per Second):} Standard metric in scientific supercomputing and machine learning (e.g., GFLOPS, TFLOPS, PFLOPS).
\end{itemize}

\subsection{Amdahl's Law}
Formulated by Gene Amdahl in 1967, \textbf{Amdahl's Law} models the theoretical maximum speedup achievable by parallelizing or enhancing a fraction of a computer system:

\begin{equation}
\label{eq:amdahls_law}
\boxed{S_{\text{overall}} = \frac{1}{(1 - f) + \dfrac{f}{s}}}
\end{equation}
where:
\begin{itemize}
    \item $f$ is the fraction of the execution workload that can be enhanced or parallelized ($0 \le f \le 1$).
    \item $(1 - f)$ is the strictly serial fraction that cannot be parallelized.
    \item $s$ is the speedup factor applied to the enhanced fraction.
\end{itemize}

\textbf{Asymptotic Limit of Parallel Speedup:}
As the number of parallel processor cores approaches infinity ($s \to \infty$):
\begin{equation}
\label{eq:amdahl_limit}
\lim_{s \to \infty} S_{\text{overall}} = \frac{1}{1 - f}
\end{equation}
\textbf{Critical Engineering Takeaway:} If just $10\%$ of an application is inherently serial ($1 - f = 0.10$), the maximum theoretical speedup achievable on an infinite number of parallel CPU cores is strictly capped at:
\[ S_{\text{max}} = \frac{1}{0.10} = 10\times \]


\section{Multicore Concurrency, Hardware Multithreading, and Flynn's Taxonomy}
\label{sec:multicore_parallelism}

\subsection{Multicore versus Multiprocessor Systems}
When power dissipation and thermal ceilings (the \textbf{Power Wall}) halted unconstrained CPU clock frequency scaling in the mid-2000s, semiconductor architects pivoted from single-core frequency scaling to multi-core parallelism.
\begin{itemize}
    \item \textbf{Symmetric Multiprocessing (SMP):} Multiple separate CPU chips connected across an external motherboard bus.
    \item \textbf{Chip Multiprocessor (CMP / Multicore):} Multiple independent execution cores integrated onto a single monolithic semiconductor die, sharing an on-die L3 cache and high-bandwidth interconnects (ring buses or 2D mesh fabrics).
\end{itemize}

\subsection{Hardware Multithreading and Simultaneous Multithreading (SMT)}
Execution units inside a high-performance out-of-order CPU core are frequently underutilized due to long-latency cache misses or data dependency stalls.
\begin{itemize}
    \item \textbf{Coarse-Grained Multithreading:} Switches execution to an alternate thread only upon a catastrophic long-latency stall (e.g., an L3 cache miss).
    \item \textbf{Fine-Grained Multithreading:} Interleaves instructions from different threads cycle-by-cycle in a round-robin schedule.
    \item \textbf{Simultaneous Multithreading (SMT / Intel Hyper-Threading):} Replicates only the architectural register state (PC, register file, flags) while \textbf{sharing the underlying physical execution units, caches, and pipelines}. Enables instructions from two or more independent threads to execute concurrently within the same clock cycle, increasing core utilization by $15$--$35\%$.
\end{itemize}

\subsection{Flynn's Taxonomy}
Michael J.~Flynn classified all computer architectures into four operational quadrants based on the multiplicity of hardware Instruction Streams and Data Streams:

\begin{table}[htbp]
\centering
\small
\caption{Flynn's Architectural Taxonomy of Computer Systems}
\label{tab:flynns_taxonomy}
\begin{tabularx}{\textwidth}{l l X X}
\toprule
\textbf{Category} & \textbf{Full Form} & \textbf{Operational Principle} & \textbf{Physical Examples} \\
\midrule
\textbf{SISD} & Single Instruction, Single Data & Classic sequential processor; 1 instruction stream operates on 1 data stream. & Classic single-core CPU (Intel 8086, PDP-11). \\
\textbf{SIMD} & Single Instruction, Multiple Data & A single machine instruction operates concurrently across a vector of data operands. & Vector processors, GPUs, CPU vector extensions (AVX-512, ARM Neon). \\
\textbf{MISD} & Multiple Instruction, Single Data & Multiple distinct instruction streams operate concurrently on the identical data stream. & Fault-tolerant redundant systems (Space Shuttle flight computers, systolic arrays). \\
\textbf{MIMD} & Multiple Instruction, Multiple Data & Multiple autonomous processors execute distinct instruction streams on distinct data. & Modern multicore CPUs, distributed HPC server clusters, cloud datacenters. \\
\bottomrule
\end{tabularx}
\end{table}


\section{Endianness: Big-Endian versus Little-Endian Architecture}
\label{sec:endianness}

\textbf{Endianness} defines the architectural byte ordering convention used to store multi-byte data types (e.g., 16-bit, 32-bit, or 64-bit integers) in byte-addressable physical memory.

\subsection{Architectural Definitions}
\begin{enumerate}[label=\bfseries\arabic*.]
    \item \textbf{Big-Endian:} The \textbf{Most Significant Byte (MSB)} is stored at the lowest memory address (natural human reading order from left to right).
    \item \textbf{Little-Endian:} The \textbf{Least Significant Byte (LSB)} is stored at the lowest memory address.
\end{enumerate}

\textbf{Concrete Example:} Consider storing the 32-bit hexadecimal integer \texttt{0x12345678} at memory address \texttt{0x1000}:
\begin{itemize}
    \item Byte breakdown: $\text{MSB} = \texttt{0x12}, \quad \text{Byte 2} = \texttt{0x34}, \quad \text{Byte 1} = \texttt{0x56}, \quad \text{LSB} = \texttt{0x78}$.
\end{itemize}

\begin{table}[htbp]
\centering
\small
\caption{Physical Memory Byte Layout for Value \texttt{0x12345678}}
\label{tab:endian_layout}
\begin{tabularx}{0.7\textwidth}{c c c}
\toprule
\textbf{Memory Address} & \textbf{Big-Endian Layout} & \textbf{Little-Endian Layout} \\
\midrule
\texttt{0x1000} & \texttt{0x12} (MSB) & \texttt{0x78} (LSB) \\
\texttt{0x1001} & \texttt{0x34} & \texttt{0x56} \\
\texttt{0x1002} & \texttt{0x56} & \texttt{0x34} \\
\texttt{0x1003} & \texttt{0x78} (LSB) & \texttt{0x12} (MSB) \\
\bottomrule
\end{tabularx}
\end{table}

\subsection{Architectural Adoption and Network Byte Order}
\begin{itemize}
    \item \textbf{Little-Endian Platforms:} Intel/AMD x86-64, ARM (default mode), RISC-V.
    \item \textbf{Big-Endian Platforms:} IBM z/Architecture, Motorola 68000, SPARC.
    \item \textbf{Network Byte Order:} The Internet Protocol (TCP/IP) suite strictly standardizes on \textbf{Big-Endian} for all packet headers. Host systems must execute conversion routines:
    \begin{itemize}
        \item \texttt{htons()} / \texttt{htonl()}: Host to Network Short (16-bit) / Long (32-bit).
        \item \texttt{ntohs()} / \texttt{ntohl()}: Network to Host Short (16-bit) / Long (32-bit).
    \end{itemize}
\end{itemize}


\section{Logic Families, Circuit Timing, and Digital Interfaces}
\label{sec:logic_families}

\subsection{TTL versus CMOS Logic Families}
Integrated circuits are realized through two primary semiconductor logic families:
\begin{enumerate}[label=\bfseries\arabic*.]
    \item \textbf{Transistor-Transistor Logic (TTL):} Constructed using Bipolar Junction Transistors (BJTs). Characterized by high switching speed, low input impedance, and high static power consumption ($V_{CC} = +5\text{ V}$).
    \item \textbf{Complementary Metal-Oxide Semiconductor (CMOS):} Constructed using complementary pairs of nMOS and pMOS Field-Effect Transistors (FETs).
    \begin{itemize}
        \item \textit{Near-Zero Static Power:} In steady state, either the pull-up pMOS or pull-down nMOS is completely off, drawing only nanoampere leakage currents.
        \item \textit{Dynamic Power:} Power is dissipated almost exclusively during logic transitions charging gate capacitances ($P = C \cdot V^2 \cdot f$).
        \item Dominates $>99\%$ of modern VLSI microprocessors, memory, and ASICs.
    \end{itemize}
\end{enumerate}

\begin{table}[htbp]
\centering
\small
\caption{Comparative Analysis: TTL versus CMOS Digital Logic Families}
\label{tab:ttl_vs_cmos}
\begin{tabularx}{\textwidth}{l X X}
\toprule
\textbf{Parameter} & \textbf{TTL (Bipolar BJT)} & \textbf{CMOS (MOSFET)} \\
\midrule
\textbf{Core Devices} & Bipolar Junction Transistors (NPN/PNP). & Complementary pMOS and nMOS pairs. \\
\textbf{Supply Voltage ($V_{DD}$)} & Fixed $+5\text{ V} \pm 10\%$. & Wide range ($0.8\text{ V}$ to $5\text{ V}$); modern VLSI uses $<1.2\text{ V}$. \\
\textbf{Static Power Dissipation} & High ($1$ to $10\text{ mW}$ per gate continuously). & Negligible (nanowatts in steady state). \\
\textbf{Noise Margin} & Moderate ($0.4\text{ V}$). & Exceptionally high (approaches $30$--$40\%$ of $V_{DD}$). \\
\textbf{Input Impedance} & Low ($\sim 10\text{ k}\Omega$). & Extremely high ($>10^{12}\ \Omega$, capacitive gate oxide). \\
\textbf{Fan-Out} & Typically 10 standard TTL loads. & $>50$ at low frequencies (limited by capacitive load). \\
\bottomrule
\end{tabularx}
\end{table}

\subsection{Fan-In and Fan-Out Specifications}
\begin{itemize}
    \item \textbf{Fan-In:} The maximum number of physical inputs a logic gate can accept without degrading electrical switching thresholds or propagation delays.
    \item \textbf{Fan-Out:} The maximum number of standard load inputs of the same logic family that a gate output can reliably drive simultaneously while maintaining valid output voltage logic levels ($V_{OH}$ and $V_{OL}$).
    \begin{equation}
    \label{eq:fan_out}
    \text{Fan-Out} = \min \left( \left\lfloor \frac{I_{OH}}{I_{IH}} \right\rfloor, \left\lfloor \frac{I_{OL}}{I_{IL}} \right\rfloor \right)
    \end{equation}
\end{itemize}

\subsection{Combinational Hazards and Glitches}
A \textbf{Hazard} is an unwanted, transient switching spike (glitch) on a combinational output line caused by unequal propagation delays along converging circuit paths:
\begin{enumerate}[label=\bfseries\arabic*.]
    \item \textbf{Static-1 Hazard:} The output is expected to remain constant at logic $1$, but momentarily drops to $0$ before settling back to $1$.
    \item \textbf{Static-0 Hazard:} The output is expected to remain constant at logic $0$, but momentarily spikes to $1$ before settling back to $0$.
    \item \textbf{Dynamic Hazard:} The output is scheduled to transition cleanly once ($0 \to 1$ or $1 \to 0$), but oscillates multiple times ($0 \to 1 \to 0 \to 1$) due to multi-level path disparities.
    \item \textbf{Mitigation (Hazard Covers):} Static-1 hazards are eliminated on a K-map by adding redundant consensus prime implicants that bridge adjacent grouping loops.
\end{enumerate}


\section{Data Converters (ADC, DAC) and the Nyquist Sampling Theorem}
\label{sec:data_converters}

Microprocessors operate strictly on discrete binary values, whereas the physical universe operates via continuous analog voltages, pressures, and temperatures.

\subsection{Analog-to-Digital Conversion (ADC)}
An ADC performs three sequential operations:
\begin{enumerate}[label=\bfseries\arabic*.]
    \item \textbf{Sampling:} Converting a continuous-time signal into a discrete-time signal at periodic sampling intervals ($T_s$).
    \item \textbf{Quantization:} Discretizing the continuous analog amplitude into one of $2^n$ finite voltage bins.
    \item \textbf{Encoding:} Converting each quantized bin level into an $n$-bit binary word.
\end{enumerate}

\textbf{ADC Step Size and Resolution:}
For an $n$-bit ADC with full-scale reference voltage $V_{\text{ref}}$:
\begin{equation}
\label{eq:adc_step_size}
\boxed{\text{Step Size } (\Delta V) = \frac{V_{\text{ref}}}{2^n - 1} \approx \frac{V_{\text{ref}}}{2^n}}
\end{equation}

\subsection{Major ADC Architectures}
\begin{itemize}
    \item \textbf{Flash ADC:} Uses $2^n - 1$ analog comparators operating in parallel across a precision resistor voltage divider. Completes conversion in \textbf{1 clock cycle} ($O(1)$ latency). Fastest ADC available (gigahertz sampling), but exponentially high hardware cost; limited to $6$--$8$ bits.
    \item \textbf{Successive Approximation Register (SAR) ADC:} Uses a binary search algorithm with an internal DAC and a single comparator. Converges in exactly $n$ clock cycles for an $n$-bit conversion ($O(n)$ latency). Industry workhorse for microcontrollers ($10$--$16$ bits).
    \item \textbf{Dual-Slope Integrating ADC:} Integrates input voltage across a capacitor, then de-integrates using a reference. Slow ($2^{n+1}$ cycles), but delivers extreme precision ($16$--$24$ bits) and perfect rejection of $50/60\text{ Hz}$ power-line noise (used in digital multimeters).
\end{itemize}

\subsection{Digital-to-Analog Conversion (DAC)}
Converts an $n$-bit binary digital value into an analog voltage:
\begin{itemize}
    \item \textbf{Binary-Weighted Resistor DAC:} Uses parallel resistors with values $R, 2R, 4R, 8R, \dots$ summed at an operational amplifier virtual ground. Poor precision for large $n$ due to extreme resistance ranges ($>1000:1$).
    \item \textbf{R-2R Ladder DAC:} Uses only \textbf{two precision resistor values} ($R$ and $2R$) in a repeating ladder network. Universal standard in integrated circuits.
    \begin{equation}
    \label{eq:r2r_dac}
    V_{\text{out}} = V_{\text{ref}} \sum_{i=1}^{n} b_i \cdot 2^{-i}
    \end{equation}
\end{itemize}

\subsection{The Nyquist-Shannon Sampling Theorem}
The fundamental mathematical law of digital signal processing states:
\begin{equation}
\label{eq:nyquist_theorem}
\boxed{f_s \ge 2 \cdot f_{\text{max}}}
\end{equation}
where $f_s$ is the sampling frequency and $f_{\text{max}}$ is the maximum frequency component present in the analog signal.
\begin{itemize}
    \item \textbf{Nyquist Rate:} The minimum mandatory sampling frequency $f_{\text{Nyquist}} = 2 f_{\text{max}}$.
    \item \textbf{Aliasing:} If $f_s < 2 f_{\text{max}}$, high-frequency signal components fold back into the lower baseband, masquerading as false low frequencies and irrevocably corrupting the reconstructed signal.
    \item \textbf{Anti-Aliasing Filter:} An analog low-pass filter placed immediately before the ADC to eliminate all frequency components above $f_s / 2$.
\end{itemize}


\section{Language Translation: Compilers, Interpreters, Assemblers, Linkers, and Loaders}
\label{sec:translation_systems}

\subsection{The Translation Spectrum}
\begin{enumerate}[label=\bfseries\arabic*.]
    \item \textbf{Compiler:} Translates high-level source code (C, C++, Rust, Fortran) entirely into native machine code or object files ahead of time (Ahead-Of-Time, AOT). Fast execution; platform-dependent executable.
    \item \textbf{Interpreter:} Analyzes, parses, and executes high-level source code line-by-line in real time (e.g., standard Python, Ruby). Portable and interactive, but incurs runtime interpretation overhead.
    \item \textbf{Hybrid / Bytecode JIT (Just-In-Time):} Compiles source into intermediate platform-neutral bytecode (e.g., Java \texttt{.class}, C\# CIL), which is executed by a Virtual Machine (JVM, .NET CLR) and dynamically compiled into native machine code during execution.
    \item \textbf{Assembler:} A low-level translator that converts assembly language mnemonics (\texttt{MOV}, \texttt{ADD}) directly into binary machine opcodes on a 1-to-1 basis.
\end{enumerate}

\subsection{Phases of the Compilation Pipeline}
\begin{enumerate}[label=\bfseries\arabic*.]
    \item \textbf{Lexical Analysis (Scanning / Tokenization):} Converts the raw input stream of characters into recognized lexical tokens (identifiers, keywords, operators, literals).
    \item \textbf{Syntax Analysis (Parsing):} Verifies token sequences against grammar rules and constructs an \textbf{Abstract Syntax Tree (AST)}.
    \item \textbf{Semantic Analysis:} Enforces language semantics, type-checking, and variable scoping rules.
    \item \textbf{Intermediate Code Generation (IR):} Produces platform-independent three-address code (e.g., LLVM IR).
    \item \textbf{Code Optimization:} Eliminates dead code, unrolls loops, propagates constants, and optimizes register allocation.
    \item \textbf{Target Code Generation:} Emits machine instructions for the target ISA.
\end{enumerate}

\subsection{Linker versus Loader}
\begin{table}[htbp]
\centering
\small
\caption{System Roles: Linker versus Loader}
\label{tab:linker_vs_loader}
\begin{tabularx}{\textwidth}{l X X}
\toprule
\textbf{Dimension} & \textbf{Linker} & \textbf{Loader} \\
\midrule
\textbf{Primary Role} & Combines multiple compiled object files (\texttt{.o}) and static libraries (\texttt{.a}) into a single executable binary. & Reads the executable file from storage into physical RAM and prepares it for CPU execution. \\
\textbf{Symbol Resolution} & Resolves external function references and global variable symbols across compilation modules. & Binds dynamic addresses and binds shared libraries (\texttt{.so}, \texttt{.dll}). \\
\textbf{Relocation} & Assigns relative virtual addresses to code sections and data sections. & Adjusts memory addresses to physical memory base allocations; sets up stack and heap. \\
\textbf{Timing} & Invoked at build time. & Invoked at runtime whenever a program is launched. \\
\bottomrule
\end{tabularx}
\end{table}


\section{Control Unit Design: Hardwired versus Microprogrammed}
\label{sec:control_unit_design}

The \textbf{Control Unit (CU)} orchestrates the processor datapath by issuing synchronized control signals (register enables, multiplexer selects, ALU opcodes) during each clock cycle.

\begin{enumerate}[label=\bfseries\arabic*.]
    \item \textbf{Hardwired Control Unit:}
    \begin{itemize}
        \item Constructed using fixed combinational logic gates, decoders, and a central state register (Finite State Machine).
        \item \textit{Performance:} Maximum operating speed; control signals propagate within gate delays.
        \item \textit{Trade-off:} Inflexible; adding or modifying an instruction mandates a complete hardware silicon redesign.
        \item \textit{Adoption:} Standard in RISC processors (ARM, RISC-V, MIPS).
    \end{itemize}
    \item \textbf{Microprogrammed Control Unit:}
    \begin{itemize}
        \item Formulated by Maurice Wilkes in 1951. Each machine instruction initiates an internal execution sequence of low-level \textbf{Microinstructions} stored in an on-chip \textbf{Control Memory (ROM)}.
        \item \textit{Flexibility:} Extremely easy to modify or expand ISA instructions by updating firmware microcode without hardware layout changes.
        \item \textit{Trade-off:} Slower execution speed due to control memory read latency.
        \item \textit{Adoption:} Standard in complex CISC architectures (Intel x86, VAX, IBM System/360).
    \end{itemize}
\end{enumerate}

\begin{table}[htbp]
\centering
\small
\caption{Architectural Comparison: Hardwired versus Microprogrammed Control Units}
\label{tab:hardwired_vs_microprogrammed}
\begin{tabularx}{\textwidth}{l X X}
\toprule
\textbf{Feature} & \textbf{Hardwired Control Unit} & \textbf{Microprogrammed Control Unit} \\
\midrule
\textbf{Implementation} & Discrete combinational gates, decoders, flip-flops. & Control Store memory (ROM) holding micro-routines. \\
\textbf{Execution Speed} & Ultra-fast (single gate-delay propagation). & Slower (constrained by ROM memory access latency). \\
\textbf{Modification Flexibility} & Extremely difficult (requires chip redesign/re-spin). & Highly flexible (microcode patch updates). \\
\textbf{Instruction Complexity} & Best suited for simple, uniform RISC instruction sets. & Ideal for complex, multi-cycle CISC instruction sets. \\
\textbf{Control Word Layout} & Generated directly by combinational decoders. & Read sequentially from control store word fields. \\
\bottomrule
\end{tabularx}
\end{table}


\section{Memory Access Methods}
\label{sec:memory_access_methods}

Computer architectures employ four distinct physical mechanisms to access data elements from storage devices:
\begin{enumerate}[label=\bfseries\arabic*.]
    \item \textbf{Sequential Access:} Data is organized into linear records. To access record $N$, the hardware read mechanism must sequentially traverse records $1$ through $N-1$ physically. (Access time is highly variable and depends on distance; e.g., magnetic tape).
    \item \textbf{Direct Access:} Blocks possess unique addresses. Access is achieved by physically positioning a read head near the target block, followed by sequential scanning within the sector. (Access time is variable; e.g., hard disk drives, optical discs).
    \item \textbf{Random Access:} Each addressable location in memory possesses a dedicated, hardwired addressing mechanism. Any arbitrary physical location can be accessed in \textbf{constant time $O(1)$}, completely independent of physical location or access sequence. (e.g., semiconductor RAM and ROM).
    \item \textbf{Associative Access:} Data is retrieved based on its \textbf{content} rather than its memory address. All memory words are compared in parallel against a target search key simultaneously. (e.g., Cache Translation Lookaside Buffers [TLB], Content-Addressable Memory [CAM]).
\end{enumerate}


\section{Content-Addressable Memory (CAM and TCAM)}
\label{sec:cam_tcam}

\subsection{Operational Mechanism of CAM}
Standard RAM takes an address as input and returns the data stored at that address:
\[ \text{Address} \implies \text{RAM} \implies \text{Data} \]
In contrast, \textbf{Content-Addressable Memory (CAM)} takes a data search key as input, searches all internal memory locations in parallel in a single clock cycle, and returns the physical matching address:
\[ \text{Data Search Key} \implies \text{CAM} \implies \text{Matching Address / Hit Flag} \]

\subsection{Binary CAM versus Ternary CAM (TCAM)}
\begin{itemize}
    \item \textbf{Binary CAM:} Stores strictly binary bits ($0$ and $1$). Performs exact bitwise pattern matching.
    \item \textbf{Ternary CAM (TCAM):} Stores three logical states: $0$, $1$, and \textbf{Don't-Care ($X$)}.
    \begin{itemize}
        \item \textit{Critical Networking Application:} High-speed core Internet routers use TCAM to execute wire-speed \textbf{Longest Prefix Match (LPM)} packet routing table lookups and Access Control List (ACL) security filtering in a single clock cycle at hundreds of millions of packets per second.
    \end{itemize}
\end{itemize}


\section{Memory Bandwidth, Bus Arbitration, and Latency}
\label{sec:memory_bandwidth}

\subsection{Peak Theoretical Memory Bandwidth Formula}
Memory bandwidth measures the maximum rate of data transfer between the CPU and the memory subsystem:

\begin{equation}
\label{eq:memory_bandwidth}
\boxed{\text{Peak Bandwidth} = \text{Bus Width (Bytes)} \times \text{Clock Frequency } (f) \times \text{Transfers Per Clock Cycle}}
\end{equation}

\textbf{Worked Example:} Compute the peak memory bandwidth of a dual-channel DDR4-3200 system.
\begin{itemize}
    \item Each memory channel has a data bus width of 64 bits ($= 8\text{ bytes}$).
    \item Dual-channel configuration provides $2 \times 8 = 16\text{ bytes}$ total bus width.
    \item DDR4-3200 operates at an effective transfer rate of $3200\text{ MT/s}$ ($3.2 \times 10^9\text{ transfers/second}$).
\end{itemize}
\[ \text{Peak Bandwidth} = 16\text{ bytes} \times 3.2 \times 10^9\text{ transfers/s} = 51.2\text{ GB/s} \]

\subsection{Bus Arbitration Mechanisms}
When multiple bus masters (CPU, DMA Controller, GPU) request control of the shared system bus simultaneously, \textbf{Bus Arbitration} determines bus ownership:
\begin{enumerate}[label=\bfseries\arabic*.]
    \item \textbf{Daisy Chaining (Serial Centralized):} A single Bus Grant line propagates serially from device to device. Devices closer to the arbiter have highest priority. Vulnerable to single-point failure.
    \item \textbf{Centralized Parallel Arbitration:} Each bus master has independent Bus Request (\texttt{BREQ}) and Bus Grant (\texttt{BGRANT}) lines connected directly to a central hardware arbiter, which enforces fair round-robin or strict priority scheduling.
    \item \textbf{Distributed / Decentralized Arbitration:} Each master participates in peer-to-peer collision detection or priority resolution across shared lines (e.g., SCSI, CAN bus).
\end{enumerate}



\section{Master Numerical Examination Problem Bank with Worked Solutions}
\label{sec:numerical_problems}

This section delivers comprehensive worked numerical problems reflecting the exact difficulty, analytical patterns, and scoring models used in Assistant Programmer, System Analyst, and IT recruitment examinations.

\subsection{Problem 1: Multi-Base Conversion}
\textbf{Problem:} Convert $(11010110.101)_2$ into Decimal, Octal, and Hexadecimal.

\textbf{Solution:}
\begin{enumerate}[label=\bfseries\alph*.]
    \item \textbf{Conversion to Decimal:}
    \begin{align*}
    \text{Integer Part: } & 1\cdot 2^7 + 1\cdot 2^6 + 0\cdot 2^5 + 1\cdot 2^4 + 0\cdot 2^3 + 1\cdot 2^2 + 1\cdot 2^1 + 0\cdot 2^0 \\
    &= 128 + 64 + 0 + 16 + 0 + 4 + 2 + 0 = 214_{10} \\
    \text{Fractional Part: } & 1\cdot 2^{-1} + 0\cdot 2^{-2} + 1\cdot 2^{-3} = 0.5 + 0 + 0.125 = 0.625_{10} \\
    \textbf{Decimal Result: } & \mathbf{214.625_{10}}
    \end{align*}
    \item \textbf{Conversion to Octal (Group in 3s from binary point):}
    \begin{itemize}
        \item Integer: $\underline{011} \ \underline{010} \ \underline{110} \implies 3 \ 2 \ 6_8$
        \item Fraction: $\underline{101} \implies 5_8$
        \item \textbf{Octal Result:} $\mathbf{326.5_8}$
    \end{itemize}
    \item \textbf{Conversion to Hexadecimal (Group in 4s from binary point):}
    \begin{itemize}
        \item Integer: $\underline{1101} \ \underline{0110} \implies \text{D} \ 6_{16}$
        \item Fraction: $\underline{1010} \implies \text{A}_{16}$
        \item \textbf{Hexadecimal Result:} $\mathbf{D6.A_{16}}$
    \end{itemize}
\end{enumerate}

\subsection{Problem 2: Two's Complement Signed Arithmetic and Overflow Detection}
\textbf{Problem:} Perform 8-bit two's complement addition for:
\begin{enumerate}[label=\bfseries\alph*.]
    \item $(+75_{10}) + (+68_{10})$
    \item $(-75_{10}) + (-68_{10})$
    \item $(+75_{10}) + (-68_{10})$
\end{enumerate}
For each case, state the Carry ($C$) flag, Overflow ($V$) flag, and mathematical validity.

\textbf{Solution:}
Recall 8-bit signed two's complement range: $[-128, +127]$.
\begin{itemize}
    \item $+75_{10} = 01001011_2$
    \item $+68_{10} = 01000100_2$
    \item $-75_{10} = \sim(01001011_2) + 1 = 10110100_2 + 1 = 10110101_2$
    \item $-68_{10} = \sim(01000100_2) + 1 = 10111011_2 + 1 = 10111100_2$
\end{itemize}

\begin{enumerate}[label=\bfseries\alph*.]
    \item \textbf{Case A: $(+75) + (+68)$:}
    \begin{center}
    \begin{tabular}{r c}
      & \texttt{01001011} \quad ($+75$) \\
    $+$ & \texttt{01000100} \quad ($+68$) \\
    \hline
    (0) & \texttt{10001111} \quad (Signed interpretation: $-113_{10}$)
    \end{tabular}
    \end{center}
    \begin{itemize}
        \item Carry out of MSB: $C_8 = 0$. Carry into MSB: $C_7 = 1$.
        \item Overflow flag: $V = C_7 \oplus C_8 = 1 \oplus 0 = \mathbf{1}$ (\textbf{OVERFLOW DETECTED}).
        \item \textit{Diagnosis:} Two positive numbers produced a negative result because theoretical sum $+143 > +127$. Result is mathematically \textbf{INVALID}.
    \end{itemize}

    \item \textbf{Case B: $(-75) + (-68)$:}
    \begin{center}
    \begin{tabular}{r c}
      & \texttt{10110101} \quad ($-75$) \\
    $+$ & \texttt{10111100} \quad ($-68$) \\
    \hline
    (1) & \texttt{01110001} \quad (Signed interpretation: $+113_{10}$)
    \end{tabular}
    \end{center}
    \begin{itemize}
        \item Carry out of MSB: $C_8 = 1$. Carry into MSB: $C_7 = 0$.
        \item Overflow flag: $V = C_7 \oplus C_8 = 0 \oplus 1 = \mathbf{1}$ (\textbf{OVERFLOW DETECTED}).
        \item \textit{Diagnosis:} Two negative numbers produced a positive result because $-143 < -128$. Result is \textbf{INVALID}.
    \end{itemize}

    \item \textbf{Case C: $(+75) + (-68)$:}
    \begin{center}
    \begin{tabular}{r c}
      & \texttt{01001011} \quad ($+75$) \\
    $+$ & \texttt{10111100} \quad ($-68$) \\
    \hline
    (1) & \texttt{00000111} \quad (Signed interpretation: $+7_{10}$)
    \end{tabular}
    \end{center}
    \begin{itemize}
        \item Carry out of MSB: $C_8 = 1$. Carry into MSB: $C_7 = 1$.
        \item Overflow flag: $V = C_7 \oplus C_8 = 1 \oplus 1 = \mathbf{0}$ (\textbf{NO OVERFLOW}).
        \item \textit{Diagnosis:} Sum is $+7_{10}$, correctly within $[-128, +127]$. Result is mathematically \textbf{VALID}.
    \end{itemize}
\end{enumerate}

\subsection{Problem 3: K-Map Minimization with Don't-Care Conditions}
\textbf{Problem:} Minimize the following 4-variable Boolean function in Sum of Products (SOP) form:
\[ F(A, B, C, D) = \sum m(0, 2, 5, 7, 8, 10, 14, 15) + d(3, 11) \]

\textbf{Solution:}
Populate the 16-cell K-map:
\begin{center}
\small
\begin{tabular}{|c|c|c|c|c|}
\hline
$AB \backslash CD$ & \textbf{00} & \textbf{01} & \textbf{11} & \textbf{10} \\
\hline
\textbf{00} & \textbf{1} ($m_0$) & 0 & $X$ ($d_3$) & \textbf{1} ($m_2$) \\
\hline
\textbf{01} & 0 & \textbf{1} ($m_5$) & \textbf{1} ($m_7$) & 0 \\
\hline
\textbf{11} & 0 & 0 & \textbf{1} ($m_{15}$) & \textbf{1} ($m_{14}$) \\
\hline
\textbf{10} & \textbf{1} ($m_8$) & 0 & $X$ ($d_{11}$) & \textbf{1} ($m_{10}$) \\
\hline
\end{tabular}
\end{center}

\begin{enumerate}[label=\bfseries\arabic*.]
    \item \textbf{Corner Quad ($m_0, m_2, m_8, m_{10}$):}
    The four extreme outer corners wrap around horizontally and vertically:
    Rows $\{00, 10\} \implies B=0$ is constant. Columns $\{00, 10\} \implies D=0$ is constant.
    Yields the prime implicant: $\mathbf{\overline{B}\,\overline{D}}$.
    \item \textbf{Quad utilizing Don't-Cares ($m_2, d_3, m_{10}, d_{11}$):}
    Alternatively, examine cells $m_5$ and $m_7$ (row $01$). Pair with $m_7$ and $m_{15}$?
    Notice cells in column $CD=11$: $d_3, m_7, m_{15}, d_{11}$ form a vertical column quad!
    In this quad, $C=1, D=1$ are constant, covering $m_7$ and $m_{15}$ cleanly.
    Yields: $\mathbf{CD}$.
    \item \textbf{Remaining $1$s: $m_5$ and $m_{14}$:}
    \begin{itemize}
        \item For $m_5$ ($0101$): Can be paired with $m_7$ ($0111$) in row $01$.
        In row $01$, $A=0, B=1$ are constant. In columns $01$ and $11$, $D=1$ is constant.
        Yields: $\mathbf{\overline{A}BD}$.
        \item For $m_{14}$ ($1110$): Can be paired with $m_{15}$ ($1111$) in row $11$.
        In row $11$, $A=1, B=1$ are constant. In columns $11$ and $10$, $C=1$ is constant.
        Yields: $\mathbf{ABC}$.
    \end{itemize}
    \item \textbf{Minimized Boolean Equation:}
    \[ \mathbf{F = \overline{B}\,\overline{D} + CD + \overline{A}BD + ABC} \]
\end{enumerate}

\subsection{Problem 4: Ripple Carry versus Carry Lookahead Adder Propagation Delay}
\textbf{Problem:} A 32-bit adder is constructed using logic gates with the following propagation delays:
\begin{itemize}
    \item XOR gate delay: $t_{\text{XOR}} = 1.2\text{ ns}$
    \item AND gate delay: $t_{\text{AND}} = 0.8\text{ ns}$
    \item OR gate delay: $t_{\text{OR}} = 0.8\text{ ns}$
\end{itemize}
Calculate:
\begin{enumerate}[label=\bfseries\alph*.]
    \item The worst-case propagation delay of a 32-bit Ripple Carry Adder (RCA).
    \item The worst-case propagation delay of a 32-bit Carry Lookahead Adder (CLA) constructed using 4-bit CLA blocks with a secondary block-lookahead carry unit.
\end{enumerate}

\textbf{Solution:}
\begin{enumerate}[label=\bfseries\alph*.]
    \item \textbf{Ripple Carry Adder (RCA):}
    In each Full Adder, carry generation requires $t_{\text{carry}} = t_{\text{AND}} + t_{\text{OR}} = 0.8 + 0.8 = 1.6\text{ ns}$.
    The carry-in must ripple sequentially through all 32 stages:
    \[ T_{\text{RCA}} = 32 \times t_{\text{carry}} = 32 \times 1.6\text{ ns} = \mathbf{51.2\text{ ns}} \]
    (Maximum possible clock frequency is restricted to $f_{\text{max}} = 1 / 51.2\text{ ns} \approx 19.5\text{ MHz}$).
    \item \textbf{Carry Lookahead Adder (CLA):}
    All generate ($G_i$) and propagate ($P_i$) signals are evaluated in parallel: $t_{P,G} = \max(t_{\text{AND}}, t_{\text{XOR}}) = 1.2\text{ ns}$.
    The lookahead carry logic evaluates all carries across a two-level AND-OR tree:
    \[ t_{\text{tree}} = t_{\text{AND}} + t_{\text{OR}} = 0.8 + 0.8 = 1.6\text{ ns} \]
    Sum bits are evaluated at the final stage via XOR: $t_{\text{sum}} = t_{\text{XOR}} = 1.2\text{ ns}$.
    \[ T_{\text{CLA}} = t_{P,G} + (2 \times t_{\text{tree}}) + t_{\text{sum}} = 1.2 + (2 \times 1.6) + 1.2 = \mathbf{5.6\text{ ns}} \]
    \textbf{Speedup:} The CLA is $\frac{51.2}{5.6} \approx \mathbf{9.14\times\text{ faster}}$ than the RCA.
\end{enumerate}

\subsection{Problem 5: Cache Memory Address Decomposition}
\textbf{Problem:} A computer system has a 32-bit byte-addressable physical memory space and a \textbf{256 KiB}, \textbf{8-way set-associative cache} with a \textbf{64-byte block (line) size}.
Determine:
\begin{enumerate}[label=\bfseries\alph*.]
    \item The number of Byte Offset bits.
    \item The number of Set Index bits.
    \item The number of Tag bits.
    \item The total size of the cache memory in bits, including Valid bits and Dirty bits.
\end{enumerate}

\textbf{Solution:}
\begin{enumerate}[label=\bfseries\alph*.]
    \item \textbf{Byte Offset ($b$):}
    \[ b = \log_2(\text{Block Size}) = \log_2(64) = \mathbf{6\text{ bits}} \]
    \item \textbf{Number of Sets ($S$) and Set Index ($s$):}
    \[ \text{Total Cache Lines} = \frac{\text{Total Capacity}}{\text{Line Size}} = \frac{256\text{ KiB}}{64\text{ B}} = \frac{256 \times 1024}{64} = 4096\text{ lines} \]
    \[ \text{Number of Sets } (S) = \frac{\text{Total Lines}}{\text{Associativity } N} = \frac{4096}{8} = 512\text{ sets} = 2^9\text{ sets} \]
    \[ s = \log_2(S) = \log_2(512) = \mathbf{9\text{ bits}} \]
    \item \textbf{Tag Bits ($t$):}
    \[ t = \text{Address Bits} - s - b = 32 - 9 - 6 = \mathbf{17\text{ bits}} \]
    \item \textbf{Total Cache Storage Overhead in Bits:}
    Each line stores:
    \begin{itemize}
        \item Data: $64\text{ bytes} \times 8\text{ bits/byte} = 512\text{ bits}$.
        \item Tag: $17\text{ bits}$.
        \item Valid bit: $1\text{ bit}$.
        \item Dirty bit: $1\text{ bit}$.
        \item Total bits per line $= 512 + 17 + 1 + 1 = 531\text{ bits}$.
    \end{itemize}
    Total cache size $= 4096\text{ lines} \times 531\text{ bits/line} = 2,174,976\text{ bits} = \mathbf{265.5\text{ KiB}}$.
\end{enumerate}

\subsection{Problem 6: Two-Level Cache AMAT Calculation}
\textbf{Problem:} A high-performance processor features a split L1 cache and a unified L2 cache with the following parameters:
\begin{itemize}
    \item L1 instruction miss rate: $2\%$, L1 data miss rate: $8\%$.
    \item Instruction references constitute $70\%$ of all accesses; data references constitute $30\%$.
    \item L1 hit latency: $1\text{ ns}$ (1 clock cycle).
    \item L2 unified local miss rate: $12\%$.
    \item L2 hit latency: $6\text{ ns}$.
    \item Main memory access penalty: $90\text{ ns}$.
\end{itemize}
Calculate the Average Memory Access Time (AMAT).

\textbf{Solution:}
\begin{enumerate}[label=\bfseries\arabic*.]
    \item \textbf{Combined Average L1 Miss Rate ($m_{\text{L1}}$):}
    \[ m_{\text{L1}} = (0.70 \times 0.02) + (0.30 \times 0.08) = 0.014 + 0.024 = 0.038 = 3.8\% \]
    \item \textbf{L2 Miss Penalty:}
    \[ \text{Penalty}_{\text{L2}} = T_{\text{hit\_L2}} + (m_{\text{L2\_local}} \times T_{\text{memory}}) = 6\text{ ns} + (0.12 \times 90\text{ ns}) = 6 + 10.8 = 16.8\text{ ns} \]
    \item \textbf{Overall System AMAT:}
    \[ \text{AMAT} = T_{\text{hit\_L1}} + (m_{\text{L1}} \times \text{Penalty}_{\text{L2}}) = 1\text{ ns} + (0.038 \times 16.8\text{ ns}) = 1 + 0.6384 = \mathbf{1.6384\text{ ns}} \]
\end{enumerate}

\subsection{Problem 7: CPU Execution Time and CPI Optimization}
\textbf{Problem:} A benchmark program of $2 \times 10^9$ instructions executes on a $3.2\text{ GHz}$ processor. The instruction distribution is:
\begin{itemize}
    \item ALU operations: $40\%$, requiring $1\text{ cycle}$ each.
    \item Load operations: $25\%$, requiring $4\text{ cycles}$ each.
    \item Store operations: $15\%$, requiring $2\text{ cycles}$ each.
    \item Branch operations: $20\%$, requiring $3\text{ cycles}$ each.
\end{itemize}
Determine:
\begin{enumerate}[label=\bfseries\alph*.]
    \item The weighted average CPI.
    \item Total CPU execution time in seconds.
    \item The MIPS rating of the processor on this workload.
\end{enumerate}

\textbf{Solution:}
\begin{enumerate}[label=\bfseries\alph*.]
    \item \textbf{Weighted Average CPI:}
    \[ \text{CPI}_{\text{avg}} = (0.40 \times 1) + (0.25 \times 4) + (0.15 \times 2) + (0.20 \times 3) = 0.40 + 1.00 + 0.30 + 0.60 = \mathbf{2.30} \]
    \item \textbf{Total CPU Execution Time:}
    \[ T_{\text{CPU}} = \frac{\text{IC} \times \text{CPI}}{f} = \frac{(2 \times 10^9) \times 2.30}{3.2 \times 10^9\text{ Hz}} = \frac{4.60 \times 10^9}{3.2 \times 10^9} = \mathbf{1.4375\text{ seconds}} \]
    \item \textbf{MIPS Rating:}
    \[ \text{MIPS} = \frac{f}{\text{CPI} \times 10^6} = \frac{3200\text{ MHz}}{2.30} \approx \mathbf{1391.3\text{ MIPS}} \]
\end{enumerate}

\subsection{Problem 8: Amdahl's Law Multicore Scaling}
\textbf{Problem:} An image processing application spends $85\%$ of its total execution time inside an image filtering routine that can be parallelized across multiple CPU cores, while the remaining $15\%$ consists of sequential file I/O that cannot be parallelized.
\begin{enumerate}[label=\bfseries\alph*.]
    \item What is the speedup achieved by executing the application on an 8-core processor?
    \item What is the maximum theoretical speedup achievable even with an infinite number of CPU cores?
\end{enumerate}

\textbf{Solution:}
Here, parallel fraction $f = 0.85$, serial fraction $(1 - f) = 0.15$.
\begin{enumerate}[label=\bfseries\alph*.]
    \item \textbf{Speedup on an 8-Core Processor ($s = 8$):}
    \[ S = \frac{1}{(1 - f) + \dfrac{f}{s}} = \frac{1}{0.15 + \dfrac{0.85}{8}} = \frac{1}{0.15 + 0.10625} = \frac{1}{0.25625} \approx \mathbf{3.902\times} \]
    (An 8-core CPU delivers less than a $4\times$ real-world speedup due to serial overhead).
    \item \textbf{Maximum Theoretical Speedup ($s \to \infty$):}
    \[ S_{\text{max}} = \frac{1}{1 - f} = \frac{1}{0.15} = \mathbf{6.67\times} \]
\end{enumerate}

\subsection{Problem 9: Intel 8086 Physical Address Calculation}
\textbf{Problem:} An 8086 microprocessor executes instructions with the following register contents:
\begin{itemize}
    \item $\text{CS} = 4\text{B}20_{16}, \quad \text{IP} = 1\text{C}34_{16}$
    \item $\text{SS} = 7000_{16}, \quad \text{SP} = 2\text{FF}0_{16}$
\end{itemize}
Calculate the 20-bit physical addresses generated for:
\begin{enumerate}[label=\bfseries\alph*.]
    \item The next instruction fetch.
    \item The current top of the stack.
\end{enumerate}

\textbf{Solution:}
Formula: $\text{Physical Address} = (\text{Segment} \times 16_{10}) + \text{Offset} = (\text{Segment} \ll 4) + \text{Offset}$.
\begin{enumerate}[label=\bfseries\alph*.]
    \item \textbf{Instruction Fetch Address:}
    \begin{align*}
    \text{CS} \times 10_{16} &= 4\text{B}200_{16} \\
    \text{IP} &= \phantom{0}1\text{C}34_{16} \\
    \text{Physical Address} &= 4\text{B}200_{16} + 1\text{C}34_{16} = \mathbf{4\text{CE}34_{16}}
    \end{align*}
    \item \textbf{Top of Stack Address:}
    \begin{align*}
    \text{SS} \times 10_{16} &= 70000_{16} \\
    \text{SP} &= \phantom{0}2\text{FF}0_{16} \\
    \text{Physical Address} &= 70000_{16} + 2\text{FF}0_{16} = \mathbf{72\text{FF}0_{16}}
    \end{align*}
\end{enumerate}

\subsection{Problem 10: ADC Resolution and Sampling Rate}
\textbf{Problem:} A 12-bit ADC has an analog input reference range of $0\text{ V}$ to $+5.0\text{ V}$.
\begin{enumerate}[label=\bfseries\alph*.]
    \item Calculate the step size (voltage resolution) in millivolts.
    \item If the analog input voltage is $V_{\text{in}} = 3.25\text{ V}$, determine the generated 12-bit binary digital output code.
    \item If the input signal contains frequency components up to $15\text{ kHz}$, calculate the minimum Nyquist sampling rate and the sampling interval $T_s$.
\end{enumerate}

\textbf{Solution:}
\begin{enumerate}[label=\bfseries\alph*.]
    \item \textbf{Voltage Step Size ($\Delta V$):}
    \[ \Delta V = \frac{V_{\text{ref}}}{2^n - 1} = \frac{5.0\text{ V}}{2^{12} - 1} = \frac{5.0}{4095} \approx 0.001221\text{ V} = \mathbf{1.221\text{ mV}} \]
    \item \textbf{Digital Output Code ($D$):}
    \[ D = \left\lfloor \frac{V_{\text{in}}}{\Delta V} \right\rfloor = \left\lfloor \frac{3.25\text{ V}}{0.001221\text{ V}} \right\rfloor = 2661_{10} \]
    Converting $2661_{10}$ to 12-bit binary:
    \begin{itemize}
        \item $2661 / 256 = 10 \implies \text{Hex } \text{A}$
        \item $2661 \pmod{256} = 101 \implies 101 / 16 = 6, \text{ Remainder } 5$
        \item Hexadecimal: $\text{A}65_{16}$
        \item Binary: $\mathbf{1010 \ 0110 \ 0101_2}$
    \end{itemize}
    \item \textbf{Nyquist Sampling Rate and Interval:}
    \[ f_{\text{Nyquist}} = 2 \times f_{\text{max}} = 2 \times 15\text{ kHz} = \mathbf{30\text{ kHz} = 30,000\text{ samples/second}} \]
    \[ T_s = \frac{1}{f_{\text{Nyquist}}} = \frac{1}{30 \times 10^3\text{ s}^{-1}} \approx \mathbf{33.33\ \mu\text{s}} \]
\end{enumerate}


\section{Theoretical Viva and Recruitment Examination Question Bank}
\label{sec:viva_questions}

\begin{enumerate}[label=\bfseries Q\arabic*.]
    \item \textbf{What is the primary difference between Computer Architecture and Computer Organization?} \\
    \textbf{Answer:} Computer Architecture defines the programmer-visible attributes and functional specifications (ISA, addressing modes, register widths, data types). Computer Organization (microarchitecture) describes how those specifications are physically implemented in hardware (datapath, control units, pipeline stages, cache structures). Architecture answers \textit{``What does the CPU do?''}; Organization answers \textit{``How is it built?''}

    \item \textbf{Why does modern digital hardware utilize two's complement representation rather than one's complement or sign-magnitude?} \\
    \textbf{Answer:} Two's complement possesses two decisive engineering advantages: (1) It provides a unique, unambiguous representation for zero ($00000000_2$), avoiding the dual $+0/-0$ anomalies of other formats. (2) Binary subtraction ($A - B$) is executed natively using standard binary addition circuitry ($A + \sim B + 1$), eliminating the need for separate subtraction hardware or end-around carry logic.

    \item \textbf{How is arithmetic overflow detected in two's complement addition hardware?} \\
    \textbf{Answer:} In an $n$-bit parallel adder, overflow is detected by computing the Exclusive-OR of the carry-in and carry-out of the Most Significant Bit (sign bit): $V = C_{n-1} \oplus C_n$. If $V = 1$, an overflow has occurred. Conceptually, overflow occurs when two numbers of like sign produce a sum of opposite sign.

    \item \textbf{What is a Universal Logic Gate, and why are NAND and NOR classified as universal?} \\
    \textbf{Answer:} A universal gate is a logic gate capable of implementing any arbitrary Boolean switching function without requiring any other gate type. NAND and NOR are universal because each can synthesize the complete foundational set of Boolean primitives: NOT, AND, and OR.

    \item \textbf{State De Morgan's First and Second Theorems.} \\
    \textbf{Answer:}
    First Theorem: $\overline{A + B} = \overline{A} \cdot \overline{B}$ (The complement of a logical sum equals the product of the individual complements). \\
    Second Theorem: $\overline{A \cdot B} = \overline{A} + \overline{B}$ (The complement of a logical product equals the sum of the individual complements).

    \item \textbf{Why is Gray Code strictly required on Karnaugh Map axes instead of standard binary?} \\
    \textbf{Answer:} Gray code enforces the \textbf{unit-distance property}, ensuring that adjacent physical cells differ in exactly one variable literal. This enables visual minimization via the Boolean identity $A + \overline{A} = 1$, physically eliminating varying literals from groupings.

    \item \textbf{What is the difference between a Latch and a Flip-Flop?} \\
    \textbf{Answer:} A latch is \textbf{level-sensitive} (transparent while its enable line is high, updating outputs continuously). A flip-flop is \textbf{edge-triggered} (samples inputs and updates outputs only during the instantaneous rising or falling transition of a clock signal).

    \item \textbf{Explain the Race-Around condition in JK flip-flops and how it is eliminated.} \\
    \textbf{Answer:} In a level-triggered JK flip-flop with $J=1, K=1$, if the clock pulse width ($t_p$) exceeds the internal gate propagation delay ($t_{pd}$), the output will toggle continuously back and forth throughout the clock pulse, leaving the final resting state indeterminate. It is eliminated by using \textbf{Master-Slave JK flip-flops} or \textbf{edge-triggered} circuit topologies.

    \item \textbf{Define Setup Time and Hold Time.} \\
    \textbf{Answer:} Setup time ($t_{\text{setup}}$) is the minimum duration the data input must remain stable \textit{before} the active clock transition. Hold time ($t_{\text{hold}}$) is the minimum duration the data input must remain stable \textit{after} the active clock transition. Violating either causes \textbf{metastability}.

    \item \textbf{What is the Von Neumann Bottleneck, and how do modern architectures mitigate it?} \\
    \textbf{Answer:} The throughput ceiling caused by the shared bus between the CPU and memory, where memory transfer bandwidth lags far behind CPU execution speeds. Mitigations include multi-level on-chip caching (L1/L2/L3), split Harvard L1 caches (L1-I and L1-D), aggressive hardware prefetching, and out-of-order execution.

    \item \textbf{Differentiate between Moore and Mealy Finite State Machines.} \\
    \textbf{Answer:} In a Moore machine, outputs depend \textit{exclusively on the present internal state}. In a Mealy machine, outputs depend on \textit{both the present internal state and current external inputs}. Mealy machines generally require fewer states, while Moore machines exhibit cleaner glitch-free timing closure.

    \item \textbf{What are the three classes of Pipeline Hazards, and how is each resolved?} \\
    \textbf{Answer:}
    (1) \textbf{Structural Hazards} (hardware resource contention): Resolved via separate L1 instruction/data caches and multi-ported register files. \\
    (2) \textbf{Data Hazards} (data dependencies: RAW, WAR, WAW): Resolved via hardware forwarding/bypassing and compiler instruction scheduling. \\
    (3) \textbf{Control Hazards} (branch instruction uncertainty): Resolved via dynamic branch prediction (2-bit saturating counters) and branch target buffers (BTBs).

    \item \textbf{Explain the difference between RISC and CISC architectures.} \\
    \textbf{Answer:} RISC emphasizes simple, fixed-length 32-bit instructions executing in $\approx 1$ cycle using a load/store architecture and hardwired control. CISC emphasizes variable-length, multi-cycle instructions with memory-to-memory operations and microprogrammed control memory to minimize code size.

    \item \textbf{How does the Intel 8086 compute a 20-bit physical address from 16-bit registers?} \\
    \textbf{Answer:} Memory is partitioned into logical segments. The 8086 shifts the 16-bit Segment Base left by 4 bits (multiplying by 16) and adds the 16-bit Offset: $\text{Physical Address} = (\text{Segment} \ll 4) + \text{Offset}$.

    \item \textbf{Explain the 3 Cs model of cache misses.} \\
    \textbf{Answer:}
    \textbf{Compulsory (Cold):} First-time access to a block; unavoidable. \\
    \textbf{Capacity:} Working data set exceeds cache size; blocks evicted and re-fetched. \\
    \textbf{Conflict:} Multiple memory blocks map to the identical cache set in direct or set-associative caches. \\
    (The 4th C in multiprocessors is \textbf{Coherence misses} caused by peer core invalidations).

    \item \textbf{What is the MESI cache coherence protocol?} \\
    \textbf{Answer:} A 4-state protocol governing multi-core caches: \textbf{Modified} (exclusively owned, dirty relative to DRAM), \textbf{Exclusive} (exclusively owned, clean), \textbf{Shared} (replicated in peer caches, clean), and \textbf{Invalid} (stale data).

    \item \textbf{Compare SRAM and DRAM at the silicon circuit level.} \\
    \textbf{Answer:} SRAM uses a 6-transistor (6T) bistable CMOS latch; it is static, ultra-fast ($<2\text{ ns}$), low density, and used for CPU caches. DRAM uses a 1-transistor 1-capacitor (1T-1C) cell; it is dynamic, requires periodic refresh ($64\text{ ms}$) due to charge leakage, possesses extreme packing density, and is used for main memory.

    \item \textbf{What is the Nyquist Sampling Theorem, and what occurs if it is violated?} \\
    \textbf{Answer:} The theorem mandates that sampling frequency must satisfy $f_s \ge 2 f_{\text{max}}$. Violating it causes \textbf{aliasing}, where high-frequency signals fold into lower baseband frequencies, permanently corrupting reconstructed data.

    \item \textbf{What is Endianness? Differentiate Big-Endian and Little-Endian.} \\
    \textbf{Answer:} Endianness is the byte ordering convention for multi-byte values in memory. Big-Endian stores the Most Significant Byte at the lowest memory address (network byte order). Little-Endian stores the Least Significant Byte at the lowest address (x86-64, ARM default).

    \item \textbf{Differentiate between Memory-Mapped I/O and Isolated (Port) I/O.} \\
    \textbf{Answer:} Memory-Mapped I/O shares the unified memory address space and uses standard memory instructions (\texttt{MOV}, \texttt{LOAD}) to access peripheral registers. Isolated I/O uses a distinct address space and requires dedicated machine instructions (\texttt{IN}, \texttt{OUT}) and control lines (\texttt{IOR}, \texttt{IOW}).
\end{enumerate}


\section{Master Comparison Tables of Digital Computing}
\label{sec:master_comparison_tables}

\begin{table}[htbp]
\centering
\small
\caption{Comprehensive Synthesis: Core Computer Architecture Comparison Matrix}
\label{tab:master_architecture_matrix}
\begin{tabularx}{\textwidth}{l X X}
\toprule
\textbf{Architectural Domain} & \textbf{Paradigm A} & \textbf{Paradigm B} \\
\midrule
\textbf{System Memory Model} & \textbf{Von Neumann:} Unified memory and shared bus for instructions and data. Simple, lower cost, vulnerable to bottleneck. & \textbf{Harvard:} Separate memory and independent buses for instructions and data. High concurrent throughput, higher silicon area. \\
\midrule
\textbf{Logic Circuit Class} & \textbf{Combinational:} Output strictly depends on current inputs. Zero memory, acyclic, no clock required. (Adders, MUX). & \textbf{Sequential:} Output depends on current inputs and stored internal state. Contains feedback and flip-flops. (Counters, FSM). \\
\midrule
\textbf{Storage Cell Triggering} & \textbf{Latch:} Level-sensitive. Output follows input continuously while enabled (transparent). Vulnerable to race-around. & \textbf{Flip-Flop:} Edge-triggered. Output updates strictly at clock voltage transitions. Immune to race-around. \\
\midrule
\textbf{State Machine Model} & \textbf{Moore Machine:} Output depends exclusively on present state ($Y=g(S)$). Glitch-free clock-synchronized transitions. & \textbf{Mealy Machine:} Output depends on present state and current inputs ($Y=g(S,X)$). Faster response, fewer states, glitch-prone. \\
\midrule
\textbf{Instruction Set Style} & \textbf{RISC:} Fixed 32-bit instructions, single-cycle execution, load/store architecture, hardwired control. (ARM, RISC-V). & \textbf{CISC:} Variable 1--15 byte instructions, multi-cycle execution, memory operands, microcoded control store. (x86-64). \\
\midrule
\textbf{Microprocessor Family} & \textbf{Intel 8085:} 8-bit data bus, 16-bit address (64 KB), flat linear memory, unified core, no instruction prefetch queue. & \textbf{Intel 8086:} 16-bit data bus, 20-bit address (1 MB), segmented memory, dual BIU/EU with 6-byte instruction prefetch queue. \\
\midrule
\textbf{Semiconductor RAM} & \textbf{SRAM:} 6-transistor latch, static storage, ultra-fast ($<2\text{ ns}$), high power, low density. Used for on-chip cache. & \textbf{DRAM:} 1-transistor 1-capacitor, dynamic charge, periodic refresh required, high density, low cost. Used for main memory. \\
\midrule
\textbf{Control Unit Design} & \textbf{Hardwired Control:} Combinational gates and state registers. Maximum speed ($O(1)$ gate delay), rigid, inflexible. & \textbf{Microprogrammed Control:} Microinstructions in control ROM. Flexible (firmware updates), slower due to ROM read latency. \\
\bottomrule
\end{tabularx}
\end{table}


\section{Formula Quick-Revision Sheet and 35 Examination Traps}
\label{sec:formula_cheat_sheet}

\subsection{Essential Formulas Cheat Sheet}
\begin{enumerate}[label=\bfseries\arabic*.]
    \item \textbf{Clock Speed and Period:} $T = \frac{1}{f} \quad \Longleftrightarrow \quad f = \frac{1}{T}$
    \item \textbf{Addressable Memory Capacity:} $\text{Capacity} = 2^k \times W$ (where $k = \text{address lines}$, $W = \text{word width}$).
    \item \textbf{Two's Complement Value:} $\text{Value} = -b_{n-1} 2^{n-1} + \sum_{i=0}^{n-2} b_i 2^i$
    \item \textbf{Signed Integer Range ($n$ bits):} $[-2^{n-1}, +(2^{n-1} - 1)]$
    \item \textbf{Arithmetic Overflow Condition:} $V = C_{n-1} \oplus C_n$
    \item \textbf{Full Adder Logic Equations:}
    \[ S = A \oplus B \oplus C_{\text{in}}, \qquad C_{\text{out}} = AB + C_{\text{in}}(A \oplus B) \]
    \item \textbf{Carry Lookahead Generator:} $G_i = A_i B_i, \quad P_i = A_i \oplus B_i, \quad C_{i+1} = G_i + P_i C_i$
    \item \textbf{Flip-Flop Characteristic Equations:}
    \begin{align*}
    \text{SR: } & Q^+ = S + \overline{R}Q \quad (S \cdot R = 0) & \text{D: } & Q^+ = D \\
    \text{JK: } & Q^+ = J\overline{Q} + \overline{K}Q & \text{T: } & Q^+ = T \oplus Q
    \end{align*}
    \item \textbf{Maximum Operating Clock Frequency:} $f_{\text{max}} = \frac{1}{t_{cq} + t_{\text{comb}} + t_{\text{setup}}}$
    \item \textbf{Pipeline Speedup ($k$ stages, $n$ instructions):} $S = \frac{n \cdot k}{k + n - 1} \xrightarrow{n \to \infty} k$
    \item \textbf{Intel 8086 Physical Address:} $\text{Physical Address} = (\text{Segment} \ll 4) + \text{Offset} = (\text{Segment} \times 16) + \text{Offset}$
    \item \textbf{Average Memory Access Time (AMAT):} $\text{AMAT} = \text{Hit Time} + (\text{Miss Rate} \times \text{Miss Penalty})$
    \item \textbf{Cache Address Partitioning:}
    \[ \text{Offset } b = \log_2(\text{Line Size}), \quad \text{Set Index } s = \log_2(\text{Total Sets}), \quad \text{Tag } t = W - s - b \]
    \item \textbf{The Iron Law of CPU Time:} $T_{\text{CPU}} = \frac{\text{Instruction Count} \times \text{CPI}}{f_{\text{clock}}}$
    \item \textbf{Amdahl's Law Speedup:} $S = \frac{1}{(1 - f) + \dfrac{f}{s}} \xrightarrow{s \to \infty} \frac{1}{1 - f}$
    \item \textbf{Nyquist Sampling Rate:} $f_{\text{Nyquist}} \ge 2 \cdot f_{\text{max}}$
    \item \textbf{ADC Step Size:} $\Delta V = \frac{V_{\text{ref}}}{2^n - 1}$
    \item \textbf{Peak Memory Bandwidth:} $\text{Bandwidth} = \text{Bus Width (Bytes)} \times f_{\text{bus}} \times \text{Transfers/Cycle}$
\end{enumerate}

\subsection{35 Critical Examination Traps and Pitfalls}
\begin{enumerate}[label=\bfseries\arabic*.]
    \item \textbf{Confusing Carry Flag ($C$) with Overflow Flag ($V$):} Carry indicates unsigned arithmetic wraparound beyond $2^n - 1$; Overflow indicates signed two's complement sign-bit corruption ($V = C_{\text{in}} \oplus C_{\text{out}}$ of MSB).
    \item \textbf{Sign-Extension Errors:} When extending an 8-bit negative number to 16 bits, replicate the sign bit (MSB), not zeros. For example, $-5 = \texttt{0xFB} \to \texttt{0xFFFB}$, not $\texttt{0x00FB}$.
    \item \textbf{Two's Complement Asymmetric Range:} For 8 bits, the range is $-128$ to $+127$. Negating $-128$ produces $+128$, which cannot be represented and triggers an immediate overflow.
    \item \textbf{K-Map Gray Code Axis Reversal:} Writing axes as $00, 01, 10, 11$ instead of $00, 01, 11, 10$ invalidates cell adjacency.
    \item \textbf{Group Sizes Must Be Powers of Two:} Grouping 3, 5, or 6 cells on a K-map is invalid; only $1, 2, 4, 8, 16$ are permitted.
    \item \textbf{Forgetting Toroidal Wrap-Around:} Failing to group the 4 outer corners ($m_0, m_2, m_8, m_{10}$) into a single quad ($\overline{B}\,\overline{D}$).
    \item \textbf{Pure Don't-Care Groups:} Forming a group consisting exclusively of $X$s is strictly redundant and incorrect.
    \item \textbf{Full Adder NAND Count:} Synthesizing a Full Adder using 2-input NAND gates requires \textbf{9 NAND gates}, not 5 (Half Adder requires 5).
    \item \textbf{Carry Lookahead Gate Delays:} CLA generates all carry signals in \textbf{2 gate delays} (one AND level, one OR level), independent of word width.
    \item \textbf{Priority Encoder Inactive State:} An active $D_0$ cannot be distinguished from all-zero idle inputs without checking the \textbf{Valid ($V$)} output bit.
    \item \textbf{Latch vs Flip-Flop Sensitivity:} Latches are level-sensitive; Flip-Flops are edge-triggered.
    \item \textbf{Race-Around Condition Prerequisite:} Race-around occurs in level-triggered JK flip-flops only when $J=1, K=1$ and clock pulse width $t_p > t_{pd}$.
    \item \textbf{Master-Slave Clock Phasing:} The Master latch is enabled when $\text{CLK}=1$; the Slave latch updates outputs when $\text{CLK}=0$.
    \item \textbf{Counter Modulus Formula:} $n$ flip-flops can realize a maximum modulus of $\text{MOD} = 2^n$. To count up to $N$, the required flip-flop count is $n = \lceil \log_2 N \rceil$.
    \item \textbf{Ring vs Johnson Counter States:} An $n$-flip-flop Ring Counter provides $n$ states; an $n$-flip-flop Johnson Counter provides $2n$ states.
    \item \textbf{Moore vs Mealy Output Timing:} Moore machine outputs change synchronously with clock edges; Mealy machine outputs change asynchronously as soon as inputs transition.
    \item \textbf{Von Neumann vs Harvard Unified Memory:} Von Neumann has a single unified bus for instructions and data; Harvard has separate independent buses.
    \item \textbf{8086 Physical Address Multiplier:} The segment base must be shifted left by 4 bits (multiplied by $16_{10}$ or $10_{16}$), not 8 bits or 2 bits.
    \item \textbf{8086 Bus vs Execution Unit:} BIU handles instruction prefetching (6-byte queue) and address generation; EU handles instruction decoding and execution.
    \item \textbf{Load/Store Architecture Constraint:} In pure RISC ISAs, ALU instructions \textbf{never} accept memory operands; all arithmetic occurs strictly in registers.
    \item \textbf{Memory-Mapped vs Port-Mapped Address Space:} Memory-Mapped I/O reduces available RAM address space; Port-Mapped I/O preserves the full RAM address space.
    \item \textbf{Pipeline Latency vs Throughput:} Pipelining improves instruction \textbf{throughput} (rate of completion), but does \textbf{not} decrease individual instruction execution \textbf{latency}.
    \item \textbf{Pipeline Theoretical Speedup Bound:} The maximum theoretical speedup of a $k$-stage pipeline is strictly bounded by $k$.
    \item \textbf{RAW Hazard Frequency:} RAW (Read-After-Write) is the only data hazard that can occur in a simple, in-order execution pipeline.
    \item \textbf{Load-Use Hazard Bubble:} Even with full hardware data forwarding, a load instruction immediately followed by an instruction consuming that data requires \textbf{one stall cycle}.
    \item \textbf{Cache Line Index vs Total Lines:} In an $N$-way set-associative cache, the index selects the \textbf{Set}, not the individual line ($S = \text{Total Lines} / N$).
    \item \textbf{Direct-Mapped Associativity:} A Direct-Mapped cache is mathematically a \textbf{1-Way Set-Associative cache}.
    \item \textbf{Fully Associative Index Bits:} A Fully Associative cache has \textbf{zero Set Index bits}; the entire address (minus byte offset) is the Tag.
    \item \textbf{Write-Back Dirty Bit:} Write-Back caches must maintain a \textbf{Dirty Bit} per line; Write-Through caches do not require a dirty bit.
    \item \textbf{SRAM vs DRAM Volatility:} Both SRAM and DRAM are volatile; however, DRAM requires periodic capacitor refresh, whereas SRAM does not.
    \item \textbf{Amdahl's Law Serial Bottleneck:} If $20\%$ of a program is serial, the speedup can never exceed $1/0.20 = 5\times$, regardless of core count.
    \item \textbf{Nyquist Sampling Criterion:} Sampling must be strictly greater than or equal to \textbf{twice the highest frequency} ($f_s \ge 2 f_{\text{max}}$), not the bandwidth or average frequency.
    \item \textbf{ADC Step Count:} An $n$-bit ADC divides the reference voltage into $2^n - 1$ steps, not $2^n$ steps (e.g., 10-bit is 1023 steps).
    \item \textbf{Endianness LSB Placement:} Little-Endian places the Least Significant Byte at the lowest address; Big-Endian places the Most Significant Byte at the lowest address.
    \item \textbf{Linker vs Loader Timing:} The Linker operates at build time (combines object files); the Loader operates at launch time (loads binary into RAM).
\end{enumerate}
